% Lecture notes on the classification of homogeneous plane continua,
% following the formal (Lean 4) proof in this repository.
% Volume 8: Circles of pseudo-arcs and the classification theorem.
%
% Compile with:  pdflatex main && pdflatex main   (two runs for the contents and references)
\documentclass[11pt,oneside]{book}

\usepackage[T1]{fontenc}
\usepackage{lmodern}
\usepackage[margin=1.15in]{geometry}
\usepackage{amsmath,amssymb,amsthm}
\usepackage{enumitem}
\usepackage{booktabs}
\usepackage{longtable}
\usepackage{tikz}
\usepackage[hidelinks]{hyperref}

% ---------------------------------------------------------------- theorem styles
\theoremstyle{plain}
\newtheorem{theorem}{Theorem}[chapter]
\newtheorem{lemma}[theorem]{Lemma}
\newtheorem{proposition}[theorem]{Proposition}
\newtheorem{corollary}[theorem]{Corollary}
\newtheorem{claim}[theorem]{Claim}
\theoremstyle{definition}
\newtheorem{definition}[theorem]{Definition}
\newtheorem{example}[theorem]{Example}
\newtheorem{exercise}[theorem]{Exercise}
\theoremstyle{remark}
\newtheorem{remark}[theorem]{Remark}

% ---------------------------------------------------------------- notation
\newcommand{\N}{\mathbb{N}}
\newcommand{\Z}{\mathbb{Z}}
\newcommand{\Q}{\mathbb{Q}}
\newcommand{\R}{\mathbb{R}}
\newcommand{\C}{\mathbb{C}}
\newcommand{\diam}{\operatorname{diam}}
\newcommand{\dist}{\operatorname{dist}}
\newcommand{\interior}{\operatorname{int}}
\newcommand{\cl}[1]{\overline{#1}}
\newcommand{\abs}[1]{\left\lvert #1\right\rvert}
\newcommand{\set}[1]{\left\{ #1 \right\}}
\newcommand{\fl}[1]{\left\lfloor #1\right\rfloor}
\newcommand{\cel}[1]{\left\lceil #1\right\rceil}
% the unit interval and the circle
\newcommand{\I}{[0,1]}
\newcommand{\Sph}{S^1}
% the fibres of an arc of continua, and the preimage of an interval
\newcommand{\fib}[2]{#1^{-1}(#2)}
\newcommand{\G}[2]{G[#1,#2]}
% the order relation of Volume 1 (r and s strictly between i and j, r nearer to i)
\newcommand{\Btw}{\mathrm{Btw}}
% the star image H_i d(p,i) of Bing and Jones
\newcommand{\st}{\operatorname{st}}
% runs of a circular description
\newcommand{\run}{\operatorname{R}}
% the joint positions and their interpolation
\newcommand{\Ph}{\widehat{P}}
% the covering map of the punctured plane
\newcommand{\cexp}{\operatorname{E}}

% Names of declarations in the Lean development (typewriter, underscores allowed).
\newcommand{\lean}[1]{\texttt{\detokenize{#1}}}
% A Lean name ending in a subscript digit.
\newcommand{\leansub}[2]{\texttt{\detokenize{#1}}\ensuremath{_{#2}}}
% A paragraph pointing to the formal counterpart of a result.
\newcommand{\formal}[1]{\par\smallskip\noindent{\small\textit{Formal counterpart:} #1}\par\smallskip}
% References to numbered results of the earlier volumes.
\newcommand{\vone}[1]{Volume~1, #1}
\newcommand{\vtwo}[1]{Volume~2, #1}
\newcommand{\vthree}[1]{Volume~3, #1}
\newcommand{\vfour}[1]{Volume~4, #1}
\newcommand{\vfive}[1]{Volume~5, #1}
\newcommand{\vsix}[1]{Volume~6, #1}
\newcommand{\vseven}[1]{Volume~7, #1}

\setlist{itemsep=2pt,topsep=4pt}

\title{\textbf{Homogeneous Plane Continua}\\[6pt]
\Large Lecture notes following the formal proof\\[18pt]
\large Volume 8\\[4pt]
\Large Circles of pseudo-arcs\\ and the classification theorem}
\author{}
\date{}

\begin{document}

\frontmatter
\maketitle
\tableofcontents

\mainmatter
% The introduction is Chapter 0, as in Volumes 1--7.
\setcounter{chapter}{-1}
%======================================================================
\chapter{Introduction}\label{ch:intro}
%======================================================================

\section{What this volume is about}

This is the last volume of the notes. It proves the classification of homogeneous plane continua (\vone{Theorem~0.1}):

\begin{theorem}[Theorem~\ref{thm:classification}]
A nondegenerate plane continuum is homogeneous if and only if it is a simple closed curve, a pseudo-arc, or a circle of pseudo-arcs.
\end{theorem}

and its plane form (Theorem~\ref{thm:classification-plane}): there are three plane continua, a simple closed curve, a pseudo-arc and a circle of
pseudo-arcs, each homogeneous, no two homeomorphic, and every nondegenerate homogeneous plane continuum is homeomorphic to exactly one of them.

Volume~6 settled the indecomposable case: a nondegenerate homogeneous indecomposable plane continuum is a pseudo-arc (\vsix{Theorem~7.4}). Volume~7
settled most of the decomposable case: a decomposable homogeneous plane continuum $M$ is a simple closed curve, or there is a continuous, open, surjective
map $f\colon M\to\Sph$ whose fibres are pseudo-arcs and are terminal in $M$ (\vseven{Theorem~9.8}). What is missing is the theory of the third class.

A \emph{circle of pseudo-arcs}, in the sense of Bing and Jones (1959), is a circle-like continuum $X$ with a continuous, open, surjective map $X\to\Sph$
whose fibres are pseudo-arcs. Three things must be proved about it.
\begin{enumerate}
\item \textbf{Uniqueness and homogeneity.} Any two circles of pseudo-arcs are homeomorphic, and each is homogeneous. Bing and Jones derive both from their
Theorem~10: between two ``continuous snake-like arcs of pseudo-arcs'' over $\I$, any homeomorphisms of the end fibres extend to a homeomorphism that
carries each fibre onto the fibre over the same point.
\item \textbf{Recognition.} The map of Volume~7 makes $M$ a circle of pseudo-arcs. Volume~7 does not show that $M$ is circle-like, and an open map onto
$\Sph$ with pseudo-arc fibres does not force it by itself. The map of Volume~7 is \emph{atomic}, and with atomicity the halves of $M$ over the two
half-circles are snake-like. Theorem~10 then makes $M$ homeomorphic to any circle of pseudo-arcs, hence circle-like (Domokos and Prajs, Theorem~6.1).
\item \textbf{Existence in the plane.} Some plane continuum is a circle of pseudo-arcs. Bing and Jones assert this (their \S7) but do not construct the
example. These notes build one by a different construction, a circular analogue of the plane pseudo-arc of \vone{Chapter~5}.
\end{enumerate}

\section{The chapters}

\begin{itemize}
\item \textbf{Chapter~\ref{ch:circles}: Arcs and circles of pseudo-arcs.} The definitions, the two halves of a map onto the circle, and three facts that
distinguish the circle, the pseudo-arc and circles of pseudo-arcs.
\item \textbf{Chapter~\ref{ch:atomic}: Atomic arcs are snake-like.} Bing's proof of his 1951 Theorem~11, with the fibres of an atomic map in place of his
layers: an atomic map onto $\I$ with snake-like hereditarily indecomposable fibres has a snake-like domain. Hence the halves of an atomic map are
snake-like.
\item \textbf{Chapter~\ref{ch:chains}: Chains in the sense of Bing and Jones.} Chains whose nonadjacent links are a positive distance apart; chain maps;
private points; runs; folding a chain along a walk; the form of Bing and Jones's Theorem~6 that is used.
\item \textbf{Chapter~\ref{ch:arcs}: Snake-like arcs of continua.} Bing and Jones's Theorems~3, 4 and~5 and Step~1(b) of Theorem~10. Theorem~5 is false
as printed; it is proved in a repaired form.
\item \textbf{Chapter~\ref{ch:thm78}: Theorems 7 and 8.} Chains following a chain map, on one pseudo-arc (Theorem~7) and on a whole arc of pseudo-arcs
(Theorem~8 and its corollary).
\item \textbf{Chapter~\ref{ch:zigzag}: The homeomorphism from a zigzag.} Two sequences of chains with chain maps going back and forth define a
homeomorphism (Bing and Jones, pp.~188--189).
\item \textbf{Chapter~\ref{ch:thm10}: Theorem 10.} One step of the construction, made alternately on the two sides, and the induction.
\item \textbf{Chapter~\ref{ch:uniqueness}: Uniqueness, homogeneity and recognition.} Gluing two half homeomorphisms; Domokos and Prajs, Theorem~6.1.
\item \textbf{Chapter~\ref{ch:circdescr}: Circular descriptions.} Nested circular chains whose links carry arcs of the base circle; a recognition theorem
for the continuum they describe.
\item \textbf{Chapter~\ref{ch:sweep}: Sweep patterns.} Explicit integer sequences satisfying the conditions of Chapter~\ref{ch:circdescr}, and the
combinatorial ``sweep lemma'' that makes them crooked.
\item \textbf{Chapter~\ref{ch:plane}: A circle of pseudo-arcs in the plane.} The circular analogue of the planar step of \vone{Proposition~5.14}, by a
shear on the universal cover of the punctured plane, and the assembly of a circular description in $\R^2$.
\item \textbf{Chapter~\ref{ch:classification}: The classification theorem.}
\end{itemize}

The main results of the volume are:

\begin{theorem}[Bing 1951, for atomic maps; Theorem~\ref{thm:atomic-snake}]
Let $X$ be a compact metric space and $f\colon X\to\I$ a continuous atomic map whose fibres are snake-like and hereditarily indecomposable. Then $X$ is
snake-like.
\end{theorem}

\begin{theorem}[Bing and Jones, Theorem~10; Theorem~\ref{thm:bj10}]
Let $f\colon X\to\I$ and $g\colon Y\to\I$ be continuous, open and surjective, with pseudo-arc fibres, where $X$ and $Y$ are snake-like compact metric spaces.
Let $h_0\colon\fib f0\to\fib g0$ and $h_1\colon\fib f1\to\fib g1$ be homeomorphisms. Then there is a homeomorphism $H\colon X\to Y$ with $g\circ H=f$ that
extends $h_0$ and $h_1$.
\end{theorem}

\begin{theorem}[Theorems~\ref{thm:unique}, \ref{thm:homog}, \ref{thm:dp61}]
Any two circles of pseudo-arcs are homeomorphic, and each is homogeneous. A compact metric space with an open atomic map onto $\Sph$ whose fibres are
pseudo-arcs is a circle of pseudo-arcs.
\end{theorem}

\begin{theorem}[Theorem~\ref{thm:plane-example}]
Some plane continuum is a circle of pseudo-arcs.
\end{theorem}

\section{Where the formal proof differs from the literature}\label{sec:intro-diff}

\begin{enumerate}
\item \textbf{Halves of an atomic map.} Domokos and Prajs obtain the snake-likeness of the halves from a proposition of Hagopian, Marsh and Prajs on
``atriodic folders''. The formal proof uses instead Bing's 1951 argument (Chapter~\ref{ch:atomic}); no atriodicity is needed.
\item \textbf{Chains.} Bing and Jones require nonadjacent links of a chain to be a positive distance apart; the earlier volumes only ask them to be disjoint.
The theorems are stated with the weaker hypothesis and the stronger conclusion wherever possible, and a shrinking lemma passes between the two
(Theorem~\ref{thm:shrink}).
\item \textbf{Theorems 3, 5, 7 and 8.} The chain printed in the proof of Theorem~3 is not a chain; it is repaired at its two junctions. Theorem~5 is false
as printed (Proposition~\ref{prop:bj5-false}); a repaired form is proved and used, on a window of the arc. The second sentence of Theorem~7 is false for
chains of one link and its printed proof fails for three links; it is proved for every chain of at least three links with a different auxiliary chain. In
Theorem~8 the printed links need not form a chain; they are restricted to the fibres over their intervals.
\item \textbf{Theorem 10.} Bing and Jones extend the homeomorphism to cut fibres in the interior of the arc at each step, using their Theorem~9. The formal
proof uses the homeomorphism only on the two end fibres, and passes from one sub-arc to the next through overlapping windows. Theorem~9 is therefore not
used. One step serves for both sides of the induction, and the induction starts from chains of one link.
\item \textbf{The passage to circles.} ``It follows from Theorem~10 that any two circles of pseudo-arcs are topologically equivalent and that each is
homogeneous'' is not argued by Bing and Jones; it is supplied by cutting the circle into two halves (Chapter~\ref{ch:uniqueness}).
\item \textbf{The plane example.} Bing and Jones's \S7 asserts a sequence of circular chains with nine properties without constructing it, and does not argue
that the resulting continuum has pseudo-arc fibres; their condition~(9) does not visibly give it. The example of Chapters~\ref{ch:circdescr}--\ref{ch:plane}
is built by another route.
\end{enumerate}

\section{Prerequisites and conventions}\label{sec:intro-conv}

The reader should know chains, crookedness, descriptions and the plane pseudo-arc (Volume~1); pattern following, the uniqueness and homogeneity of the
pseudo-arc (Volume~2); boundary bumping, snake-like and circle-like continua (\vthree{Chapters~1 and~5}); and the main theorems of Volumes~6 and~7.

\begin{itemize}
\item \textbf{Continua.} A continuum is compact, connected and metric; it may be a point unless ``nondegenerate'' is said. A pseudo-arc is a nondegenerate,
snake-like, hereditarily indecomposable continuum (\vone{Definition~2.14}).
\item \textbf{Chains} are as in \vone{Definition~2.1}: open links $D_0,\dots,D_{n-1}$ with $D_i\cap D_j\ne\emptyset$ if and only if $\abs{i-j}\le1$. Links are
numbered from~$0$. Bing and Jones number from~$1$; quotations are translated.
\item \textbf{Fibres.} For $f\colon X\to\I$ and $x\in\I$ the fibre $\fib fx$ is also written $p_x$ (and $q_x$ for $g\colon Y\to\I$), following Bing and Jones.
For $a\le b$ we write $\G ab=\fib f{[a,b]}$.
\item \textbf{The circle} is $\Sph=\set{z\in\C:\abs z=1}$, and $e^{i\theta}$ is written $\exp(i\theta)$. The plane carries the sup metric, as in the earlier
volumes.
\end{itemize}

\paragraph{Formal counterparts.} As in the earlier volumes, many results are followed by a line ``Formal counterpart:'' naming the corresponding declaration of the
Lean development (all names begin with \lean{HPC.}). These lines may be skipped. Appendix~\ref{app:formal} collects them.

\paragraph{Proofs in outline.} Some long constructions (Bing's surgery and Fugate's lemma in Chapter~\ref{ch:atomic}, Theorem~7 in Chapter~\ref{ch:thm78}, the
step of Theorem~10, the planar step of Chapter~\ref{ch:plane}) are given as proofs in outline, marked as such. The complete arguments are in the declarations named in the ``Formal counterpart'' lines.

\paragraph{Exercises.} Each chapter ends with exercises. None of them is used later.

%======================================================================
\chapter{Arcs and circles of pseudo-arcs}\label{ch:circles}
%======================================================================

This chapter fixes the definitions of the volume and proves the elementary facts about them: the two halves of a map onto the circle, and three
facts that tell the circle, the pseudo-arc and the circles of pseudo-arcs apart.

\section{Definitions}

\begin{definition}[Bing and Jones, p.~173]\label{def:arc-continua}
Let $X$ be a metric space and $f\colon X\to\I$.
\begin{enumerate}[label=(\alph*)]
\item $f$ is an \emph{arc of continua} if $f$ is continuous and surjective and every fibre $\fib fx$ is a continuum.
\item $f$ is a \emph{continuous arc of pseudo-arcs} if $f$ is continuous, open and surjective and every fibre $\fib fx$ is a pseudo-arc.
\end{enumerate}
\end{definition}

\formal{\lean{HPC.BJ59.IsArcOfContinua}, \lean{HPC.BJ59.IsContArcPA}.}

Bing and Jones speak of a \emph{continuous snake-like arc of pseudo-arcs}; snake-likeness of $X$ is stated separately here. ``Continuous'' refers to
the decomposition into fibres, which is continuous exactly when $f$ is open (\vseven{Lemma~9.7} for one direction).

\begin{definition}[Bing and Jones, p.~173]\label{def:circle-pa}
A metric space $X$ is a \emph{circle of pseudo-arcs} if $X$ is circle-like (\vthree{Definition~5.13}) and there is a continuous, open, surjective map
$f\colon X\to\Sph$ all of whose fibres are pseudo-arcs. Such a map, on a circle-like space, is a \emph{circle-of-pseudo-arcs map}.
\end{definition}

\formal{\lean{HPC.BJ59.IsCircleOfPA}, \lean{HPC.BJ59.IsCircPAMap}; the equivalence of ``$X$ is a circle of pseudo-arcs'' with ``$X$ carries a
circle-of-pseudo-arcs map'' is \lean{HPC.BJ59.isCircleOfPA_iff}.}

\begin{remark}\label{rem:nonsnake}
Bing and Jones call a continuum circle-like only if it is moreover not snake-like. For a circle of pseudo-arcs this is automatic: by
Corollary~\ref{cor:circle-pa-decomp} below $X$ is the union of two continua whose intersection is the union of two disjoint fibres, so $X$ is not
unicoherent, while snake-like continua are (\vthree{Theorem~5.4}). The definition above omits the condition.
\end{remark}

Recall also the atomic maps of \vseven{Definition~9.6}: $f\colon X\to Y$ is \emph{atomic} if every continuum in $X$ meeting a fibre $\fib fy$ lies in it or
contains it; an \emph{atomic circle-of-pseudo-arcs map} is a continuous, open, surjective, atomic map $X\to\Sph$ whose fibres are pseudo-arcs. Here
circle-likeness is not assumed. By \vseven{Theorem~9.8}, a decomposable homogeneous plane continuum which is not a simple closed curve carries such a map.

\begin{example}\label{ex:not-atomic}
The second projection $\I\times\I\to\I$ is open and has arcs as fibres, but it is not atomic: the continuum $\set0\times\I$ meets the fibre $\I\times\set t$
without lying in it or containing it. The projection $P\times\Sph\to\Sph$, $P$ a pseudo-arc, fails to be atomic in the same way. The identity of any
space is atomic.
\end{example}

\formal{\lean{HPC.BJ59.not_isAtomicMap_snd}, \lean{HPC.BJ59.isAtomicMap_id}.}

\section{Preimages of connected sets}

\begin{lemma}\label{lem:monotone-preimage}
Let $X$ be compact, $q\colon X\to Y$ continuous into a metric space, with every fibre connected. If $A\subseteq q(X)$ is connected, then $q^{-1}(A)$ is
connected.
\end{lemma}

\begin{proof}
Suppose $q^{-1}(A)\subseteq U\cup V$ with $U,V$ open, disjoint, both meeting $q^{-1}(A)$. Each fibre over a point of $A$ is connected, so it lies in $U$ or
in $V$. Put $A_U=\set{a\in A:q^{-1}(a)\subseteq U}$ and $A_V$ likewise. They are disjoint, cover $A$ and are nonempty. Since $q$ is a closed map ($X$ is
compact), $Y\setminus q(X\setminus U)$ is open, and $A_U=A\setminus q(X\setminus U)$ (a point of $A$ lies in $q(X)$, so its fibre is nonempty).
So $A_U$, and likewise $A_V$, is relatively open in $A$, and $A$ is not connected.
\end{proof}

\formal{\lean{HPC.BJ59.isPreconnected_preimage_of_monotone'}.}

\begin{corollary}\label{cor:G-connected}
Let $X$ be compact and $f\colon X\to\I$ an arc of continua. For $a\le b$ the set $\G ab=\fib f{[a,b]}$ is a continuum, and so is $X$.
\end{corollary}

\formal{\lean{HPC.BJ59.isPreconnected_preimage_Icc}, \lean{HPC.BJ59.isCompact_preimage_Icc}.}

\section{The two halves}

\begin{definition}\label{def:halves}
For $s\in\set{1,-1}$ let $c_s\colon\I\to\Sph$, $c_s(t)=\exp(i\pi st)$, and let $A_s=c_s(\I)$, the closed upper ($s=1$) or lower ($s=-1$) half-circle from $1$
to $-1$. The map $c_s$ is a homeomorphism of $\I$ onto $A_s$. For $f\colon X\to\Sph$ the \emph{half} of $X$ over $A_s$ is $X_s=\fib f{A_s}$, with the
\emph{half map} $f_s=c_s^{-1}\circ f\colon X_s\to\I$.
\end{definition}

\formal{\lean{HPC.Jones.halfC}, \lean{HPC.BJ59.halfParam}, \lean{HPC.BJ59.halfSet}, \lean{HPC.BJ59.halfMap}.}

Thus $A_1\cup A_{-1}=\Sph$ and $A_1\cap A_{-1}=\set{1,-1}$, and the point $\exp(-i\pi s/2)$ is not in $A_s$. The fibre of $f_s$ over $t$ is the fibre of $f$
over $c_s(t)$; the end fibres of $f_s$ are $\fib f1$ and $\fib f{-1}$, for both values of $s$.

\formal{\lean{HPC.BJ59.range_halfC_union}, \lean{HPC.BJ59.mem_range_halfC_inter}, \lean{HPC.BJ59.exists_not_mem_range_halfC},
\lean{HPC.BJ59.image_val_halfMap_fibre}, \lean{HPC.BJ59.halfFibreEquiv}.}

\begin{proposition}\label{prop:halves}
Let $X$ be compact and $f\colon X\to\Sph$ continuous, open and surjective, with pseudo-arc fibres. Then each half map $f_s\colon X_s\to\I$ is a continuous
arc of pseudo-arcs, and $X_s$ is a continuum. If $f$ is atomic, so is $f_s$.
\end{proposition}

\begin{proof}
$f_s$ is continuous and surjective, and its fibres are fibres of $f$. It is open: an open set of $X_s$ is $W\cap X_s$ with $W$ open in $X$, and
$f(W\cap X_s)=f(W)\cap A_s$ because $X_s$ is the full preimage of $A_s$; this is open in $A_s$, and $c_s$ is a homeomorphism. $X_s$ is closed, hence
compact, and connected by Lemma~\ref{lem:monotone-preimage}, $A_s$ being connected. A continuum in $X_s$ is a continuum in $X$, and the fibres of $f_s$ are
those of $f$, so atomicity passes to $f_s$.
\end{proof}

\formal{\lean{HPC.BJ59.isContArcPA_halfMap'} (with \lean{HPC.BJ59.continuous_halfMap}, \lean{HPC.BJ59.isOpenMap_halfMap},
\lean{HPC.BJ59.surjective_halfMap}), \lean{HPC.BJ59.isContinuum_halfSet}, \lean{HPC.BJ59.isAtomicMap_halfMap}.}

\begin{proposition}\label{prop:halves-snake}
If $f\colon X\to\Sph$ is a circle-of-pseudo-arcs map, each half $X_s$ is snake-like.
\end{proposition}

\begin{proof}
$X$ is compact, being circle-like, and $X_s$ is a continuum (Proposition~\ref{prop:halves}). It is a proper subcontinuum of $X$: it misses the fibre over
$\exp(-i\pi s/2)$, which is nonempty since $f$ is onto. A proper subcontinuum of a circle-like continuum is snake-like (\vthree{Theorem~5.16}), and
snake-likeness is a property of the space $X_s$ (\vthree{Theorem~5.11}).
\end{proof}

\formal{\lean{HPC.BJ59.isContArcPA_halfMap}, \lean{HPC.BJ59.isSnakeLike_halfSet}, from \lean{HPC.Plane.isSnakeLike_of_subset_circleLike}.}

Both kinds of maps of this volume, the circle-of-pseudo-arcs maps and the atomic ones, will be used only through the following properties.

\begin{definition}\label{def:snake-halves}
A map $f\colon X\to\Sph$ \emph{has snake-like halves} if it is continuous, open and surjective, its fibres are pseudo-arcs, and both halves $X_1$, $X_{-1}$
are snake-like.
\end{definition}

\formal{\lean{HPC.BJ59.IsSnakeHalvesMap}; a circle-of-pseudo-arcs map has snake-like halves by \lean{HPC.BJ59.IsCircPAMap.isSnakeHalvesMap}, and an
atomic one on a compact space by \lean{HPC.BJ59.IsAtomicCircPAMap.isSnakeHalvesMap} (Chapter~\ref{ch:atomic}).}

Rotations preserve all these properties: if $f$ is a circle-of-pseudo-arcs map (or an atomic one), so is $x\mapsto cf(x)$ for every $c\in\Sph$, since its
fibres are the fibres of $f$.

\formal{\lean{HPC.BJ59.IsCircPAMap.mul_left}, \lean{HPC.BJ59.IsAtomicCircPAMap.mul_left}.}

\section{The three classes are distinct}

\begin{proposition}\label{prop:circle-decomp}
The circle is decomposable.
\end{proposition}

\begin{proof}
$\Sph=A_1\cup A_{-1}$, and each $A_s$ is a continuum (the image of $\I$) different from $\Sph$ (it misses $\exp(-i\pi s/2)$).
\end{proof}

\formal{\lean{HPC.BJ59.isDecomposableSpace_circle}.}

\begin{corollary}\label{cor:circle-pa-decomp}
If $f\colon X\to\Sph$ is continuous, open and surjective with pseudo-arc fibres, and $X$ is compact, then $X=X_1\cup X_{-1}$ is decomposable, and
$X_1\cap X_{-1}=\fib f1\cup\fib f{-1}$. In particular every circle of pseudo-arcs is decomposable.
\end{corollary}

\begin{proof}
By Proposition~\ref{prop:halves} the halves are continua; each misses a nonempty fibre, so each is proper. The intersection is the preimage of
$A_1\cap A_{-1}=\set{1,-1}$.
\end{proof}

\formal{\lean{HPC.BJ59.isDecomposableSpace_of_isCircPAMap}.}

\begin{proposition}\label{prop:pa-not-in-circle}
No proper subset of the circle is a pseudo-arc.
\end{proposition}

\begin{proof}
Let $K\subsetneq\Sph$ be a nondegenerate continuum and $z\notin K$. After a rotation we may assume $z=-1$. The argument $\arg\colon\Sph\setminus\set{-1}\to
(-\pi,\pi)$ is continuous and injective, so it maps the compact set $K$ homeomorphically onto a nondegenerate compact connected subset of $\R$, that is, onto
an interval $[a,b]$ with $a<b$. An interval is decomposable ($[a,b]=[a,c]\cup[c,b]$ for $a<c<b$), so $K$ is not hereditarily indecomposable, and not a
pseudo-arc.
\end{proof}

\formal{\lean{HPC.BJ59.not_isPseudoArc_of_ne_univ_circle}, with \lean{HPC.BJ59.continuousAt_arg_circle}, \lean{HPC.BJ59.arg_injective_circle} and
\lean{HPC.BJ59.not_isHereditarilyIndecomposable_of_nontrivial_real}.}

\begin{corollary}\label{cor:circle-not-cpa}
The circle is not homeomorphic to a circle of pseudo-arcs.
\end{corollary}

\begin{proof}
Let $f\colon X\to\Sph$ be a circle-of-pseudo-arcs map and $\Phi\colon X\to\Sph$ a homeomorphism. The fibre $\fib f1$ is a pseudo-arc, and so is its image
under $\Phi$ (being a pseudo-arc is a topological property, \vthree{Theorem~5.11} and the invariance of hereditary indecomposability). This image misses
$\Phi(x)$ for any $x\in\fib f{-1}$, so it is a proper subset of $\Sph$, against Proposition~\ref{prop:pa-not-in-circle}.
\end{proof}

\formal{Inside the proof of \lean{HPC.BJ59.ho16_classification_plane}, with \lean{HPC.BJ59.isPseudoArc_of_homeomorph}.}

The pseudo-arc is indecomposable, so by Proposition~\ref{prop:circle-decomp} and Corollary~\ref{cor:circle-pa-decomp} it is homeomorphic neither to the circle
nor to a circle of pseudo-arcs.

\section{Two facts about the circle}

\begin{proposition}\label{prop:circle-homog}
The circle is homogeneous, and so is every space homeomorphic to it. A simple closed curve lies in the plane.
\end{proposition}

\begin{proof}
For $x,y\in\Sph$ multiplication by $yx^{-1}$ is a homeomorphism carrying $x$ to $y$. If $e\colon W\to Y$ is a homeomorphism and $W$ is homogeneous, then
for $x,y\in Y$ a homeomorphism $g$ of $W$ with $g(e^{-1}x)=e^{-1}y$ gives $e\circ g\circ e^{-1}$. The unit circle of $\R^2$ is a simple closed curve.
\end{proof}

\formal{\lean{HPC.BJ59.isHomogeneous_circle'}, \lean{HPC.BJ59.isHomogeneous_circle}, \lean{HPC.Span.isHomogeneous_of_homeomorph},
\lean{HPC.BJ59.exists_plane_circle}.}

\section*{Exercises}

\begin{exercise}
Show that the preimage of every closed arc of $\Sph$ under a circle-of-pseudo-arcs map is a continuum, and that the preimage of every proper closed arc is
snake-like.
\end{exercise}

\begin{exercise}
Show that Lemma~\ref{lem:monotone-preimage} fails without compactness of $X$. (Consider the map $[0,1)\cup\set2\to\I$ that is the identity on $[0,1)$ and
sends $2$ to $1$.)
\end{exercise}

\begin{exercise}
Let $f\colon X\to\Sph$ be an atomic circle-of-pseudo-arcs map. Show that $X$ contains no arc. (An arc in $X$ either lies in a fibre or contains one.)
\end{exercise}

%======================================================================
\chapter{Atomic arcs are snake-like}\label{ch:atomic}
%======================================================================

Let $f\colon M\to\Sph$ be the map of \vseven{Theorem~9.8}: open, atomic, with pseudo-arc fibres. To compare $M$ with a circle of pseudo-arcs we cut it
into its two halves (Chapter~\ref{ch:circles}), and Theorem~10 of Bing and Jones needs each half to be snake-like. This chapter proves that. The tool
is Bing's 1951 theorem that a continuum built from snake-like hereditarily indecomposable ``layers'' is snake-like; the fibres of an atomic map play the
role of the layers.

Throughout, $X$ is a compact metric space. A continuum $K$ is \emph{$\varepsilon$-coverable} if some chain of mesh $<\varepsilon$ (an
\emph{$\varepsilon$-chain}) covers it; $K$ is snake-like if it is $\varepsilon$-coverable for every $\varepsilon>0$ (\vthree{Chapter~5}). A chain $F$
\emph{refines} a chain $D$ if each link of $F$ lies in some link of $D$.

\formal{\lean{HPC.Plane.EpsCoverable}, \lean{HPC.Plane.isEpsChain_of_refines}.}

\section{Atomic maps onto the interval}

\begin{lemma}[Saturation]\label{lem:saturation}
Let $f\colon X\to Y$ be atomic and $L\subseteq X$ a continuum. Either $L$ lies in one fibre, or $L=f^{-1}(f(L))$; in the second case $L$ contains every
fibre it meets.
\end{lemma}

\begin{proof}
Suppose $L$ lies in no fibre. If $x\in f^{-1}(f(L))$, pick $l\in L$ with $f(l)=f(x)$. The continuum $L$ meets the fibre $\fib f{f(l)}$ and does not lie in
it, so it contains it, and $x\in L$.
\end{proof}

\formal{\lean{HPC.BJ59.subset_fibre_or_saturated}, \lean{HPC.BJ59.fibre_subset_of_saturated}.}

From now on $f\colon X\to\I$ is continuous and atomic, and its fibres are continua. Then $\G ab=\fib f{[a,b]}$ is a continuum for $a\le b$, and $X$ itself is
a continuum (Corollary~\ref{cor:G-connected}).

\formal{\lean{HPC.BJ59.isPreconnected_preimage_of_fibres}, \lean{HPC.BJ59.isContinuum_univ_of_fibres}.}

Recall Bing's \emph{condition B} for a point $p$ of a continuum $M$: of any two subcontinua of $M$ containing $p$, one contains the other.

\begin{lemma}[Terminal points]\label{lem:condB}
Suppose the fibres of $f$ are hereditarily indecomposable, and let $a\le b$. Every point $\mu$ with $f(\mu)\in\set{a,b}$ satisfies condition B in $\G ab$.
\end{lemma}

\begin{proof}
Let $H,K\subseteq\G ab$ be continua containing $\mu$, and let $F=\fib f{f(\mu)}$. By Lemma~\ref{lem:saturation} each of $H$, $K$ lies in $F$ or is saturated.
If both lie in $F$, then $H\cup K$ is a continuum in the hereditarily indecomposable $F$, so one of $H$, $K$ contains the other (otherwise $H\cup K$ would
be decomposable). If $H\subseteq F$ and $K$ is saturated, then $K\supseteq F\supseteq H$ by the second part of the saturation lemma; symmetrically in the
opposite case. If both are saturated, $f(H)$ and $f(K)$ are intervals in $[a,b]$ containing the end point $f(\mu)$ of $[a,b]$, so they are nested, and so
are their preimages $H$ and $K$.
\end{proof}

\formal{\lean{HPC.BJ59.condB_preimage_Icc}, with \lean{HPC.BJ59.nested_of_common_bound}.}

\begin{proposition}\label{prop:atomic-hu}
If the fibres of $f$ are snake-like, $X$ is hereditarily unicoherent: for subcontinua $H$, $L$ with $H\cup L$ a continuum, $H\cap L$ is connected.
\end{proposition}

\begin{proof}
Let $K=H\cup L$. Since $K$ is connected, $H\cap L\neq\emptyset$. If $K$ lies in a fibre, use that snake-like continua are hereditarily unicoherent
(\vthree{Theorem~5.4}). Otherwise, if $H$ lies in a fibre $\fib fy$, then $L$ meets $\fib fy$ (at a point of $H\cap L$), and $L\not\subseteq\fib fy$ (else
$K\subseteq\fib fy$), so $L\supseteq\fib fy\supseteq H$ and $H\cap L=H$. Likewise if $L$ lies in a fibre. In the remaining case both are saturated, and
$H\cap L=f^{-1}(f(H)\cap f(L))$ is the preimage of the intersection of two intervals, which is connected by Lemma~\ref{lem:monotone-preimage}.
\end{proof}

\formal{\lean{HPC.BJ59.isHereditarilyUnicoherent_of_atomic}, with \lean{HPC.Plane.isHereditarilyUnicoherent_of_isSnakeLike}.}

\section{Bing's minimal counterexample}

\begin{lemma}[Brouwer reduction]\label{lem:brouwer}
If a continuum $M\subseteq X$ is not $\varepsilon$-coverable, it contains a subcontinuum $S$ which is not $\varepsilon$-coverable, while every proper
subcontinuum of $S$ is.
\end{lemma}

\begin{proof}
Order the subcontinua of $M$ that are not $\varepsilon$-coverable by inclusion. A decreasing sequence $S_1\supseteq S_2\supseteq\cdots$ of them has
intersection $S_\infty$, a continuum (\vthree{Chapter~1}). If a chain covered $S_\infty$, its union $W$ would be an open set containing
$\bigcap_nS_n$, so $S_n\subseteq W$ for some $n$ by compactness; hence $S_\infty$ is not $\varepsilon$-coverable either. In a compact metric space this
sequential condition suffices for a minimal element (the form of Zorn's lemma of \vthree{Chapter~2}).
\end{proof}

\formal{\lean{HPC.Plane.exists_minimal_notCoverable}.}

The heart of Bing's argument is a way of moving a point to the end of a chain. Bing proved it for a point satisfying condition B (his 1951 Theorem~10);
Fugate extended it to terminal subcontinua. A subcontinuum $A$ of $M$ is \emph{terminal} if for subcontinua $K,L$ of $M$ both meeting $A$,
$K\subseteq A\cup L$ or $L\subseteq K\cup A$. For $A=\set p$ this is condition B.

\formal{\lean{HPC.Plane.IsTerminalSub}, \lean{HPC.Plane.isTerminalSub_singleton_iff}.}

\begin{theorem}[Bing 1951, Theorem~10; Fugate 1966, Lemma~1]\label{thm:fugate1}
Let $M$ be a continuum, $A$ a terminal subcontinuum of $M$ and $D$ a chain covering $M$. There is a chain $F$ covering $M$ and refining $D$ whose last link
contains a point of $A$ outside the closure of the union of the other links.
\end{theorem}

The proof is long; we give its structure. Call a chain $F$ a \emph{Theorem-10 chain for $(K,p)$ under $D$} if it covers $K$, refines $D$, has $p$ in its
first link and $p$ in the closure of no other link. Theorem~\ref{thm:fugate1} asks for a Theorem-10 chain for $(M,p)$ with some $p\in A$ (then reverse it).

\formal{\lean{HPC.Plane.IsT10Chain}.}

\begin{lemma}[Bing's surgery]\label{lem:surgery}
Let $D=(D_0,\dots,D_{m-1})$ cover the continuum $M$, let $p\in M$, and let $K\subseteq M$. Suppose $F=(F_0,\dots,F_{k-1})$ is a Theorem-10 chain for
$(K,p)$ under $D$, and that some $y\ne p$ lies in $F_{k-1}\cap D_0$, such that every subcontinuum of $M$ containing $p$ and not $y$ lies in $K$. Then there
is a Theorem-10 chain for $(M,p)$ under $D$.
\end{lemma}

\begin{proof}[Proof in outline]
Take a small open ball $N$ about $y$ with $\cl N\subseteq F_{k-1}\cap D_0$ and $p\notin\cl N$. The component of $p$ in the compact set $M\setminus N$ lies in
$K$, hence in the union of $F$. A component of a compact space is the intersection of the clopen sets containing it (\vthree{Chapter~1}), so $M\setminus N$
splits into disjoint compact sets $H\ni p$ and $H'$ with $H$ inside the union of $F$. Let $U\supseteq H$ and $V\supseteq H'$ be open with disjoint
closures and $p\notin\cl V$. The sequence
\[
F_0\cap U,\ \dots,\ F_{k-2}\cap U,\ (F_{k-1}\cap U)\cup N,\ D_0\cap V,\ (D_1\cap V)\setminus\cl N,\ \dots,\ (D_{m-1}\cap V)\setminus\cl N
\]
covers $M$, its links are open and refine $D$, and links two or more places apart are disjoint. Removing empty links and trimming (\vone{Lemma~2.5}
and its companion for such ``weak chains'') gives a chain with the required properties; $p$ is still only in the first link.
\end{proof}

\formal{\lean{HPC.Plane.surgery_of_sep} (Bing's original form with condition B is \lean{HPC.Plane.theorem10_1951_of_lastLink}); the trimming is
\lean{HPC.Plane.exists_chain_of_weak}.}

\begin{lemma}[The composant step]\label{lem:composant-step}
If $A$ is a proper terminal subcontinuum of $M$ and every link of a chain $D$ covering $M$ meets $M$, some proper subcontinuum of $M$ containing $A$ meets
every link of $D$.
\end{lemma}

\begin{proof}
Fix $z\in M\setminus A$. For each link $D_i$, boundary bumping (\vthree{Chapter~1}) gives a subcontinuum $K_i\subseteq M$ meeting $A$ and $D_i$, and missing
$z$. By terminality the sets $A\cup K_i$ are linearly ordered by inclusion; the largest one is a proper subcontinuum of $M$ (it misses $z$) containing $A$
and meeting every link.
\end{proof}

\formal{\lean{HPC.Plane.exists_proper_superset_meeting_all_links}, with \lean{HPC.Plane.IsTerminalSub.mono}.}

\begin{proof}[Proof of Theorem~\ref{thm:fugate1} in outline]
Strong induction on the number $m$ of links of $D$. Suppose the statement fails for some $M\supseteq A$ under $D$; by Zorn's lemma take a minimal such
$M'$ among the subcontinua of $M$ containing $A$ (terminality passes to them). Then every link of $D$ meets $M'$ (else apply the induction hypothesis to the
shorter chain). If $A$ met the first or last link of $D$, removing a small closed ball about a point $p\in A$ there from the other links would give a
Theorem-10 chain directly. So $A$ misses the end links. Lemma~\ref{lem:composant-step} gives a proper subcontinuum $K\supseteq A$ of $M'$ meeting every
link; by minimality $K$ has a Theorem-10 chain $F$ from some $p\in A$. Choose $F$ so that its last link is as low as possible; it contains a point
$y\in K\cap(D_0\cup D_{m-1})$, and $y\notin A$. By terminality, every subcontinuum of $M'$ containing $p$ and not $y$ lies in $K$, and
Lemma~\ref{lem:surgery} (or its mirror form) gives a Theorem-10 chain for $M'$, a contradiction.
\end{proof}

\formal{\lean{HPC.Plane.fugate_lemma1}, from \lean{HPC.Plane.lemma1_aux}, \lean{HPC.Plane.t10Chain_of_mem_first},
\lean{HPC.Plane.t10Chain_of_mem_last}. Fugate assumes $M$ atriodic and hereditarily unicoherent; the formal proof does not use either hypothesis.}

\section{Gluing two chains}

\begin{proposition}[Fugate, p.~392]\label{prop:glue-halves}
Let $M$ be a hereditarily unicoherent continuum, $M=A\cup C$ with $A$, $C$ continua, and $\mu\in A\cap C$. Suppose that every subcontinuum $Q\subseteq C$
which meets $A\cap C$ but does not lie in it contains $A\cap C$ (``Claim 1''). If $A$ is covered by an $\varepsilon$-chain $E=(E_0,\dots,E_{x-1})$ with
$\mu\in E_{x-1}$, and $C$ by an $\varepsilon$-chain $F=(F_0,\dots,F_{y-1})$ with $\mu\in F_{y-1}$, then $M$ is $\varepsilon$-coverable.
\end{proposition}

\begin{proof}
Take an open $O\ni\mu$ with $\cl O\subseteq E_{x-1}\cap F_{y-1}$, and let $E^*$ be the union of $E$.

\emph{Claim 2: no connected $Z\subseteq M\setminus O$ meets both $A\setminus O$ and $M\setminus E^*$.} Otherwise let $N=\cl Z$, a continuum in $M\setminus O$
containing a point $a\in A$ and a point $z\notin E^*$; then $z\notin A$, so $z\in C$. Since $N$ meets $A$ and $C$ and lies in $A\cup C$, it meets $A\cap C$.
By hereditary unicoherence $Q=N\cap C$ is a continuum; it lies in $C$, meets $A\cap C$, and contains $z\notin A\cap C$. By Claim~1, $A\cap C\subseteq Q\subseteq
N$. But $\mu\in A\cap C$ and $\mu\in O$, which $N$ misses.

By Claim~2 and the separation of two closed sets in a compact space when no connected set meets both (\vthree{Chapter~1}), $M\setminus O=H\cup L$ with $H$,
$L$ disjoint and closed, $A\setminus O\subseteq H$ and $M\setminus E^*\subseteq L$. Let $S\supseteq H$, $T\supseteq L$ be disjoint open sets. Consider
\[
(E_0\cap S)\setminus\cl O,\ \dots,\ (E_{x-2}\cap S)\setminus\cl O,\ E_{x-1}\cap S,\ (F_{y-1}\cap T)\cup O,\ F_{y-2}\cap T,\ \dots,\ F_0\cap T.
\]
Each link lies in a link of $E$ or of $F$ (note $O\subseteq F_{y-1}$), so the mesh is $<\varepsilon$. Links two or more places apart are disjoint: within one
half by the chain property, across the middle because $S\cap T=\emptyset$ and the first $x-1$ links miss $\cl O$. The links cover $M$: a point of $O$ is in
the middle link; a point of $H$ lies in $E^*\cap S$, and if it lies in $\cl O$ it lies in $E_{x-1}\cap S$; a point of $L$ lies in $T$ and in $C$ (if it were
in $A$ it would be in $A\setminus O\subseteq H$), hence in some $F_j\cap T$. Trimming the weak chain gives an $\varepsilon$-chain covering $M$.
\end{proof}

\formal{\lean{HPC.Plane.epsCoverable_of_halves}, with Claim~2 as \lean{HPC.Plane.claim2}. Fugate uses $\varepsilon/2$-chains; no mesh is lost.}

\section{The theorem}

\begin{theorem}[Bing 1951, Theorem~11, for atomic maps]\label{thm:atomic-snake}
Let $X$ be a compact metric space and $f\colon X\to\I$ a continuous atomic map whose fibres are snake-like and hereditarily indecomposable. Then $X$ is
snake-like.
\end{theorem}

\begin{proof}
$X$ is a continuum. Suppose it is not $\varepsilon$-coverable, and let $S$ be a minimal non-$\varepsilon$-coverable subcontinuum (Lemma~\ref{lem:brouwer}).
$S$ lies in no fibre, since fibres are snake-like; by Lemma~\ref{lem:saturation}, $S=f^{-1}(f(S))=\G ab$ with $f(S)=[a,b]$, and $a<b$. Let $c=(a+b)/2$,
$A=\G ac$ and $C=\G cb$. These are continua with $A\cup C=S$ and $A\cap C=\fib fc$, and both are proper subcontinua of $S$ (they miss $\fib fb$ and $\fib
fa$), hence $\varepsilon$-coverable. Fix $\mu\in\fib fc$.

By Lemma~\ref{lem:condB}, $\mu$ satisfies condition B in $A$ and in $C$. Apply Theorem~\ref{thm:fugate1} to $A$ with the terminal subcontinuum $\set\mu$ and an $\varepsilon$-chain
covering $A$: we get an $\varepsilon$-chain covering $A$ with $\mu$ in its last link. Likewise for $C$. Claim~1 of Proposition~\ref{prop:glue-halves} is
atomicity at the fibre $A\cap C=\fib fc$, and $S$ is hereditarily unicoherent (Proposition~\ref{prop:atomic-hu}). So $S$ is $\varepsilon$-coverable, a
contradiction.
\end{proof}

\formal{\lean{HPC.BJ59.isSnakeLike_of_atomic'}; the case of pseudo-arc fibres is \lean{HPC.BJ59.isSnakeLike_of_atomic}. Openness of $f$ is not needed.
Bing's own theorem, for his ``layers'', is \lean{HPC.Plane.theorem11_1951}.}

\begin{corollary}\label{cor:atomic-halves}
If $X$ is compact and $f\colon X\to\Sph$ is an atomic circle-of-pseudo-arcs map, each half $f_s\colon X_s\to\I$ is a continuous arc of pseudo-arcs and each
half $X_s$ is snake-like. Hence $f$ has snake-like halves.
\end{corollary}

\begin{proof}
By Proposition~\ref{prop:halves}, $f_s$ is a continuous, open, atomic arc of pseudo-arcs on the compact space $X_s$. Its fibres are pseudo-arcs, so snake-like
and hereditarily indecomposable, and Theorem~\ref{thm:atomic-snake} applies.
\end{proof}

\formal{\lean{HPC.BJ59.isContArcPA_halfMap_of_atomic}, \lean{HPC.BJ59.isSnakeLike_halfSet_of_atomic},
\lean{HPC.BJ59.IsAtomicCircPAMap.isSnakeHalvesMap}.}

\begin{remark}
Domokos and Prajs obtain Corollary~\ref{cor:atomic-halves} from a proposition of Hagopian, Marsh and Prajs, which requires the halves to be atriodic. The
argument above needs no atriodicity.
\end{remark}

\section*{Exercises}

\begin{exercise}
Show that Lemma~\ref{lem:condB} fails for points of an interior fibre of $\G ab$: give two subcontinua through such a point neither of which contains the
other.
\end{exercise}

\begin{exercise}
Show that the projection $P\times\I\to\I$, $P$ a pseudo-arc, satisfies every hypothesis of Theorem~\ref{thm:atomic-snake} except atomicity, and that
$P\times\I$ is not snake-like. (For three points $p_1,p_2,p_3\in P$, the continuum $P\times\set0\cup\set{p_1,p_2,p_3}\times\I$ is a
triod.)
\end{exercise}

\begin{exercise}
Prove Lemma~\ref{lem:composant-step} for $A=\set p$ directly from condition B.
\end{exercise}

%======================================================================
\chapter{Chains in the sense of Bing and Jones}\label{ch:chains}
%======================================================================

The proof of Theorem~10 manipulates chains on two spaces at once and maps of their link indices. This chapter collects the facts about chains that it uses.
Throughout, $X$ is a metric space and chains are as in \vone{Definition~2.1}, with links numbered from $0$. For a chain $D=(D_0,\dots,D_{n-1})$ a
\emph{run} is a set of consecutive links $D_u,\dots,D_v$ ($u\le v$); we say the run $[u,v]$ \emph{covers} a set $C$ if $C\subseteq D_u\cup\dots\cup D_v$.

\section{Chains whose nonadjacent links are apart}

\begin{definition}[Bing and Jones, p.~172]\label{def:bjchain}
A chain $D$ is a \emph{BJ-chain} if any two links $D_i$, $D_j$ with $\abs{i-j}\ge2$ are a positive distance apart: there is $\delta>0$ with
$\dist(x,y)\ge\delta$ for $x\in D_i$, $y\in D_j$.
\end{definition}

\formal{\lean{HPC.BJ59.IsBJChain}. These are Fugate's \emph{taut} chains: \lean{HPC.BJ59.isBJChain_iff_isTaut}.}

In a BJ-chain nonadjacent links have disjoint closures. Hence:

\begin{lemma}\label{lem:adj-closure}
If $D$ is a BJ-chain and a point lies in $\cl{D_a}\cap\cl{D_b}$, then $\abs{a-b}\le1$.
\end{lemma}

\formal{\lean{HPC.BJ59.IsBJChain.adjacent_of_closure}.}

\begin{lemma}\label{lem:adjacent-meet}
Let $K$ be connected and let $U_0,\dots,U_{n-1}$ be open sets covering $K$, each meeting $K$, with $U_i\cap U_j=\emptyset$ for $\abs{i-j}\ge2$. Then
$U_i\cap U_{i+1}\neq\emptyset$ for every $i$; so $(U_0,\dots,U_{n-1})$ is a chain.
\end{lemma}

\begin{proof}
Otherwise $U_0\cup\dots\cup U_i$ and $U_{i+1}\cup\dots\cup U_{n-1}$ are disjoint open sets covering $K$ and both meeting it.
\end{proof}

\formal{\lean{HPC.BJ59.adjacent_meet_of_cover}.}

\begin{theorem}[Shrinking to a BJ-chain]\label{thm:shrink}
Let $K$ be a continuum, $D$ a chain covering $K$ with every link meeting $K$, and $C_i\subseteq D_i$ compact sets. There is a BJ-chain $D'$ with the same
number of links covering $K$, every link meeting $K$, with $\cl{D'_i}\subseteq D_i$ and $C_i\subseteq D'_i$ for every $i$.
\end{theorem}

\begin{proof}
Choose $k_i\in K\cap D_i$ and enlarge $C_i$ to $C_i\cup\set{k_i}$. The compact set $K$ is covered by the open sets $D_i$, so it is covered by compact sets
$F_i\subseteq D_i$ (the shrinking lemma for finite open covers of a compact metric space). Let $Q_i=F_i\cup C_i$, a compact subset of $D_i$. For
$\abs{i-j}\ge2$, $Q_i$ and $Q_j$ are disjoint compact sets, so they have positive distance; let $d>0$ be smaller than all these distances and smaller than
the distance from each $Q_i$ to the complement of $D_i$. Put $D'_i=\set{x\in D_i:\dist(x,Q_i)<d/3}$. Then $\cl{D'_i}\subseteq\set{x:\dist(x,Q_i)\le
d/3}\subseteq D_i$, $C_i\subseteq D'_i$, the $D'_i$ cover $K$ and each meets $K$ (at $k_i$), and for $\abs{i-j}\ge2$ points of $D'_i$ and $D'_j$ are at distance
at least $d/3$. By Lemma~\ref{lem:adjacent-meet}, $D'$ is a chain.
\end{proof}

\formal{\lean{HPC.BJ59.exists_bjChain_shrink}, from \lean{HPC.BJ59.exists_bjChain_shrink'}. The argument is Barrett's, as used for taut chains in
\lean{HPC.Plane.exists_taut_epsChain}.}

\section{Private points}

\begin{definition}\label{def:private}
A chain $D$ has \emph{private points} if every link $D_j$ contains a point lying in the closure of no other link.
\end{definition}

\formal{\lean{HPC.BJ59.HasPrivatePoints}.}

\begin{proposition}\label{prop:private-interior}
Let $D$ be a BJ-chain covering a connected set $P$, with every link meeting $P$. Every interior link $D_i$ ($0<i<n-1$) contains a point of $P$ outside the
closures of the other links.
\end{proposition}

\begin{proof}
Otherwise every point of $P\cap D_i$ lies in the closure of another link. Let $A$ be the union of the closures of the links $D_j$, $j<i$, and $B$ that of
the links $D_j$, $j>i$. They are closed, and disjoint by the BJ property ($j<i<k$ forces $\abs{j-k}\ge2$). $P\subseteq A\cup B$, and $P$ meets both (at
points of $D_{i-1}$ and $D_{i+1}$), so $P$ is not connected.
\end{proof}

\formal{\lean{HPC.BJ59.exists_privatePoint_interior}.}

\begin{proposition}\label{prop:private-shrink}
Let $K$ be a nondegenerate continuum and $D$ a BJ-chain covering $K$ with every link meeting $K$. There is a BJ-chain $D'$ with $D'_i\subseteq D_i$ covering
$K$ and having private points in $K$.
\end{proposition}

\begin{proof}
A nondegenerate connected set has no isolated points, so we may choose distinct points $z_i\in K\cap D_i$. Choose $r>0$ with $B(z_i,r)\subseteq D_i$ and
$2r<\dist(z_i,z_j)$ for all $i\neq j$, and put $D'_i=D_i\setminus\bigcup_{j\ne i}\cl B(z_j,r/2)$. A point of $D_i$ removed from $D'_i$ lies in some
$\cl B(z_j,r/2)\subseteq D_j$ and in no other of these balls, so it lies in $D'_j$; hence $D'$ covers whatever $D$ covers. $z_i\in D'_i$, and $z_i$ is not in
the closure of $D'_j$ for $j\neq i$, since $D'_j$ misses $B(z_i,r/2)$. Subsets of the links of a BJ-chain keep the distance condition, and
Lemma~\ref{lem:adjacent-meet} shows $D'$ is a chain.
\end{proof}

\formal{\lean{HPC.BJ59.exists_privatePoint_shrink}, with \lean{HPC.BJ59.exists_injective_choice}.}

\section{Runs}

\begin{lemma}\label{lem:runs}
Let $D$ be a chain with $n$ links.
\begin{enumerate}[label=(\alph*)]
\item If $C\neq\emptyset$ is covered by $D$, some run $[r,s]$ covers $C$ while neither $[r+1,s]$ nor $[r,s-1]$ does.
\item If $S\neq\emptyset$ is connected and covered by $D$, the links meeting $S$ form a run covering $S$.
\item If $C$ is connected and the run $[r,s]$ covers $C$ minimally as in (a), each link $D_j$, $r\le j\le s$, contains a point of $C$ lying in no other
link of the run.
\end{enumerate}
\end{lemma}

\begin{proof}
(a) Take a shortest covering run. (b) If $k$ lies strictly between two indices of links meeting $S$ and $D_k$ misses $S$, the unions of the links before and
after $k$ separate $S$. (c) Otherwise $C\cap D_j$ lies in the union $U$ of the links $D_k$, $r\le k<j$, together with the union $V$ of the links $D_k$,
$j<k\le s$. These are disjoint open sets covering $C$. If $C$ missed $U$, the run $[r+1,s]$ would cover $C$ (if $j>r$, as $D_r\subseteq U$; if $j=r$, as $C\subseteq
V$); so $C$ meets $U$, likewise $V$, and $C$ is not connected.
\end{proof}

\formal{\lean{HPC.BJ59.exists_minimal_run}, \lean{HPC.BJ59.exists_run_meeting}, \lean{HPC.BJ59.exists_point_only_of_minimal}.}

\section{Chain maps}

\begin{definition}[Bing and Jones, p.~173]\label{def:chainmap}
A map $H\colon\set{0,\dots,n-1}\to\set{0,\dots,m-1}$ is a \emph{chain map} if $\abs{H(i)-H(j)}\le1$ whenever $\abs{i-j}\le1$. If $D$ has $n$ links and $E$ has
$m$ links, we say $H$ is a chain map of $D$ into (onto, if surjective) $E$, thinking of $H$ as sending $D_i$ to $E_{H(i)}$; containments $D_i\subseteq
E_{H(i)}$ are stated separately when they hold.
\end{definition}

\formal{\lean{HPC.BJ59.IsChainMap}.}

\begin{lemma}[Discrete intermediate values]\label{lem:chainmap-ivt}
If $H$ is a chain map and $i\le i'$, every value between $H(i)$ and $H(i')$ is taken by $H$ on $[i,i']$. So the image of a run is a run.
\end{lemma}

\formal{\lean{HPC.BJ59.chainMap_ivt}, from \lean{HPC.BJ59.exists_eq_of_steps}.}

Bing and Jones's Theorem~6 says that two chain maps into the same chain, each reaching its next-to-last link, agree up to adjacency somewhere. The proof
of Theorem~10 uses the following form instead.

\begin{theorem}[Theorem~6, the form used]\label{thm:bj6}
Let $D$ be a chain, $C$ a set, and $R,R'$ chain maps defined on a run $[t,u]$ of indices, with values in the indices of $D$. Suppose each link $D_{R(a)}$,
$t\le a\le u$, meets $C$ and these links cover $C$; and the same for $R'$. Then $\abs{R(a)-R'(a)}\le1$ for some $a\in[t,u]$.
\end{theorem}

\begin{proof}
Let $a_0$ minimize $R$ on $[t,u]$ and $b_0$ minimize $R'$; put $v=R(a_0)$, $v'=R'(b_0)$. A point of $D_v\cap C$ lies in some $D_{R'(a)}$, which then meets $D_v$,
so $v'\le R'(a)\le v+1$. Symmetrically $v\le v'+1$. The function $\varphi=R-R'$ changes by at most $2$ between consecutive indices. Now $\varphi(a_0)=v-R'(a_0)\le
v-v'\le1$ and $\varphi(b_0)=R(b_0)-v'\ge v-v'\ge-1$. If $\varphi(a_0)\ge-1$ we are done at $a_0$. Otherwise $\varphi(a_0)\le-2$ and $\varphi(b_0)\ge-1$; walking
from $a_0$ towards $b_0$, the first index where $\varphi\ge-1$ follows one where $\varphi\le-2$, so there $-1\le\varphi\le0$.
\end{proof}

\formal{\lean{HPC.BJ59.adjacent_of_properCover}; Bing and Jones's form is \lean{HPC.BJ59.bj_theorem6} (with \lean{HPC.BJ59.sign_persist}), which is
not used.}

\section{Folding a chain along a walk}

\begin{definition}\label{def:fold}
Let $E$ be a chain with $N$ links and $\pi\colon\set{0,\dots,N-1}\to\set{0,\dots,m-1}$ a chain map with a section $\mu$ ($\pi\circ\mu=\mathrm{id}$) that is
also a chain map. The \emph{folded chain} has links $\bigcup_{\pi(k)=i}E_k$, $i<m$.
\end{definition}

It is a chain: if folded links $i$ and $j$ meet, some $E_k$ and $E_l$ with $\pi(k)=i$, $\pi(l)=j$ meet, so $\abs{k-l}\le1$ and $\abs{i-j}\le1$; conversely
$E_{\mu(i)}$ and $E_{\mu(j)}$ meet when $\abs{i-j}\le1$. If $E$ is a BJ-chain so is the folded chain (finitely many pairs of links, each at positive
distance).

\formal{\lean{HPC.BJ59.foldChain}, \lean{HPC.BJ59.isBJChain_foldChain}.}

\begin{lemma}[The unfolded walk]\label{lem:unfold}
Let $H$ be a chain map of $\set{0,\dots,m-1}$ onto $\set{0,\dots,n-1}$ with $H(r)=0$, $H(s)=n-1$, $r\le s$. The walk
\[
r,\ r-1,\ \dots,\ 1,\ 0,\ 1,\ \dots,\ m-1,\ m-2,\ \dots,\ s
\]
has steps of at most one and passes through every index; composed with $H$ it is a pattern from the first link of $D$ to the last (\vtwo{Chapter~3}).
\end{lemma}

\formal{\lean{HPC.BJ59.unfoldIdx}, \lean{HPC.BJ59.unfoldFin}, \lean{HPC.BJ59.unfoldMain}, \lean{HPC.BJ59.isPattern_unfold}.}

Thus if a chain $E'$ follows the pattern $H\circ(\text{walk})$ in a chain $D$, folding $E'$ along the walk (with the middle pass $0,1,\dots,m-1$ as the section)
gives a chain with $m$ links whose $k$-th link lies in $D_{H(k)}$: a chain that \emph{follows} $H$. This is Bing and Jones's device
$e_i=e'_i+f_i+g_i$ in the proof of their Theorem~7.

\section{Fibres and fine chains}

\begin{lemma}\label{lem:fibre-spread}
Let $X$ be compact and $f\colon X\to\I$ continuous and open, with nondegenerate fibres. There is $\kappa>0$ such that every fibre contains two points at
distance at least $\kappa$.
\end{lemma}

\begin{proof}
For each $t$ choose $p_t\ne q_t$ in $\fib ft$ and let $d_t=\dist(p_t,q_t)$. The set $O_t=f(B(p_t,d_t/3))\cap f(B(q_t,d_t/3))$ is open (as $f$ is open) and
contains $t$. Finitely many $O_t$ cover $\I$; let $\kappa$ be the least of the corresponding $d_t/3$. If $s\in O_t$, the fibre over $s$ has points within
$d_t/3$ of $p_t$ and of $q_t$, at distance at least $d_t/3\ge\kappa$.
\end{proof}

\formal{\lean{HPC.BJ59.exists_fibre_spread}.}

\begin{definition}\label{def:fibre-fine}
A chain $E$ on $Y$ is \emph{fibre-fine} for $g\colon Y\to\I$ if it has one link, or every run covering a fibre $\fib gt$ has at least three links.
\end{definition}

\formal{\lean{HPC.BJ59.IsFibreFine}.}

\begin{lemma}\label{lem:fibre-fine}
If every fibre of $g$ contains two points at distance at least $\kappa$, every chain of mesh $<\kappa/3$ is fibre-fine.
\end{lemma}

\begin{proof}
If a run of at most two links covered $\fib gt$, two points of the fibre at distance $\ge\kappa$ would lie in one link or in two meeting links, hence at
distance $<2\kappa/3$.
\end{proof}

\formal{\lean{HPC.BJ59.isFibreFine_of_isEpsChain}.}

\section*{Exercises}

\begin{exercise}
Give a chain on $\R$ which is not a BJ-chain. (Use open intervals whose nonadjacent members are disjoint but touch at a point.)
\end{exercise}

\begin{exercise}
Show that Proposition~\ref{prop:private-interior} fails for the end links: a BJ-chain of three links properly covering $\I$ may have its first link inside
the closure of the second.
\end{exercise}

\begin{exercise}
Show that the conclusion of Theorem~\ref{thm:bj6} can fail if only one of $R$, $R'$ has all its links meeting $C$.
\end{exercise}

%======================================================================
\chapter{Snake-like arcs of continua}\label{ch:arcs}
%======================================================================

This chapter treats Sections~3--5 of Bing and Jones: chains running from one end of a snake-like arc of continua to the other, open sets that separate the
two ends, and the repaired form of their Theorem~5. Throughout, $X$ is a compact metric space and $f\colon X\to\I$ is an arc of continua
(Definition~\ref{def:arc-continua}); $g_0=\fib f0$ and $g_1=\fib f1$ are the end fibres and $\G ab=\fib f{[a,b]}$, a continuum for $a\le b$
(Corollary~\ref{cor:G-connected}).

\section{Theorem 3: chains from end to end}

\begin{theorem}[Bing and Jones, Theorem~3]\label{thm:bj3}
If $X$ is snake-like, then for every $\varepsilon>0$ there is an $\varepsilon$-chain covering $X$ whose first link meets $g_0$ and whose last link meets
$g_1$.
\end{theorem}

\begin{proof}
We may assume $\varepsilon$ is small compared with the distance from $g_0$ to $g_1$. Take an $\varepsilon$-chain $E$ covering $X$. The links meeting $g_0$ form a run $E^0$ (Lemma~\ref{lem:runs}(b)), with union $V_0$; likewise the links meeting
$g_1$ form a run $E^1$ with union $V_1$; for small $\varepsilon$ these runs are disjoint and not adjacent. Choose $0<a<b<1$ with $\G0a\subseteq V_0$ and $\G
b1\subseteq V_1$ (possible since $f$ is closed). Let $F=(F_0,\dots,F_{N-1})$ be a chain covering $X$ of mesh $\gamma$, smaller than a Lebesgue number of $E$ and
than the distances from $g_0$ to $\G a1$ and from $g_1$ to $\G0b$; orient it so that its links meeting $g_0$ come before those meeting $g_1$.

The links of $F$ before the first link meeting $g_0$ lie in $\G0a$: the continuum $\G a1$ meets $g_1$, and if it met such a link it would meet every link
between it and a link meeting $g_1$, including a link meeting $g_0$, against the choice of $\gamma$. So these links lie in $V_0$. Let $F_{v+1}$ be the first link
after the links meeting $g_0$ that is not contained in $V_0$, and $F_{w-1}$ the last link before the links meeting $g_1$ that is not contained in $V_1$. Then
$F_v$ lies in a link of $E^0$ and meets a set outside $V_0$, so it lies in an end link of $E^0$; orient $E^0=(e^0_0,\dots,e^0_{n_0})$ so that $F_v\subseteq
e^0_{n_0}$. Likewise $F_w$ lies in the first link $e^1_0$ of $E^1=(e^1_0,\dots,e^1_{n_1})$. For small $\gamma$, $w\ge v+4$. Put $A=F_0\cup\dots\cup F_v$ and
$B=F_w\cup\dots\cup F_{N-1}$; then $A\subseteq V_0$ and $B\subseteq V_1$. The sequence
\[
A\cap e^0_0,\ \dots,\ A\cap e^0_{n_0-1},\ (A\cap e^0_{n_0})\cup F_{v+1},\ F_{v+2},\ \dots,\ F_{w-2},\ (B\cap e^1_0)\cup F_{w-1},\ B\cap e^1_1,\ \dots,\ B\cap
e^1_{n_1}
\]
covers $X$, has mesh $<\varepsilon$, starts with a link meeting $g_0$ and ends with one meeting $g_1$. It is a chain: the only nontrivial point is that
$F_{v+1}$ meets $A$ only inside $F_v\subseteq e^0_{n_0}$, so merged into the last $A$-link it creates no new intersections; likewise at the other end. Empty links
are removed.
\end{proof}

\formal{\lean{HPC.BJ59.theorem3_arc} (with explicit hypotheses), \lean{HPC.BJ59.bj_theorem3}; the chain is \lean{HPC.BJ59.chain_assembly}.}

\begin{remark}
Bing and Jones print the chain with $A\cap e^0_{n_0}$ and $F_{v+1}$ as two separate links. Then $A\cap e^0_{n_0-1}$ may meet $F_{v+1}$ (both can meet $F_v$), and
the sequence is not a chain. Merging the two links repairs it, at both junctions.
\end{remark}

Bing and Jones's Theorem~2 (for $0<a<b<c<d<1$, in every fine enough chain covering $X$, links meeting $\G0a$, $\G bc$ and $\G d1$ occur in this order or the
reverse) is also proved, but not used below.

\formal{\lean{HPC.BJ59.theorem2_arc}, \lean{HPC.BJ59.bj_theorem2}.}

\section{Theorem 4: open sets separate the ends}

\begin{theorem}[Bing and Jones, Theorem~4]\label{thm:bj4}
Let $X$ be snake-like and $U\neq\emptyset$ an open set missing $g_0\cup g_1$. Then $X\setminus U=A\cup B$ with $A$, $B$ disjoint and closed, $g_0\subseteq A$
and $g_1\subseteq B$.
\end{theorem}

\begin{proof}
Pick $p\in U$ and, by Theorem~\ref{thm:bj3}, a chain $D=(D_0,\dots,D_m)$ covering $X$ from $g_0$ to $g_1$ with mesh less than the distance from $p$ to $X\setminus U$.
The link $D_l$ containing $p$ lies in $U$. Let $A=(X\setminus U)\setminus(D_{l+1}\cup\dots\cup D_m)$ and $B=(X\setminus U)\setminus(D_0\cup\dots\cup D_{l-1})$. They
are closed. Since $D_l\subseteq U$, every point of $X\setminus U$ lies in a link before or after $l$, and no point lies in both (such links would be
nonadjacent), so $A\cup B=X\setminus U$ and $A\cap B=\emptyset$. The connected set $g_0$ misses $U$, so it lies in $A$ or in $B$; it meets $D_0$, and $l\ne0$
since $D_0$ meets $g_0$, which misses $U$; so $g_0\subseteq A$. Likewise $g_1\subseteq B$.
\end{proof}

\formal{\lean{HPC.BJ59.bj_theorem4}, from \lean{HPC.BJ59.theorem4_of_chain} and \lean{HPC.BJ59.theorem4_explicit}.}

\begin{corollary}[Located form]\label{cor:bj-cut}
Let $X$ be snake-like, $0<\alpha<\beta<1$, and $V\neq\emptyset$ open with $\alpha<f<\beta$ on $V$. Then $X\setminus V=L\cup R$ with $L$, $R$ disjoint and
closed, $f<\beta$ on $L$ and $f>\alpha$ on $R$.
\end{corollary}

\begin{proof}
Take $L=A$, $R=B$ from Theorem~\ref{thm:bj4}. The continuum $\G\beta1$ misses $V$ and contains $g_1$, so it lies in $B$; hence $f<\beta$ on $A$. Likewise
$\G0\alpha\subseteq A$.
\end{proof}

\formal{\lean{HPC.BJ59.bj_cut}.}

Two compactness facts about open sets meeting every fibre will be needed in Chapter~\ref{ch:thm78}.

\begin{lemma}[Deep points]\label{lem:deep}
Let $f$ be open and let $d$ be an open set meeting every fibre. There is $\rho>0$ such that every fibre contains a point $y$ with $B(y,3\rho)\subseteq d$.
If moreover $d\neq X$, there are open sets $U$, $U'$ with $\cl U\subseteq d$, $\cl{U'}\subseteq U$, and $U'$ meeting every fibre.
\end{lemma}

\begin{proof}
For each $t$ choose $y_t\in\fib ft$ and $r_t>0$ with $B(y_t,r_t)\subseteq d$. The open sets $f(B(y_t,r_t/2))$ cover $\I$; take a finite subcover and let $3\rho$
be the least corresponding $r_t/2$. A fibre over a point of $f(B(y_t,r_t/2))$ contains a point $y$ within $r_t/2$ of $y_t$, and $B(y,r_t/2)\subseteq d$. For the
second statement let $U'=\set{y:\dist(y,X\setminus d)>2\rho}$ and $U=\set{y:\dist(y,X\setminus d)>\rho}$; the deep points lie in $U'$.
\end{proof}

\formal{\lean{HPC.BJ59.exists_deep_radius}, \lean{HPC.BJ59.exists_deep_sets}.}

\section{Theorem 5}

Bing and Jones's Theorem~5 reads: if the fibres are indecomposable, $E$ is a chain properly covering each fibre (every link meets every fibre), and $U$ is a
nonempty open set missing $g_0\cup g_1$ whose closure lies in a link of $E$, then the decomposition $X\setminus U=A\cup B$ of Theorem~4 can be chosen so that,
for every $t$, $A\cap\fib ft$ meets every link of $E$ or $B\cap\fib ft$ does.

\begin{proposition}[Theorem~5 is false as printed]\label{prop:bj5-false}
Let $f\colon X\to\I$ be continuous and $E$ a chain of $n\ge3$ links covering $X$, each link meeting every fibre. Merging the last link into the
next-to-last gives a chain $E'$ with the same properties, and there is an open $U$ satisfying the hypotheses of Theorem~5 for $E'$, containing no fibre,
such that some fibre meets some link of $E'$ only inside $U$. For this $E'$ and $U$ the conclusion of Theorem~5 fails for every decomposition $X\setminus
U=A\cup B$.
\end{proposition}

\begin{proof}
Let $E'=(e_0,\dots,e_{n-3},e_{n-2}\cup e_{n-1},e_{n-1})$. Let $K$ be the closure of $\fib f{1/2}\cap e_{n-1}$; it is compact, lies in $e_{n-2}\cup e_{n-1}$ and
misses $g_0\cup g_1$. For small $\delta>0$ let $U$ be the $\delta$-neighbourhood of $K$ minus $\cl{e_0}$. Then $\cl U$ lies in the merged link and misses
$g_0\cup g_1$; $U$ contains $\fib f{1/2}\cap e_{n-1}$ minus $\cl{e_0}$, which (as $n\ge3$) is all of it; and $U$ misses $e_0$, which meets every fibre, so no
fibre lies in $U$. The fibre over $1/2$ meets the last link of $E'$ only inside $U$, so neither $A\cap\fib f{1/2}$ nor $B\cap\fib f{1/2}$ meets it.
\end{proof}

\formal{\lean{HPC.BJ59.bj_theorem5_printed_false_general}; a concrete instance on $\I$ is \lean{HPC.BJ59.bj_theorem5_printed_false}.}

The repair replaces the hypothesis on $\cl U$ by one that excludes exactly this situation.

\begin{theorem}[Theorem~5, repaired]\label{thm:bj5}
Let $X$ be snake-like, $f$ open, and the fibres indecomposable. Let $E$ be a chain and $U\ne\emptyset$ an open set missing $g_0\cup g_1$, such that no fibre meets
a link of $E$ only inside $\cl U$. Then there is a decomposition $X\setminus U=A\cup B$ as in Theorem~4 such that for every $t$, $A\cap\fib ft$ meets every link
of $E$, or $B\cap\fib ft$ does.
\end{theorem}

\begin{proof}
Let $p\in U$, and for $k\ge1$ let $D^k$ be chains as in Theorem~\ref{thm:bj3} with mesh $\delta_k\to0$ and a link $d_k\ni p$ inside $U$, and $A_k$, $B_k$ the
decomposition built from $D^k$ in the proof of Theorem~\ref{thm:bj4}. Suppose every $A_k,B_k$ fails: there are $t_k$ and links $e_{\alpha_k}$, $e_{\beta_k}$ of
$E$ with $A_k\cap\fib f{t_k}\cap e_{\alpha_k}=\emptyset$ and $B_k\cap\fib f{t_k}\cap e_{\beta_k}=\emptyset$. Passing to a subsequence, $\alpha_k=\alpha$,
$\beta_k=\beta$ and $t_k\to c$.

Let $P_k$ be $\fib f{t_k}$ minus the links after $d_k$, and $Q_k$ the fibre minus the links before $d_k$. They are closed, cover the fibre, and meet only inside
$d_k$. By clopen boundary bumping, $P_k\cup\set p$ is $\delta_k$-chained, and so is $Q_k\cup\set p$. Their topological limits superior $P$, $Q$ are therefore
continua; they lie in $\fib fc$ by continuity of $f$, and they cover $\fib fc$ because $f$ is open. As $\fib fc$ is indecomposable, $P=\fib fc$ or $Q=\fib
fc$; say $P$. By hypothesis $\fib fc\cap e_\alpha$ has a point $q\notin\cl U$. Near $q$, for some large $k$, $P_k$ meets $e_\alpha$ outside $U$, at a point of
$A_k\cap\fib f{t_k}\cap e_\alpha$: a contradiction. The case $Q=\fib fc$ is the same with $e_\beta$.
\end{proof}

\formal{\lean{HPC.BJ59.theorem5_repaired}, \lean{HPC.BJ59.bj_theorem5}; the chained parts are \lean{HPC.BJ59.chainedAt_union_point}. The hypotheses
``$\cl U$ in a link'' and ``$E$ covers $X$'' are not used.}

\section{Windows}

In the proof of Theorem~10, Theorem~5 is applied not to the whole arc but to a window $\G ab$, and to a run of a chain.

\begin{lemma}[Transfer to a window]\label{lem:window-transfer}
If $f$ is a continuous arc of pseudo-arcs on the snake-like compact space $X$ and $a<b$, then $\G ab$ with the map $x\mapsto(f(x)-a)/(b-a)$ is again a
continuous arc of pseudo-arcs on a snake-like compact space, and its fibre over $s$ is $\fib f{a+s(b-a)}$.
\end{lemma}

\formal{\lean{HPC.BJ59.window_transfer}, with \lean{HPC.BJ59.winScale}, \lean{HPC.BJ59.winInv}.}

\begin{theorem}[Theorem~5 on a window]\label{thm:bj5-window}
Let $f$ be a continuous arc of pseudo-arcs on a snake-like compact $X$, $a<b$, $F$ a chain and $[t,u]$ a run of $F$ covering $\G ab$ whose links meet every fibre
over $[a,b]$. Let $U\neq\emptyset$ be open with $a<f<b$ on $U$, $\cl U$ inside a link of the run, and suppose no fibre over $[a,b]$ meets a link of the run only
inside $\cl U$. Then $X\setminus U=A\cup B$ with $A$, $B$ disjoint and closed, $\G0a\subseteq A$, $\G b1\subseteq B$, and for every $s\in[a,b]$ either
$A\cap\fib fs$ meets every link of the run, or $B\cap\fib fs$ does.
\end{theorem}

\begin{proof}
Apply Theorem~\ref{thm:bj5} on the window (Lemma~\ref{lem:window-transfer}) to the run, traced on $\G ab$, and add $\G0a$ to $A$ and $\G b1$ to $B$.
\end{proof}

\formal{\lean{HPC.BJ59.bj_theorem5_window}, \lean{HPC.BJ59.bj_theorem5_window_aux}.}

\begin{lemma}\label{lem:window-set}
In the situation of Theorem~\ref{thm:bj5-window}, for every link $F_y$ of the run there is an open $U\ne\emptyset$ with $a<f<b$ on $U$ and $\cl U\subseteq F_y$
such that no fibre over $[a,b]$ meets a link of the run only inside $\cl U$.
\end{lemma}

\begin{proof}[Proof in outline]
Take $U$ a small ball about a point of $F_y$ over $(a+b)/2$, so small that its image lies in a short interval about $(a+b)/2$ and its diameter is less than a
lower bound, obtained by compactness and openness of $f$ as in Lemma~\ref{lem:fibre-spread}, for the size of the traces of the links of the run on the fibres
over that interval.
\end{proof}

\formal{\lean{HPC.BJ59.exists_window_set}.}

\section{Step 1(b): short sub-arcs}

\begin{lemma}\label{lem:meets-all}
A connected set covered by the run $[u,v]$ of a chain and meeting $D_u$ and $D_v$ meets every link of the run.
\end{lemma}

\formal{\lean{HPC.BJ59.meets_all_of_meets_ends}.}

\begin{theorem}[Bing and Jones, p.~182, Step~1(b)]\label{thm:step1b}
Let $f$ be open and $E$ a chain covering $X$. There is $\delta>0$ such that whenever $b-a\le\delta$ and the run $[u,v]$ covers $\G ab$ while neither $[u+1,v]$ nor
$[u,v-1]$ does, every link of the run meets every fibre over $[a,b]$.
\end{theorem}

\begin{proof}
\emph{Step 1.} About each $c\in\I$ there is an open $N_c\ni c$ such that every point over $N_c$ lies in a link meeting every fibre over $N_c$. Indeed, for each of
the finitely many links $E_w$ meeting $\fib fc$, $f(E_w)$ is an open set containing $c$ ($f$ is open); and the union of these links is an open set containing
$\fib fc$, so it contains $\fib fx$ for $x$ near $c$ ($f$ is closed). Intersect these finitely many neighbourhoods of $c$.

\emph{Step 2.} Let $2\delta$ be a Lebesgue number of the cover $\set{N_c}$, so $[a,b]\subseteq N=N_c$ for some $c$. Let $x\in[a,b]$; we show $\fib fx$ meets $E_u$.
Since $[u+1,v]$ does not cover $\G ab$, some $z\in\G ab$ lies in $E_u$ and in no other link of the run. By Step~1, $z\in E_w$ for a link $E_w$ meeting every fibre
over $N$, so $\abs{w-u}\le1$. If $w=u$ we are done. If $w=u+1\le v$, $z$ would lie in another link of the run. So $w=u-1$. Then $E_{u-1}$ meets $\fib fx$ at a
point $y\in\G ab$, which lies in a link of the run adjacent to $E_{u-1}$, that is, in $E_u$. Likewise $\fib fx$ meets $E_v$, and Lemma~\ref{lem:meets-all}
finishes the proof.
\end{proof}

\formal{\lean{HPC.BJ59.bj_step1b}, with Step~1 as \lean{HPC.BJ59.exists_nhd_step1b}.}

\section*{Exercises}

\begin{exercise}
Show that Theorem~\ref{thm:bj4} fails for $X=\I\times\I$ with the first projection if ``snake-like'' is dropped: an open disc in the middle does not separate the
two ends.
\end{exercise}

\begin{exercise}
Check that the sets $A$ and $B$ in the proof of Theorem~\ref{thm:bj4} are closed.
\end{exercise}

\begin{exercise}
Write out Proposition~\ref{prop:bj5-false} for $X=\I$, $f=\mathrm{id}$ and a chain of three links. (Here every fibre is a point.)
\end{exercise}

%======================================================================
\chapter{Theorems 7 and 8}\label{ch:thm78}
%======================================================================

A chain map $H$ of a chain $Y$ with $m$ links onto a chain $D$ with $n$ links says how the links of a finer chain should run through the links of $D$.
Theorem~7 realizes $H$ on one pseudo-arc and Theorem~8 on a whole snake-like arc of pseudo-arcs. We say a chain $E=(E_0,\dots,E_{m-1})$ \emph{follows} $H$ in
$D$ if $E_i\subseteq D_{H(i)}$ for every $i$.

\section{Theorem 7 on one pseudo-arc}

\begin{theorem}[Bing and Jones, Theorem~7, first sentence]\label{thm:bj7}
Let $P$ be a pseudo-arc in a metric space $X$, $D$ a chain with $n$ links covering $P$ with every link meeting $P$, and $H$ a chain map of $\set{0,\dots,m-1}$
onto $\set{0,\dots,n-1}$. There is a BJ-chain $E$ following $H$ in $D$, covering $P$, with every link meeting $P$.
\end{theorem}

\begin{proof}[Proof in outline]
\emph{Inside $P$.} The traces $D_j\cap P$ form a chain of the compact space $P$. Choose $r$, $s$ with $H(r)=0$, $H(s)=n-1$; reversing the numbering of $D$ if
necessary, $r\le s$.
\begin{enumerate}[label=(\alph*)]
\item Shrink the trace chain, link by link, to a chain $T$ from a point $p$ to a point $q$ ($p$ only in $T_0$, $q$ only in $T_{n-1}$) with $p$ and $q$ in different
composants of $P$: choose $p\in T_0\setminus\cl{T_1}$, then $q$ in the open set $T_{n-1}\setminus\cl{T_{n-2}}$ outside the composant of $p$ (composants of a
pseudo-arc are dense and disjoint, \vtwo{Chapter~1}), and remove small closed balls about $p$ and $q$ from the neighbouring links.
\item By \vtwo{Chapter~7}, $P$ has a description from $p$ to $q$; placing $T$ in front gives a description whose first chain is $T$.
\item By pattern following (\vtwo{Theorem~3.1}) the pattern $H\circ w$ in $T$, where $w$ is the walk of Lemma~\ref{lem:unfold}, is followed by some chain $E'$ of
the description: $E'$ covers $P$ and its $k$-th link lies in $T_{H(w(k))}$.
\item Fold $E'$ along $w$ (Definition~\ref{def:fold}): the result is a chain of $P$ with $m$ links following $H$ in $T$.
\end{enumerate}
\emph{Back in $X$.} Each link of the folded chain is open in $P$, so it is $P\cap O_i$ for an open $O_i\subseteq D_{H(i)}$ of $X$; the $O_i$ need not form a chain of
$X$, but their traces on $P$ do, and the shrinking argument of Theorem~\ref{thm:shrink} turns them into a BJ-chain of $X$ inside the $O_i$.
\end{proof}

\formal{\lean{HPC.BJ59.bj_theorem7'} and \lean{HPC.BJ59.bj_theorem7}; inside $P$, \lean{HPC.BJ59.bj7_compact}, with \lean{HPC.BJ59.traceChain},
\lean{HPC.BJ59.exists_chain_isFrom}, \lean{HPC.BJ59.prependDesc}, \lean{HPC.BJ59.exists_follow_pattern}; the lift uses
\lean{HPC.BJ59.exists_bjChain_shrink'}.}

Bing and Jones assume the description of step~(b) without comment. The second sentence of their Theorem~7 adds prescribed sets in two links.

\begin{theorem}[Theorem~7, second sentence]\label{thm:bj7-ends}
In Theorem~\ref{thm:bj7} suppose $n\ge3$, $H(r)=0$, $H(s)=n-1$, and let $A\subseteq P\cap D_0$ and $B\subseteq P\cap D_{n-1}$ be closed. Then $E$ may be chosen with
moreover $A\subseteq E_r$ and $B\subseteq E_s$.
\end{theorem}

\begin{proof}[Proof in outline]
Work in $P$ with the trace chain $T$, and assume $A$, $B$ nonempty.
\emph{The collar chain.} Choose $\rho>0$ with $\set{x:\dist(x,A)\le3\rho}\subseteq T_0$ and a point $w\in T_0\cap T_1$ with $\dist(w,A)>3\rho$ (if $T_0\cap T_1\subseteq
A$, then $T_0$ would be closed and open in $P$). Replace $T_0$, $T_1$ by the four sets
\[
\set{\dist(\cdot,A)<2\rho},\quad \set{\rho<\dist(\cdot,A)<3\rho},\quad T_0\cap\set{\dist(\cdot,A)>2\rho},\quad T_1\cap\set{\dist(\cdot,A)>3\rho},
\]
and likewise at the other end for $B$. This gives a chain $M$ of $n+4$ links refining $T$ by $(0,0,0,1,2,\dots,n-2,n-1,n-1,n-1)$, in which $A$ meets only the first
link and $B$ only the last.
\emph{The buffered pattern.} Extend the walk $w$ by two more copies of $r$ in front and of $s$ at the end, and let the pattern in $M$ be $0,1$, then $H(w(k))+2$ along
the walk, then $n+2,n+3$. Follow it in $M$ (as in steps (a)--(c) above) and fold along the extended walk. Every link of the following chain that meets $A$ lies in
$M_0$, so it sits at a position of the pattern with value $0$, which the walk sends to $r$; hence $A$ lies in the $r$-th folded link. Likewise $B$ lies in the
$s$-th. Lift to $X$ as before, keeping the compact sets $A$, $B$ inside the links (Theorem~\ref{thm:shrink}).
\end{proof}

\formal{\lean{HPC.BJ59.bj_theorem7_ends'}; inside $P$, \lean{HPC.BJ59.bj7_ends_compact}, with \lean{HPC.BJ59.collar_data},
\lean{HPC.BJ59.exists_collar_chain}, \lean{HPC.BJ59.bufIdx}, \lean{HPC.BJ59.bufPat}.}

\begin{remark}\label{rem:bj7-false}
As printed, the second sentence of Theorem~7 has no restriction on $n$. It fails for $n=1$: take $D$ one link, $m=3$, $A=B=\set p$, $r=0$, $s=2$; then $p$ would lie
in the two nonadjacent links $E_0$ and $E_2$. The printed proof replaces $D$ by a chain $U_0,U_1,U_2,\dots$ built from the differences $D_1\setminus D_0$, $D_2\setminus
D_1$; for $n=3$ two of these sets, whose places are not adjacent, must meet. The collar chain above avoids this and works for every $n\ge3$.
\end{remark}

\formal{\lean{HPC.BJ59.bj_theorem7_ends_printed_false}, \lean{HPC.BJ59.bj_theorem7_proof_n3}.}

\section{Theorem 8 on an arc of pseudo-arcs}

From now on $X$ is a snake-like compact metric space and $f\colon X\to\I$ a continuous arc of pseudo-arcs (open, with pseudo-arc fibres).

\begin{theorem}[Bing and Jones, Theorem~8]\label{thm:bj8}
Let $D$ be a chain of $n\ge3$ links covering $X$, each link meeting every fibre, and $H$ a chain map of $\set{0,\dots,m-1}$ onto $\set{0,\dots,n-1}$. There is a BJ-chain
$E$ covering $X$ and following $H$ in $D$, each link of $E$ meeting every fibre.
\end{theorem}

\begin{proof}
\emph{The case $H(0)=0$, $H(m-1)=n-1$.}

\emph{Step 1: $U_A$, $U_B$.} By Lemma~\ref{lem:deep} there are $\rho>0$ and open sets $U'_A$, $U_A$ with $\cl{U'_A}\subseteq U_A$, $\cl{U_A}\subseteq D_0$, $U'_A$ meeting every
fibre; likewise $U'_B$, $U_B$ in $D_{n-1}$. As $n\ge3$, $D_0\cap D_{n-1}=\emptyset$.

\emph{Step 2: $E(a)$.} For each $a\in\I$, Theorem~\ref{thm:bj7-ends} on the pseudo-arc $p_a=\fib fa$ with $r=0$, $s=m-1$, $A=\cl{U_A}\cap p_a$, $B=\cl{U_B}\cap p_a$ gives a chain
$E(a)$ following $H$ in $D$, covering $p_a$, with $\cl{U_A}\cap p_a\subseteq E(a)_0$ and $\cl{U_B}\cap p_a\subseteq E(a)_{m-1}$.

\emph{Step 3: $I_a$.} The set of points of $X$ not covered by $E(a)$, together with the points of $\cl{U_A}$ outside $E(a)_0$ and those of $\cl{U_B}$ outside $E(a)_{m-1}$, is
closed and misses $p_a$; its image is closed and misses $a$. Let $I_a$ be the complement of that image: an open neighbourhood of $a$ such that for $x\in I_a$,
$p_x\subseteq\bigcup E(a)$, $\cl{U_A}\cap p_x\subseteq E(a)_0$ and $\cl{U_B}\cap p_x\subseteq E(a)_{m-1}$.

\emph{Step 4: the cover.} Let $\lambda$ be a Lebesgue number of $\set{I_a}$ and cut $\I$ into $T$ segments so short that each segment, enlarged slightly, lies in some
$I_{a_i}$ ($i=1,\dots,T$). At each of the $T-1$ cuts $c_i$ choose a short open interval $J_i\ni c_i$ inside $(0,1)$, the $J_i$ disjoint and increasing. The set
$V_i=(U_A\cup U_B)\cap\fib f{J_i}$ is open and nonempty and lies over $J_i$, so Corollary~\ref{cor:bj-cut} splits $X\setminus V_i$ into disjoint closed sets $L_i$ (over
$[0,\sup J_i)$) and $R_i$ (over $(\inf J_i,1]$). Put $R_0=L_T=X$ and
\[
U_i=(R_{i-1}\cap L_i)\cup U_A\cup U_B=\bigl((X\setminus L_{i-1})\cap(X\setminus R_i)\bigr)\cup U_A\cup U_B .
\]
The $U_i$ are open and cover $X$ (a point lies in $U_i$ for the least $i$ with the point in $L_i$); for $i<k$, $U_i\cap U_k\subseteq U_A\cup U_B$ since $L_i\cap
R_{k-1}=\emptyset$; and $U_i\setminus(U_A\cup U_B)$ lies over $(\inf J_{i-1},\sup J_i)\subseteq I_{a_i}$.

\emph{Step 5: the links.} Put
\[
e_k=\bigcup_{i=1}^T E(a_i)_k\cap U_i\cap\fib f{I_{a_i}}\qquad(k<m).
\]
They are open, cover $X$, $e_k\subseteq D_{H(k)}$, $U_A\subseteq e_0$ and $U_B\subseteq e_{m-1}$. Nonadjacent links are disjoint: two pieces with the same $i$ are disjoint since
$E(a_i)$ is a chain; pieces with $i\neq j$ meet only in $U_A\cup U_B$; at a point of $U_A$ over $x\in I_{a_i}\cap I_{a_j}$ both $E(a_i)_0$ and $E(a_j)_0$ contain the point,
so both indices are at most $1$; $U_B$ is symmetric.

\emph{Step 6.} Each fibre meets $U_A\subseteq e_0$ and $U_B\subseteq e_{m-1}$ and is connected, so it meets every $e_k$ (Lemma~\ref{lem:meets-all}); in particular adjacent
links meet and $(e_k)$ is a chain.

\emph{Step 7.} Shrink to a BJ-chain by Theorem~\ref{thm:shrink}, keeping the compact sets $\cl{U'_A}$ in the first link and $\cl{U'_B}$ in the last; Step~6 applies
again.

\emph{The general case.} Choose $r$, $s$ with $H(r)=0$, $H(s)=n-1$; reversing, $r\le s$. Apply the first case to the pattern $H\circ w$ along the walk $w$ of
Lemma~\ref{lem:unfold} and fold along $w$. Folding keeps every fibre meeting every link and keeps the BJ property.
\end{proof}

\formal{\lean{HPC.BJ59.bj_theorem8}, from \lean{HPC.BJ59.bj_theorem8_le} and \lean{HPC.BJ59.bj_theorem8_ends}.}

\begin{remark}
Bing and Jones give no argument for $U_A$, $U_B$ (and print $d_2$ for $d_n$), derive the cover $U_1,\dots,U_t$ from Theorem~4 without detail, and take
$e_k=\bigcup_iE(a_i)_k\cap U_i$. These sets need not form a chain: $E(a_i)_k$ is controlled only over $I_{a_i}$, and over a fibre outside $I_{a_i}$ it can reach into
$U_A\subseteq U_i\cap U_j$ and meet $E(a_j)_0$ there. Restricting each piece to $\fib f{I_{a_i}}$ repairs this.
\end{remark}

\section{The corollary to Theorem 8}

\begin{lemma}\label{lem:walk-core}
Let $a_p,\dots,a_q$ be integers with $\abs{a_{j+1}-a_j}\le1$, $a_p\le b+1$ and $a_q\ge b+4$. Then for some $k$ with $p\le k\le q-2$, the three values $a_k,a_{k+1},a_{k+2}$
lie in $[b,b+3]$.
\end{lemma}

\begin{proof}
Let $k$ be the last index $\le q$ with $a_k\le b+1$; then $k<q$, and after $k$ the values are $\ge b+2$, increasing by at most one per step, so $a_{k+1}=b+2$ and
$a_{k+2}\le b+3$; and $a_k=b+1$.
\end{proof}

\formal{\lean{HPC.BJ59.walk_core}.}

\begin{corollary}[Bing and Jones, Corollary to Theorem~8]\label{cor:bj8}
Let $E$ be a chain of $n\ge3$ links covering $X$, each link meeting every fibre, and $H_1$ a chain map of $\set{0,\dots,\ell-1}$ onto $\set{0,\dots,n-1}$. There is
$\varepsilon>0$ such that for every $\varepsilon$-chain $F$ covering $X$ there is a chain map $R$ of $F$ onto $\set{0,\dots,\ell-1}$ with $F_j\subseteq E_{H_1(R(j))}$ for every
$j$, and such that, whenever a run $F_v,\dots,F_w$ covers a fibre, every $d<\ell$ is the value of $R$ at three consecutive indices in $[v,w]$.
\end{corollary}

\begin{proof}
Apply Theorem~\ref{thm:bj8} to the chain map $4i+c\mapsto H_1(i)$ ($c<4$) of $4\ell$ indices: a chain $G$ of $4\ell$ links with $G_{4i+c}\subseteq E_{H_1(i)}$, every link
meeting every fibre. Let $\varepsilon$ be smaller than a Lebesgue number of $G$ and, if $\ell=1$, than a third of the uniform spread $\kappa$ of the fibres
(Lemma~\ref{lem:fibre-spread}). For each link $F_j$ choose $a(j)$ with $F_j\subseteq G_{a(j)}$ (since consecutive links of $F$ meet, $G_{a(j)}$ and $G_{a(j+1)}$ meet, so $a$ changes by at most one
between consecutive indices) and put $R(j)=\fl{a(j)/4}$. A run covering a fibre contains points of that fibre in $G_0$ and in
$G_{4\ell-1}$, so along it $a$ runs from $\le1$ to $\ge4\ell-2$ in steps of at most one; Lemma~\ref{lem:walk-core} with $b=4d$ gives three consecutive indices with
$a\in[4d,4d+3]$, that is, $R=d$.
\end{proof}

\formal{\lean{HPC.BJ59.bj_theorem8_cor}. The irreducibility of $F$ assumed by Bing and Jones is not needed.}

The proof of Theorem~10 uses the corollary on a sub-arc $\G ab$, transferred by Lemma~\ref{lem:window-transfer}; there $E$ may also have a single link, and then $R$ is the
zigzag $0,0,0,1,1,1,\dots,\ell-1,\ell-1,\ell-1,\ell-1,\dots,0,0,0,\dots$, which works once $\varepsilon$ is small compared with the spread of the fibres.

\formal{\lean{HPC.BJ59.bj_theorem8_cor_sub}, \lean{HPC.BJ59.bj_theorem8_cor_sub_three}, \lean{HPC.BJ59.bj_theorem8_cor_sub_one}, \lean{HPC.BJ59.zigN}.}

\section*{Exercises}

\begin{exercise}
Show that the surjectivity of $H$ cannot be dropped from Theorem~\ref{thm:bj7}. (Use Proposition~\ref{prop:private-interior}.)
\end{exercise}

\begin{exercise}
Show that the conclusion of Lemma~\ref{lem:walk-core} fails if $a_q\ge b+4$ is weakened to $a_q\ge b+3$.
\end{exercise}

\begin{exercise}
Why does Step~4 of the proof of Theorem~\ref{thm:bj8} need $J_i\subseteq(0,1)$?
\end{exercise}

%======================================================================
\chapter{The homeomorphism from a zigzag}\label{ch:zigzag}
%======================================================================

The homeomorphism of Theorem~10 is the limit of a ``zigzag'': chains on the two spaces, refining each other alternately through chain maps. This chapter proves
that such a zigzag defines a homeomorphism (Bing and Jones, pp.~188--189), independently of how it was built. It is the analogue of \vtwo{Chapter~5}, where
two descriptions of pseudo-arcs gave a homeomorphism.

\section{The data}

Let $X$, $Y$ be metric spaces and $P\subseteq X$, $Q\subseteq Y$ compact. Suppose that for every $i\ge0$ we are given:
\begin{itemize}
\item a chain $D^i$ in $X$ covering $P$, every link meeting $P$, of mesh $<\varepsilon_i$;
\item a BJ-chain $E^i$ in $Y$ covering $Q$, every link meeting $Q$, of mesh $<\varepsilon_i$;
\item a chain map $H_i$ of $D^i$ onto $E^i$ and a chain map $K_i$ of $E^{i+1}$ into $D^i$;
\end{itemize}
where $\varepsilon_i\to0$, and the two conditions
\begin{description}
\item[\textnormal{(B)}] $D^{i+1}_j\subseteq D^i_{K_i(H_{i+1}(j))}$ for every link $D^{i+1}_j$;
\item[\textnormal{(B$'$)}] $E^{i+1}_j\subseteq E^i_{H_i(K_i(j))}$ for every link $E^{i+1}_j$.
\end{description}
\[
\begin{array}{ccccccc}
D^0 & \longleftarrow & D^1 & \longleftarrow & D^2 & \longleftarrow & \cdots\\
\downarrow{\scriptstyle H_0} & \nwarrow{\scriptstyle K_0} & \downarrow{\scriptstyle H_1} & \nwarrow{\scriptstyle K_1} & \downarrow{\scriptstyle H_2} & & \\
E^0 & \longleftarrow & E^1 & \longleftarrow & E^2 & \longleftarrow & \cdots
\end{array}
\]
For $p\in P$ let
\[
\st_i(p)=\bigcup\set{E^i_{H_i(j)}:\ p\in D^i_j},
\]
the union of the images of the links of $D^i$ containing $p$ (Bing and Jones write $H_id(p,i)$).

\formal{\lean{HPC.BJ59.starImage}.}

\section{The homeomorphism}

\begin{lemma}\label{lem:star-props}
For $p\in P$ and every $i$:
\begin{enumerate}[label=(\alph*)]
\item $\st_i(p)$ is one link of $E^i$ or the union of two adjacent ones; it is nonempty, meets $Q$, and has diameter $<2\varepsilon_i$;
\item $\st_{i+1}(p)\subseteq\st_i(p)$.
\end{enumerate}
\end{lemma}

\begin{proof}
(a) Any two links of $D^i$ containing $p$ meet, so they are adjacent; three links cannot be pairwise adjacent, so $p$ lies in one or two adjacent links. Their images
under the chain map $H_i$ are one link or two adjacent links. The links meet $Q$, and two meeting sets of diameter $<\varepsilon_i$ have a union of diameter $<2\varepsilon_i$.

(b) Let $p\in D^{i+1}_j$ and $k=K_i(H_{i+1}(j))$. By (B), $p\in D^i_k$, so $E^i_{H_i(k)}\subseteq\st_i(p)$. By (B$'$) applied to the link $E^{i+1}_{H_{i+1}(j)}$,
$E^{i+1}_{H_{i+1}(j)}\subseteq E^i_{H_i(K_i(H_{i+1}(j)))}=E^i_{H_i(k)}$.
\end{proof}

\formal{\lean{HPC.BJ59.mem_starImage}, \lean{HPC.BJ59.starImage_succ_subset}, \lean{HPC.BJ59.starImage_dist_le}.}

\begin{theorem}[Bing and Jones, pp.~188--189]\label{thm:zigzag}
Under the hypotheses above there is a homeomorphism $h\colon P\to Q$ characterized by: for $p\in P$ and $q\in Q$, $q=h(p)$ if and only if $q\in\cl{\st_i(p)}$ for every $i$.
\end{theorem}

\begin{proof}
\emph{Definition.} By Lemma~\ref{lem:star-props} the sets $\cl{\st_i(p)}\cap Q$ are nonempty, compact and decreasing, with diameters $\le2\varepsilon_i\to0$. Their intersection is
a single point $h(p)$, and the characterization holds by construction.

\emph{Continuity.} Fix $p$ and $i$, and let $W$ be the intersection of the links of $D^i$ containing $p$, an open neighbourhood of $p$. For $p'\in W\cap P$ the sets $\st_i(p)$ and
$\st_i(p')$ have a link in common, so $h(p)$ and $h(p')$ lie in the closure of a set of diameter $<4\varepsilon_i$.

\emph{Onto.} Let $q\in Q$ and $i\ge0$, and let $E^i_k\ni q$. As $H_i$ is onto, $k=H_i(j)$ for some $j$, and $D^i_j$ contains a point $p\in P$. Then $q$ and $h(p)$ both lie in
$\cl{\st_i(p)}$, so $\dist(q,h(p))\le2\varepsilon_i$. Hence $h(P)$ is dense in $Q$; it is compact, so $h(P)=Q$.

\emph{One-to-one.} Let $h(p)=h(p')=y$ and $i\ge1$. Since $y\in\cl{\st_i(p)}\cap\cl{\st_i(p')}$ and the closure of a finite union is the union of the closures, $y\in\cl{E^i_a}\cap
\cl{E^i_b}$ with $a=H_i(j)$, $p\in D^i_j$, and $b=H_i(j')$, $p'\in D^i_{j'}$. Since $E^i$ is a BJ-chain, $\abs{a-b}\le1$ (Lemma~\ref{lem:adj-closure}), so $\abs{K_{i-1}(a)-K_{i-1}(b)}\le1$.
By (B), $p\in D^{i-1}_{K_{i-1}(a)}$ and $p'\in D^{i-1}_{K_{i-1}(b)}$. These two links are equal or adjacent, so $\dist(p,p')<2\varepsilon_{i-1}$. As $i$ is arbitrary, $p=p'$.

A continuous bijection of a compact space onto a Hausdorff space is a homeomorphism.
\end{proof}

\formal{\lean{HPC.BJ59.zigzag_homeomorph}. Only the chains $E^i$ need to be BJ-chains.}

\begin{remark}
Bing and Jones's injectivity argument passes through the sets $K_{i-1}\st_i(p)$, whose closures contain $p$ and $q$, and concludes $\rho(p,q)<4/2^i$; with their meshes the bound
should be in terms of $\varepsilon_{i-1}$. In the sentence ``$e_j^{i+1}\subseteq H_iK_i(e_j^{i+1})=H_i(d^{i+1}_k)\subseteq H_id(p,i)$'' of p.~188 the map $H_i$ is applied to a link
of $D^{i+1}$; what is meant is $H_i(K_iH_{i+1}(d^{i+1}_k))$, as in the proof of Lemma~\ref{lem:star-props}(b).
\end{remark}

\section{Fibres}

\begin{proposition}\label{prop:zigzag-fibre}
In the situation of Theorem~\ref{thm:zigzag}, let $\varphi\colon X\to T$ and $\psi\colon Y\to T$ be maps to a set $T$ such that $Q\cap\psi^{-1}(t)$ is closed for every $t$.
Suppose that whenever $p\in P$ lies in $D^i_j$, the link $E^i_{H_i(j)}$ contains a point $q\in Q$ with $\psi(q)=\varphi(p)$. Then $\psi(h(p))=\varphi(p)$ for all $p\in P$.
\end{proposition}

\begin{proof}
For each $i$, $h(p)$ lies in the closure of one of the finitely many links making up $\st_i(p)$, say $E^i_{H_i(j)}$ with $p\in D^i_j$. That link contains $q_i\in Q$ with
$\psi(q_i)=\varphi(p)$, and $\dist(h(p),q_i)\le\varepsilon_i$. So $h(p)$ is a limit of points of the closed set $Q\cap\psi^{-1}(\varphi(p))$.
\end{proof}

\formal{\lean{HPC.BJ59.zigzag_fibre}.}

In Theorem~10, $P=X$, $Q=Y$, $\varphi=f$, $\psi=g$: the hypothesis of Proposition~\ref{prop:zigzag-fibre} is the first half of Condition~A of Chapter~\ref{ch:thm10}, and the
conclusion says $g\circ h=f$; as $h$ is onto, $h(p_x)=q_x$ for every $x$.

\section*{Exercises}

\begin{exercise}
Show that without Condition (B$'$) the sets $\st_i(p)$ need not decrease, and that $h$ may then fail to be well defined.
\end{exercise}

\begin{exercise}
Show that if every $H_i$ and $K_i$ is a bijection preserving the order of the links, then $h$ maps each link-trace $D^i_j\cap P$ into $\cl{E^i_{H_i(j)}}$.
\end{exercise}

\begin{exercise}
Where in the proof of Theorem~\ref{thm:zigzag} is the surjectivity of $H_i$ used? Construct an example with $H_i$ not onto in which $h$ is not onto.
\end{exercise}

%======================================================================
\chapter{Theorem 10}\label{ch:thm10}
%======================================================================

\begin{theorem}[Bing and Jones, Theorem~10]\label{thm:bj10}
Let $X$, $Y$ be snake-like compact metric spaces and $f\colon X\to\I$, $g\colon Y\to\I$ continuous arcs of pseudo-arcs. Let $h_0\colon\fib f0\to\fib g0$ and
$h_1\colon\fib f1\to\fib g1$ be homeomorphisms. There is a homeomorphism $H\colon X\to Y$ with $g\circ H=f$ which extends $h_0$ and $h_1$.
\end{theorem}

\formal{\lean{HPC.BJ59.bjThm10}, a proof of the statement \lean{HPC.BJ59.BJThm10}.}

The homeomorphism is produced by Theorem~\ref{thm:zigzag}, from chains built by one construction (the \emph{step}) made alternately on the two sides. We write
$p_x=\fib fx$, $q_x=\fib gx$, $P(J)=\fib fJ$ and $Q(J)=\fib gJ$. A family of links \emph{properly covers} a set if each of its links meets the set and their union
contains it.

\section{The conditions carried along}

\begin{definition}[Bing and Jones, p.~188]\label{def:conditions}
Let $D$ be a chain on $X$, $E$ a chain on $Y$, and $H$ a map of the indices of $D$ to those of $E$.
\begin{description}
\item[\textnormal{(A)}] $H$ is a chain map of $D$ onto $E$; if $D_j$ meets $p_x$ then $E_{H(j)}$ meets $q_x$; and if a run of $D$ properly covers $p_x$, the images
of its links cover $q_x$.
\item[\textnormal{(C)}] If $x\in p_0\cap D_j$ then $h_0(x)\in E_{H(j)}$; if $x\in p_1\cap D_j$ then $h_1(x)\in E_{H(j)}$.
\end{description}
\end{definition}

\formal{\lean{HPC.BJ59.IsFibreChainMap} (Condition A), \lean{HPC.BJ59.EndCompat} (Condition C).}

Recall private points (Definition~\ref{def:private}) and fibre-fineness (Definition~\ref{def:fibre-fine}). We also need a name for the conclusion of the
corollary to Theorem~8: a map $R$ of the indices of a chain $F$ \emph{runs through $[\sigma,\sigma']$ on the run $[t,u]$} if every $j\in[\sigma,\sigma']$ is the value of $R$ at
three consecutive indices of $[t,u]$.

\formal{\lean{HPC.BJ59.IsPatternOn}.}

\section{The step}

\begin{theorem}[The step]\label{thm:step}
Let $X$, $Y$, $f$, $g$, $h_0$, $h_1$ be as in Theorem~\ref{thm:bj10}. Let $D$ be a chain covering $X$ with private points, $E$ a chain covering $Y$ which is
fibre-fine for $g$, and $H$ satisfying (A) and (C). For every $\delta>0$ there are a BJ-chain $E'$ covering $Y$, of mesh $<\delta$, with private points, and a map
$K$ of its indices to those of $D$ such that:
\begin{enumerate}[label=(\roman*)]
\item (A$'$) $K$ satisfies (A) with the roles of $X$ and $Y$ exchanged;
\item (C$'$) $K$ satisfies (C) for $h_0^{-1}$, $h_1^{-1}$;
\item (B$'$) $E'_k\subseteq E_{H(K(k))}$ for every $k$.
\end{enumerate}
\end{theorem}

\formal{The statement is \lean{HPC.BJ59.BJStep}, proved as \lean{HPC.BJ59.bj_step}.}

Applied with $X$ and $Y$ exchanged (to $E'$, $D$, $K$, $h_0^{-1}$, $h_1^{-1}$), the step returns a chain $D'$ on $X$ and a chain map $H'$ of $D'$ onto $E'$
with (A), (C) and $D'_j\subseteq D_{K(H'(j))}$, which is Condition~(B) of Chapter~\ref{ch:zigzag}. So one statement serves both halves of the induction.

The construction follows Bing and Jones's Steps~2 and~3 (pp.~183--188), with two differences: the homeomorphisms $h_0$, $h_1$ are used only on the end fibres,
and at an interior cut the construction passes from the map of one sub-arc to the map of the next through an overlap of the sub-arcs. We first assemble the
pieces.

\begin{lemma}[The map at an end; Step~2(b)]\label{lem:end-piece}
Let $c\in\set{0,1}$ and let $h\colon p_c\to q_c$ be a homeomorphism. Let $W=[w_1,w_2]$ be a run of $D$ covering $p_c$ from which neither end link can be dropped. There
is $\varepsilon>0$ such that for every $\varepsilon$-chain $F$ covering $Y$ there is a chain map $R$ of $F$ into $W$ with
\begin{itemize}
\item $h^{-1}(\cl{F_a}\cap q_c)\subseteq D_{R(a)}$ for every link $F_a$;
\item every $j\in W$ is the value of $R$ at three consecutive links meeting $q_c$.
\end{itemize}
\end{lemma}

\begin{proof}
By Lemma~\ref{lem:runs}(c) every $j\in W$ has a point $\varphi(z)$, $z\in q_c$, $\varphi=h^{-1}$, lying in $D_j$ and in no other link of $W$. Let $\lambda$ be a Lebesgue number of the
cover of $p_c$ by the links of $W$, and choose $\varepsilon$ below a third of $\diam q_c$ and such that $\varphi$ maps subsets of $q_c$ of diameter $<5\varepsilon$ to sets of
diameter $<\lambda$. Let $T$ be the run of links of $F$ meeting $q_c$. For $a\in T$ let $S_a$ be the set of $j\in W$ with $\varphi(\cl{F_a}\cap q_c)\subseteq D_j$: a run of one
or two indices, nonempty. Call $a$ \emph{decided} if $S_a=\set j$. If $a$ is decided with value $j$, then $j\in S_{a'}$ for every $a'\in T$ within distance $4$ of $a$ (the
closures of the links from $a$ to $a'$ meet $q_c$ in a set of diameter $<5\varepsilon$, whose image lies in one link of $W$, necessarily $D_j$); so decided positions at most four
apart have the same value. Put $R(a)=j$ if some decided position within distance $2$ of $a$ has value $j$, and $R(a)=\min S_a$ otherwise; extend $R$ outside $T$ by its values at the
ends of $T$. Then $R(a)\in S_a$, and $R$ is a chain map. For $j\in W$, take $z$ as above and $a\in T$ with $z\in F_a$: $a$ is decided with value $j$, so $R=j$ at the
positions of $T$ within distance $2$ of $a$, at least three of them.
\end{proof}

\formal{\lean{HPC.BJ59.exists_end_piece}, from \lean{HPC.BJ59.exists_pattern}.}

\begin{lemma}[The map of a sub-arc]\label{lem:subarc-piece}
Let $a<b$ and suppose every run of $D$ covering $P[a,b]$ minimally has all its links meeting every fibre over $[a,b]$ (Theorem~\ref{thm:step1b}). Then there is a run
$[r,s]$ of $D$ properly covering every $p_x$, $x\in[a,b]$, and $\varepsilon>0$ such that for every $\varepsilon$-chain $F$ covering $Y$ there is a chain map $R$ of $F$ into $[r,s]$
with
\begin{description}
\item[\textnormal{(R1)}] $F_k\cap Q[a,b]\subseteq E_{H(R(k))}$ for every $k$;
\item[\textnormal{(R2)}] for $x\in[a,b]$, $R$ runs through $[r,s]$ on every run of $F$ covering $q_x$.
\end{description}
\end{lemma}

\begin{proof}
Take a minimal covering run $[r,s]$ of $P[a,b]$. By (A), its image $H([r,s])$ (a run, Lemma~\ref{lem:chainmap-ivt}) consists of links meeting every $q_x$, $x\in[a,b]$, and covers
$Q[a,b]$. By fibre-fineness it has one link or at least three. Apply the corollary to Theorem~8 on the sub-arc $Q[a,b]$ (Corollary~\ref{cor:bj8} and the remark after
it) with this run as $E$ and $H$ restricted to $[r,s]$ as $H_1$.
\end{proof}

\formal{\lean{HPC.BJ59.exists_subarc_piece}.}

\begin{lemma}[Stability near a fibre]\label{lem:tau}
If a run of a chain properly covers $\fib\varphi c$ for a continuous open $\varphi$ onto $\I$, it properly covers $\fib\varphi x$ for all $x$ near $c$. If the closure of
every link $F_k$ meets $\fib\varphi c$ inside an open set $V_k$, there is $\tau>0$ with $F_k\cap\fib\varphi{[c-\tau,c+\tau]}\subseteq V_k$ for all $k$.
\end{lemma}

\formal{\lean{HPC.BJ59.exists_tau_run}, \lean{HPC.BJ59.exists_tau_incl}.}

\begin{proof}[Proof of Theorem~\ref{thm:step} in outline]
\emph{(a) Sub-arcs.} Let $\delta_D$ be given by Theorem~\ref{thm:step1b} for $f$ and $D$. Choose $N_0\ge1$ and $\eta>0$ with $1/N_0+2\eta\le\delta_D$ and $\eta<1/(4N_0)$. Put
$c_i=i/N_0$ and $J_i=[c_i-\eta,c_{i+1}+\eta]\cap\I$ for $i<N_0$; consecutive $J_i$ overlap around $c_{i+1}$.

\emph{(b), (c) The pieces.} Lemma~\ref{lem:subarc-piece} on each $J_i$ gives runs $[r_i,s_i]$ and numbers $\varepsilon_i$; Lemma~\ref{lem:end-piece} at $c=0$, $1$ (with $h=h_0$,
$h_1$) gives runs $W^0$, $W^1$ and numbers $\varepsilon^0$, $\varepsilon^1$.

\emph{(d) The chain $F$.} Let $F$ be an $\varepsilon$-chain covering $Y$, with $\varepsilon$ below all $\varepsilon_i$, $\varepsilon^c$ and so small that $g$ maps each link of $F$
to a set of diameter $<\eta/4$. We obtain chain maps $R_i\colon F\to[r_i,s_i]$ with (R1), (R2) on $J_i$, and $R^c\colon F\to W^c$ as in Lemma~\ref{lem:end-piece}.

\emph{(e) Windows.} Let $T^0$, $T^1$ be the runs of $F$ meeting $q_0$, $q_1$, and $T_i$ ($0<i<N_0$) the run meeting $q_{c_i}$. By Lemma~\ref{lem:tau} choose $\tau\in(0,\eta/4)$ such
that these runs properly cover the fibres within $\tau$ of $0$, $1$, $c_i$; $W^c$ properly covers $p_x$ for $\abs{x-c}\le\tau$; and $F_a\cap Q(\set{\abs{x-c}\le\tau})\subseteq
E_{H(R^c(a))}$ for $F_a\in T^c$ (at $x=c$ this follows from Lemma~\ref{lem:end-piece} and (C)). The windows are $L^0=[\tau/4,\tau/2]$, $L^1=[1-\tau/2,1-\tau/4]$ and
$L_i=[c_i,c_i+\tau/2]$. At each window two maps meet: $R^0$ and $R_0$ at $L^0$, $R_{i-1}$ and $R_i$ at $L_i$, $R_{N_0-1}$ and $R^1$ at $L^1$. Both carry the run $T$ of
that window onto runs of $D$ properly covering the fibre at the centre, so Theorem~\ref{thm:bj6} gives a \emph{transition link} $F_y\in T$ at which their values are adjacent.
Lemma~\ref{lem:window-set} gives a transition set $U\subseteq F_y$ over the interior of the window, and Theorem~\ref{thm:bj5-window} splits $Y\setminus U$ into closed sides
$A$ (left) and $B$ (right) such that, for each $x$ in the window, every link of $T$ meets $A\cap q_x$ or every link meets $B\cap q_x$. Intersecting consecutive sides gives
closed, pairwise disjoint regions
\[
Z^0=A^0,\qquad Z_i=B_i\cap A_{i+1}\ (i<N_0,\ B_0:=B^0,\ A_{N_0}:=A^1),\qquad Z^1=B^1,
\]
covering the complement of the transition sets, with $Z_i\subseteq Q(J_i)$; region $Z_i$ carries the map $R_i$, and $Z^c$ the map $R^c$.

\emph{(f) $E'$ and $K$.} Choose $\mu>0$ below $\delta$, a Lebesgue number of $F$, the distances between distinct regions, between each transition set and the regions not next
to it, between the two sides of each window and from each transition set to the complement of its link $F_y$, and so small that a set of diameter $<\mu$ has image under $g$
of diameter $<\tau/4$. Then a set of diameter $<\mu$ meets at most one transition set, and none if it meets $q_0\cup q_1$. Let $G$ be a BJ-chain of mesh $<\mu$ covering $Y$
(Theorem~\ref{thm:shrink}). For a link $e$ of $G$: if $e$ meets a transition set $U\subseteq F_y$, then $e\subseteq F_y$, and $K(e)$ is the left map at $y$ if $e$ meets the side $A$,
otherwise the right map at $y$; if not, $e$ meets exactly one region, and $K(e)$ is that region's map at a link of $F$ containing $e$.

\emph{(g) Verification.}
\begin{enumerate}
\item $K$ is a chain map: inside a region adjacent links of $G$ lie in equal or adjacent links of $F$; a link meeting $U$ and an adjacent link not meeting it meet on the same side
and use the same map; two links meeting $U$ lie in $F_y$, where the two maps are adjacent.
\item (B$'$) follows from (R1) in the regions $Z_i$ and at interior transitions (where $e\subseteq F_y\cap Q(J_{i-1}\cap J_i)$), and from the choice of $\tau$ in $Z^c$.
\item (C$'$): a link meeting $q_c$ meets no transition set, so it lies in $Z^c$, and Lemma~\ref{lem:end-piece} applies.
\item (A$'$), first clause: if $e$ meets $q_x$, $x$ lies in the range of the map used, whose values all name links meeting $p_x$.
\item (A$'$), second clause. Let a run $C$ of $G$ properly cover $q_x$. We find a run $S$ of $D$ properly covering $p_x$, a map $M$ among the $R_i$, $R^c$, a run $T$ of $F$
and a set $\Gamma\subseteq q_x$ meeting every link of $T$, such that every $j\in S$ is the value of $M$ at three consecutive links $F_{a-1},F_a,F_{a+1}$ of $T$ and every
link of $G$ meeting $\Gamma$ takes its value from $M$. Then for $j\in S$, a point $z\in\Gamma\cap F_a$ lies in some $e\in C$, the link of $F$ used for $e$ contains $z$, so it is
$F_{a-1}$, $F_a$ or $F_{a+1}$, and $K(e)=j$. Between windows, $q_x$ lies in one region $Z_i$: take $M=R_i$, $\Gamma=q_x$, $S=[r_i,s_i]$ and (R2). In a window, the dichotomy of
Theorem~\ref{thm:bj5-window} says that every link of $T$ meets $A\cap q_x$ (take the left map, $\Gamma=A\cap q_x$) or every link meets $B\cap q_x$ (the right map); (R2) at the
centre of the window, or Lemma~\ref{lem:end-piece} at the ends, gives the three consecutive links. Near $0$ and $1$ use $R^c$ and $W^c$.
\item $K$ is onto: for a private point $z$ of $D_j$ and $x=f(z)$, the run $S$ of item~5 covers $z$, so $j\in S$ is a value of $K$.
\item Private points: shrink $G$ by Proposition~\ref{prop:private-shrink} and keep $K$. Every property survives, since links only shrink and a run of shrunk links properly
covering $q_x$ comes from a run of $G$ properly covering it.\qedhere
\end{enumerate}
\end{proof}

\formal{The gluing (e)--(g) is \lean{HPC.BJ59.bj_glue}; the pieces are \lean{HPC.BJ59.exists_subarc_piece}, \lean{HPC.BJ59.exists_end_piece}; the windows use
\lean{HPC.BJ59.adjacent_of_properCover}, \lean{HPC.BJ59.exists_window_set}, \lean{HPC.BJ59.bj_theorem5_window}.}

\section{The induction}

\begin{proof}[Proof of Theorem~\ref{thm:bj10}]
By Lemma~\ref{lem:fibre-spread} let $\kappa_X$, $\kappa_Y$ bound the spreads of the fibres of $f$ and $g$ from below. A \emph{stage} is a triple $(D,E,H)$ of chains covering
$X$ and $Y$, BJ-chains with private points, $D$ fibre-fine for $f$ and $E$ for $g$, and $H$ satisfying (A) and (C), with both meshes below a given $\varepsilon$.

\emph{Stage $0$}: the one-link chains $(X)$, $(Y)$ and the identity map. All conditions hold trivially (a one-link chain is fibre-fine, its link has private points, (A) and
(C) hold because the only link of $E$ is $Y$).

\emph{From stage $i$ to stage $i+1$}: apply Theorem~\ref{thm:step} to $(D^i,E^i,H_i)$ with $\delta_i=\min(2^{-i-1},\kappa_X/3,\kappa_Y/3)$, obtaining $E^{i+1}$ and $K_i$;
then apply it with $X$ and $Y$ exchanged to $(E^{i+1},D^i,K_i)$, obtaining $D^{i+1}$ and $H_{i+1}$. By Lemma~\ref{lem:fibre-fine} the new chains are fibre-fine. Condition
(B$'$) of Chapter~\ref{ch:zigzag} is (iii) of the first application and (B) is (iii) of the second.

\emph{The homeomorphism.} Theorem~\ref{thm:zigzag} with $P=X$, $Q=Y$ gives a homeomorphism $H\colon X\to Y$ with $H(p)\in\cl{\st_i(p)}$ for every $i$.
Proposition~\ref{prop:zigzag-fibre}, whose hypothesis is the first clause of (A), gives $g\circ H=f$. For $x\in p_0$, Condition (C) gives $h_0(x)\in\st_i(x)$ for every $i$, so
$h_0(x)=H(x)$ by the characterization; likewise for $h_1$.
\end{proof}

\formal{\lean{HPC.BJ59.bjThm10_of_step}, with \lean{HPC.BJ59.BJStage}, \lean{HPC.BJ59.BJStage.exists_init}, \lean{HPC.BJ59.BJStage.exists_next}.}

\begin{remark}
Bing and Jones extend $h$ at every step to the fibres over the cut points $i/n$ (arbitrarily in Step~1(c), by their Theorem~9 in Step~2(i)) and build the maps near those
fibres from it. In the proof above $h_0$ and $h_1$ are used only on $p_0$ and $p_1$, and Theorem~9 is not needed. There is also no separate first step: the induction starts
from one-link chains, and for a one-link chain the corollary to Theorem~8 is the zigzag map of their Step~2(b).
\end{remark}

\formal{Bing and Jones's Theorem~9 is nevertheless proved: \lean{HPC.BJ59.bj_theorem9}.}

\section*{Exercises}

\begin{exercise}
Check that the one-link stage satisfies (A): what does ``a run of $D$ properly covering $p_x$'' mean when $D$ has one link?
\end{exercise}

\begin{exercise}
In Lemma~\ref{lem:end-piece}, show that two runs $S_a$, $S_{a+1}$ of length at most two that meet have least elements at most one apart, and deduce that $R$ is a chain map.
\end{exercise}

\begin{exercise}
Show that Theorem~\ref{thm:bj10} gives, for any homeomorphism $h_0\colon p_0\to p_0$ of the end fibre of a continuous snake-like arc of pseudo-arcs, a self-homeomorphism of $X$
preserving every fibre and extending $h_0$.
\end{exercise}

%======================================================================
\chapter{Uniqueness, homogeneity and recognition}\label{ch:uniqueness}
%======================================================================

``It follows from Theorem~10 that any two circles of pseudo-arcs are topologically equivalent and that each is homogeneous'' (Bing and Jones, p.~189). This
chapter supplies the argument: cut each circle of pseudo-arcs into its two halves, apply Theorem~10 to each half, and glue. The same argument, with
Corollary~\ref{cor:atomic-halves} in place of Proposition~\ref{prop:halves-snake}, recognizes the continua of Volume~7 as circles of pseudo-arcs.

Throughout, $X$ and $Y$ are compact metric spaces and $f\colon X\to\Sph$, $g\colon Y\to\Sph$ have snake-like halves (Definition~\ref{def:snake-halves}). Recall that
the end fibres of the half maps $f_s$ are $\fib f1$ and $\fib f{-1}$, for both $s$.

\section{Gluing two halves}

\begin{lemma}[One half]\label{lem:one-half}
Let $k_+\colon\fib f1\to\fib g1$ and $k_-\colon\fib f{-1}\to\fib g{-1}$ be homeomorphisms and $s\in\set{1,-1}$. There is a homeomorphism $H_s\colon X_s\to Y_s$ with $g\circ
H_s=f$ on $X_s$, equal to $k_+$ on $\fib f1$ and to $k_-$ on $\fib f{-1}$.
\end{lemma}

\begin{proof}
The half maps $f_s\colon X_s\to\I$ and $g_s\colon Y_s\to\I$ are continuous arcs of pseudo-arcs (Proposition~\ref{prop:halves}) on snake-like compact spaces. Their fibres over
$0$ are $\fib f1$, $\fib g1$ and over $1$ are $\fib f{-1}$, $\fib g{-1}$. Theorem~\ref{thm:bj10}, with $h_0=k_+$ and $h_1=k_-$ (transported through these identifications),
gives a homeomorphism $H_s$ with $g_s\circ H_s=f_s$, that is, $g\circ H_s=f$.
\end{proof}

\formal{\lean{HPC.BJ59.exists_half_homeomorph'}, with \lean{HPC.BJ59.halfFibreEquiv}.}

\begin{theorem}[Gluing]\label{thm:glue}
With $k_+$, $k_-$ as in Lemma~\ref{lem:one-half} there is a homeomorphism $\Phi\colon X\to Y$ with $g\circ\Phi=f$ and $\Phi=k_+$ on $\fib f1$.
\end{theorem}

\begin{proof}
Let $H_1$, $H_{-1}$ be given by Lemma~\ref{lem:one-half}. The halves $X_1$, $X_{-1}$ are closed, cover $X$, and meet in $\fib f1\cup\fib f{-1}$, where $H_1$ and $H_{-1}$
both equal $k_+$ or $k_-$. So $\Phi=H_1$ on $X_1$, $\Phi=H_{-1}$ on $X_{-1}$ is well defined, and continuous by the pasting lemma for two closed sets. It satisfies $g\circ\Phi=f$.
It is injective: if $\Phi(x)=\Phi(x')$ then $f(x)=f(x')$, so $x$, $x'$ lie in a common half, where $\Phi$ is injective. It is onto, since $Y=Y_1\cup Y_{-1}$ and $H_s$ maps
$X_s$ onto $Y_s$. A continuous bijection of compact metric spaces is a homeomorphism.
\end{proof}

\formal{\lean{HPC.BJ59.exists_homeomorph_glue'}; Theorem~\ref{thm:bj10} enters as the hypothesis \lean{HPC.BJ59.BJThm10}, discharged by \lean{HPC.BJ59.bjThm10}.}

\begin{corollary}\label{cor:snakehalves-homeo}
Any two compact metric spaces carrying maps with snake-like halves are homeomorphic.
\end{corollary}

\begin{proof}
The fibres $\fib f1$, $\fib g1$ are pseudo-arcs, hence homeomorphic (\vtwo{Theorem~8.3}); likewise $\fib f{-1}$, $\fib g{-1}$. Apply Theorem~\ref{thm:glue}.
\end{proof}

\formal{\lean{HPC.BJ59.homeomorph_of_isSnakeHalvesMap}, with \lean{HPC.Plane.theorem1_HIC'}.}

\section{Circles of pseudo-arcs}

\begin{theorem}[Bing and Jones, p.~189]\label{thm:unique}
Any two circles of pseudo-arcs are homeomorphic.
\end{theorem}

\begin{proof}
A circle of pseudo-arcs is compact (it is circle-like) and its map has snake-like halves (Proposition~\ref{prop:halves-snake}). Apply Corollary~\ref{cor:snakehalves-homeo}.
\end{proof}

\formal{\lean{HPC.BJ59.homeomorph_of_isCircleOfPA'}; with Theorem~10 as a hypothesis, \lean{HPC.BJ59.homeomorph_of_isCircleOfPA} and
\lean{HPC.BJ59.homeomorph_of_isCircPAMap}.}

\begin{theorem}[Bing and Jones, p.~189]\label{thm:homog}
Every circle of pseudo-arcs is homogeneous.
\end{theorem}

\begin{proof}
Let $f\colon X\to\Sph$ be a circle-of-pseudo-arcs map and $x,y\in X$. The rotated maps $f'=f(x)^{-1}f$ and $f''=f(y)^{-1}f$ are again circle-of-pseudo-arcs maps, with $f'(x)=1$
and $f''(y)=1$. Their fibres over $1$ are the fibres $P_x\ni x$ and $P_y\ni y$ of $f$. These are pseudo-arcs; a homeomorphism $P_x\to P_y$ followed by a homeomorphism of $P_y$
moving its image of $x$ to $y$ (\vtwo{Theorems~8.3 and~8.4}) gives $k_+\colon P_x\to P_y$ with $k_+(x)=y$. Let $k_-$ be any homeomorphism between the fibres of $f'$ and $f''$
over $-1$. Theorem~\ref{thm:glue} for $f'$ and $f''$ gives a homeomorphism $\Phi$ of $X$ with $\Phi(x)=k_+(x)=y$.
\end{proof}

\formal{\lean{HPC.BJ59.isHomogeneous_of_isCircleOfPA'}; with Theorem~10 as a hypothesis, \lean{HPC.BJ59.isHomogeneous_of_isCircleOfPA},
\lean{HPC.BJ59.isHomogeneous_of_isCircPAMap}; the pseudo-arc facts are \lean{HPC.Plane.theorem1_HIC'} and \lean{HPC.Plane.homogeneous_of_isPseudoArc}.}

\section{Recognition}

\begin{lemma}\label{lem:circlelike-homeo}
The circle is circle-like, and a compact metric space homeomorphic to a circle-like space is circle-like.
\end{lemma}

\begin{proof}
For the circle, take the $N$ open arcs of length $3\pi/N$ centred at the $N$-th roots of unity, $N\ge4$ large. For the second statement pull the links of fine circular chains back
along the homeomorphism $e\colon X\to Y$; by uniform continuity of $e$ on the compact space $X$, the pulled-back chains are fine.
\end{proof}

\formal{\lean{HPC.BJ59.isCircleLike_circle}, \lean{HPC.BJ59.isCircleLike_univ_of_homeomorph}.}

\begin{theorem}[Domokos and Prajs, Theorem~6.1]\label{thm:dp61}
Let $X$ be a compact metric space with an atomic circle-of-pseudo-arcs map $f\colon X\to\Sph$ (continuous, open, surjective, atomic, with pseudo-arc fibres). Then $X$ is
homeomorphic to every circle of pseudo-arcs. If a circle of pseudo-arcs exists, $X$ is one. Conversely, every circle of pseudo-arcs carries an atomic circle-of-pseudo-arcs
map.
\end{theorem}

\begin{proof}
By Corollary~\ref{cor:atomic-halves}, $f$ has snake-like halves; so does a circle-of-pseudo-arcs map on any circle of pseudo-arcs $Y$. By Corollary~\ref{cor:snakehalves-homeo},
$X$ is homeomorphic to $Y$. Then $X$ is circle-like (Lemma~\ref{lem:circlelike-homeo}), and with $f$ it is a circle of pseudo-arcs.

For the converse, let $C$ be the plane circle of pseudo-arcs of Theorem~\ref{thm:plane-example}. It is homogeneous (Theorem~\ref{thm:homog}) and decomposable
(Corollary~\ref{cor:circle-pa-decomp}), so by \vseven{Theorem~9.8} it is a simple closed curve or carries an atomic circle-of-pseudo-arcs map. It is not a simple closed curve
(Corollary~\ref{cor:circle-not-cpa}). An atomic circle-of-pseudo-arcs map on $C$ transports along a homeomorphism $X\to C$ (Theorem~\ref{thm:unique}) to one on $X$: fibres,
openness and atomicity are preserved.
\end{proof}

\formal{\lean{HPC.BJ59.dp22_theorem61'} and \lean{HPC.BJ59.dp22_theorem61_iff'}; with the inputs as hypotheses, \lean{HPC.BJ59.dp22_theorem61},
\lean{HPC.BJ59.isCircleOfPA_of_atomic}, \lean{HPC.BJ59.dp22_theorem61_iff}; the converse uses \lean{HPC.BJ59.exists_atomic_of_plane_isCircleOfPA} and
\lean{HPC.BJ59.IsAtomicCircPAMap.comp_homeomorph}.}

Two compact spaces with atomic circle-of-pseudo-arcs maps are homeomorphic, and each is homogeneous, by the same arguments.

\formal{\lean{HPC.BJ59.homeomorph_of_atomic}, \lean{HPC.BJ59.isHomogeneous_of_atomic}.}

\section{Circle-likeness of the decomposable continua}

Volume~7 left one gap: it did not show that a decomposable homogeneous plane continuum is circle-like.

\begin{theorem}\label{thm:ho-circlelike}
Every decomposable homogeneous plane continuum is circle-like.
\end{theorem}

\begin{proof}
Let $M$ be such a continuum. By \vseven{Theorem~9.8}, $M$ is a simple closed curve, which is circle-like by Lemma~\ref{lem:circlelike-homeo}, or $M$ carries an atomic
circle-of-pseudo-arcs map. In the second case, Theorem~\ref{thm:dp61} with the plane circle of pseudo-arcs of Theorem~\ref{thm:plane-example} shows that $M$ is a circle of
pseudo-arcs, in particular circle-like.
\end{proof}

\formal{\lean{HPC.BJ59.hoCircleLike'}, a proof of the statement \lean{HPC.BJ59.HOCircleLike}; with Theorem~10 and the plane example as hypotheses,
\lean{HPC.BJ59.hoCircleLike}.}

\section*{Exercises}

\begin{exercise}
Show that in Theorem~\ref{thm:glue} the homeomorphism can be chosen to equal any prescribed homeomorphism $k_-$ on $\fib f{-1}$ as well.
\end{exercise}

\begin{exercise}
Show that the homeomorphism $\Phi$ constructed in the proof of Theorem~\ref{thm:homog} satisfies $f\circ\Phi=f(y)f(x)^{-1}\,f$; so it permutes the
fibres of $f$ by a rotation of $\Sph$.
\end{exercise}

\begin{exercise}
Let $f\colon X\to\Sph$ be a circle-of-pseudo-arcs map and $c\in\Sph$. Show that there is a homeomorphism $\Phi$ of $X$ with $f\circ\Phi=c\,f$.
\end{exercise}

%======================================================================
\chapter{Circular descriptions}\label{ch:circdescr}
%======================================================================

The rest of the volume constructs a circle of pseudo-arcs in the plane. Bing and Jones's \S7 lists nine properties of a sequence of circular chains in the
plane, without constructing the chains and without arguing that the resulting continuum has pseudo-arc fibres; as far as one can see, their condition~(9)
makes each new column crooked in a pair of old columns, which does not make a run of two new columns crooked in two old ones. The construction here is
different. It is a circular analogue of the descriptions of \vone{Chapter~3}: this chapter gives the abstract conditions and proves that they produce a circle
of pseudo-arcs; Chapter~\ref{ch:sweep} gives explicit patterns, and Chapter~\ref{ch:plane} realizes them in the plane.

\section{The definition}

We use crookedness of integer sequences on a window, as in \vone{Definition~4.1}: a sequence $g$ is \emph{crooked on $[lo,hi]$} if whenever $lo\le i,j\le hi$,
$g(i)\le h$, $g(j)\ge k$ and $h+3\le k$, there are $r$, $s$ strictly between $i$ and $j$, with $r$ nearer to $i$ than $s$, such that $g(r)=k-1$ and $g(s)=h+1$.

\formal{\lean{HPC.Fold.PatCrookedOn}.}

\begin{definition}\label{def:circdescr}
A \emph{circular description} in a metric space $X$ consists of the following data, for every level $k\ge0$.
\begin{enumerate}[label=(\arabic*)]
\item \emph{Circular chains.} Integers $N_k\ge4$ and open preconnected sets $L_k(t)\subseteq X$, $t\in\Z$, with $L_k(t+N_k)=L_k(t)$, such that $L_k(s)$ meets $L_k(t)$ if and
only if $s-t$ is within $1$ of a multiple of $N_k$. All links lie in a fixed compact set $K$, and $\diam L_k(t)\le\varepsilon_k$ with $\varepsilon_k\to0$.
\item \emph{Lifted patterns.} Maps $\pi_k\colon\Z\to\Z$ of degree one, $\pi_k(t+N_{k+1})=\pi_k(t)+N_k$, with $\abs{\pi_k(t+1)-\pi_k(t)}\le1$, such that
\begin{itemize}
\item (nest) $\cl{L_{k+1}(t)}\subseteq L_k(\pi_k(t))$;
\item (tight) if $L_{k+1}(t)$ meets $L_k(l)$, then $\pi_k(t')=l+nN_k$ for some $t'$ with $\abs{t'-t}\le1$ and some $n\in\Z$.
\end{itemize}
\item \emph{Arcs.} Reals $o_k$ (offsets) and $\lambda_k$ (arc lengths) with $\lambda_k\to0$, $5\le\lambda_kN_k$ and $\lambda_kN_k+4\le N_k$. Link $t$ of level $k$ carries the
\emph{centre} $c_k(t)=o_k+t/N_k\in\R$ and the closed arc of length $\lambda_k$ about it (read in $\R/\Z$; the arcs are labels, not geometry). With $\delta_k=1/(2N_k)$:
\begin{itemize}
\item (arcNest) $\abs{c_{k+1}(t)-c_k(\pi_k(t))}\le(\lambda_k-\lambda_{k+1})/2$ for every $t$;
\item (surv) if $\abs{y-c_k(l)}\le\lambda_k/2-\delta_k$, there is $t$ with $\pi_k(t)=l$ and $\abs{y-c_{k+1}(t)}\le\lambda_{k+1}/2-\delta_{k+1}$;
\item (crook) for every $a\in\Z$ the sequence $u\mapsto\pi_k(a+u)$ is crooked on $[0,\cel{\lambda_{k+1}N_{k+1}}+3]$.
\end{itemize}
\end{enumerate}
The continuum \emph{described} is $M=\bigcap_k\bigcup_tL_k(t)$.
\end{definition}

\formal{\lean{HPC.BJ59.CircDescr}, \lean{HPC.BJ59.CircDescr.M}; the centres are \lean{HPC.BJ59.CircDescr.c}.}

For $y\in\R$, the \emph{run} of $y$ at level $k$ is $\run_k(y)=\set{t\in\Z:\abs{y-c_k(t)}\le\lambda_k/2}$, a set of consecutive integers with at most $\lambda_kN_k+1\le N_k-3$
elements. The \emph{fibre} over $y$ is
\[
\mathrm{fib}(y)=\bigcap_k\ \bigcup_{t\in\run_k(y)}L_k(t).
\]

\formal{\lean{HPC.BJ59.CircDescr.fib}, \lean{HPC.BJ59.CircDescr.run_iff}, \lean{HPC.BJ59.CircDescr.run_len}.}

\section{The recognition theorem}

\begin{theorem}\label{thm:circdescr}
The continuum described by a circular description is a circle of pseudo-arcs.
\end{theorem}

\formal{\lean{HPC.BJ59.CircDescr.isCircleOfPA}.}

The proof occupies the rest of the chapter. Write $U_k=\bigcup_tL_k(t)$.

\begin{lemma}[$M$ is a continuum]\label{lem:cd-continuum}
$\cl{U_{k+1}}\subseteq U_k$, so $M=\bigcap_k\cl{U_k}$ is a nested intersection of nonempty compact connected sets, hence a continuum.
\end{lemma}

\begin{proof}
Up to periodicity a level has finitely many links, so $\cl{U_{k+1}}$ is the union of the closures $\cl{L_{k+1}(t)}\subseteq L_k(\pi_k(t))$. Each $\cl{U_k}$ is compact (it lies in
$K$). $U_k$ is connected: consecutive links meet and are preconnected, so $L_k(0)\cup\dots\cup L_k(t)$ is connected for every $t\ge0$, and $U_k$ is the union of these increasing
sets.
\end{proof}

\formal{\lean{HPC.BJ59.CircDescr.closure_U_succ}, \lean{HPC.BJ59.CircDescr.M_eq}, \lean{HPC.BJ59.CircDescr.isCompact_M}, \lean{HPC.BJ59.CircDescr.isConnected_M}.}

\begin{lemma}[Consecutive links meet in $M$]\label{lem:cd-meet}
For all $k$ and $t$, $L_k(t)\cap L_k(t+1)\cap M\neq\emptyset$.
\end{lemma}

\begin{proof}
Since $\pi_k$ has degree one and steps of at most one, it has an index $s$ with $\pi_k(s)=t$ and $\pi_k(s+1)=t+1$ (take an index with value $\le t$, a later one with value $\ge t+1$, and
the last index between them with value $\le t$). Repeating at every level gives $t_k=t$ and $t_j$ ($j>k$) with $\pi_j(t_{j+1})=t_j$, $\pi_j(t_{j+1}+1)=t_j+1$. The sets
$C_j=\cl{L_j(t_j)}\cap\cl{L_j(t_j+1)}$ are compact, nonempty and nested ($C_{j+1}\subseteq L_j(t_j)\cap L_j(t_j+1)$ by (nest)). A point of their intersection lies in $U_j$ for every
$j\ge k$, and by (nest) also for $j<k$; so it lies in $M\cap L_k(t)\cap L_k(t+1)$.
\end{proof}

\formal{\lean{HPC.BJ59.CircDescr.exists_pat_pair}, \lean{HPC.BJ59.CircDescr.inter_M_nonempty}.}

\begin{lemma}[Fibres]\label{lem:cd-fibre}
\begin{enumerate}[label=(\alph*)]
\item A link of $\run_{k+1}(y)$ lies in the link of $\run_k(y)$ named by $\pi_k$; so the sets $\bigcup_{t\in\run_k(y)}L_k(t)$ decrease, and $\mathrm{fib}(y)\subseteq M$.
\item (Survival) If $\abs{y-c_k(l)}\le\lambda_k/2-\delta_k$, then $\mathrm{fib}(y)$ meets $L_k(l)$.
\item $\mathrm{fib}(y+n)=\mathrm{fib}(y)$ for $n\in\Z$; every point of $M$ lies in some fibre; and a point lying in $\mathrm{fib}(y)$ and $\mathrm{fib}(y')$ has $y-y'\in\Z$.
\end{enumerate}
\end{lemma}

\begin{proof}
(a) is (nest) and (arcNest). (b) Apply (surv) repeatedly: $t_k=l$ and $t_{j+1}$ with $\pi_j(t_{j+1})=t_j$ and $\abs{y-c_j(t_j)}\le\lambda_j/2-\delta_j$. The compact sets
$\cl{L_j(t_j)}$ are nested and nonempty; a point of their intersection lies in $L_j(t_j)$ with $t_j\in\run_j(y)$ for all $j\ge k$, hence (by (a)) in $\mathrm{fib}(y)$, and it lies in
$L_k(l)$. (c) The first statement is periodicity. For $p\in M$ and each $k$, the set of $y$ such that $p$ lies in a link of $\run_k(y)$ is closed and nonempty, and these sets
decrease by (a); a point of their intersection (taken in one period) gives a fibre containing $p$. For the last statement: at level $k$, $p$ lies in at most two consecutive links
(up to periodicity), so $y$ and $y'$ are within $\lambda_k+2/N_k$ of each other modulo $1$; this tends to $0$.
\end{proof}

\formal{\lean{HPC.BJ59.CircDescr.down}, \lean{HPC.BJ59.CircDescr.fib_subset_M}, \lean{HPC.BJ59.CircDescr.fib_meets}, \lean{HPC.BJ59.CircDescr.fib_add_int},
\lean{HPC.BJ59.CircDescr.exists_fib}, \lean{HPC.BJ59.CircDescr.eq_add_int_of_mem}.}

\begin{definition}\label{def:cd-fmap}
The map $\Phi\colon M\to\Sph$ sends $p$ to $\exp(2\pi iy)$ for any $y$ with $p\in\mathrm{fib}(y)$. By Lemma~\ref{lem:cd-fibre}(c) it is well defined, and its fibre over $\exp(2\pi
iy)$ is $\mathrm{fib}(y)$.
\end{definition}

\formal{\lean{HPC.BJ59.CircDescr.fmap}, \lean{HPC.BJ59.CircDescr.fmap_eq_iff}.}

\begin{proposition}\label{prop:cd-map}
$\Phi$ is continuous, open and surjective.
\end{proposition}

\begin{proof}
\emph{Continuous.} If $p_n\to p$ and $p\in L_k(t)$, then $p_n\in L_k(t)$ for large $n$, and the parameters of $p_n$ and $p$ lie (modulo $1$) in the union of the arcs of $t-1,t,t+1$, of
length $\lambda_k+2/N_k$. Let $k\to\infty$.

\emph{Surjective.} Consecutive centres are $1/N_k$ apart and $\lambda_k/2-\delta_k\ge1/(2N_k)$, so every $y$ is within $\lambda_k/2-\delta_k$ of a centre; apply survival.

\emph{Open.} Let $p\in M$, $r>0$, and $k$ with $2\varepsilon_k<r$. Choose $t$ with $p\in L_k(t)$ and a parameter $z$ of $p$ with $\abs{z-c_k(t)}\le\lambda_k/2$. The three shrunk arcs
$\abs{y-c_k(s)}\le\lambda_k/2-\delta_k$, $s=t-1,t,t+1$, cover $\abs{y-c_k(t)}\le\lambda_k/2+\delta_k$, so every $y$ with $\abs{y-z}<\delta_k$ is in one of them. Survival gives $q\in\mathrm{fib}(y)\cap
L_k(s)$; as $L_k(s)$ meets $L_k(t)\ni p$, $\dist(p,q)\le2\varepsilon_k<r$. So the image of $B(p,r)\cap M$ contains the image of the interval $(z-\delta_k,z+\delta_k)$.
\end{proof}

\formal{\lean{HPC.BJ59.CircDescr.continuous_fmap}, \lean{HPC.BJ59.CircDescr.surjective_fmap}, \lean{HPC.BJ59.CircDescr.isOpenMap_fmap}.}

\begin{lemma}[Runs are crooked chains]\label{lem:cd-runs}
For $y\in\R$, the links $L_k(t)$, $t\in\run_k(y)$, in increasing order, form a chain $D^k(y)$, and its traces on $M$ form a chain of $M$. The chain $D^{k+1}(y)$ refines $D^k(y)$ and is crooked
in it (\vone{Definition~2.11}).
\end{lemma}

\begin{proof}
The run has at most $N_k-3$ elements, so by the meeting rule two of its links meet exactly when their indices differ by at most one. Consecutive traces on $M$ meet, and in
particular are nonempty, by Lemma~\ref{lem:cd-meet}.

By (arcNest), $\pi_k$ maps $\run_{k+1}(y)=\set{a,\dots,a+n-1}$ into $\run_k(y)$ with no shift by a multiple of $N_k$, and (nest) gives the refinement. Tightness holds inside the runs: if
$L_{k+1}(t)$ meets $L_k(l)$, then $\pi_k(t')=l+nN_k$ with $\abs{t'-t}\le1$, and $n=0$ because $l$ and $\pi_k(t)$ lie in a run of fewer than $N_k-2$ elements. By (crook) with $a-1$ in place of
$a$, $\pi_k$ is crooked on $[a-1,a+n]$ (as $n+1\le\cel{\lambda_{k+1}N_{k+1}}+2$). Now argue as in \vone{Proposition~4.10}, with the window enlarged by one index at each end: given links $i$, $j$ of $D^{k+1}(y)$ meeting links $\alpha$, $\beta$ of
$D^k(y)$ with $\alpha+2<\beta$, tightness gives $u'$, $v'$ within one of $i$, $j$ with $\pi_k(u')=\alpha$, $\pi_k(v')=\beta$; crookedness gives $r$, $s$ strictly between $u'$ and $v'$, in the
right order, with $\pi_k(r)=\beta-1$ and $\pi_k(s)=\alpha+1$. Since $\pi_k(r)=\beta-1>\alpha+1\ge\pi_k(i)$ we have $r\ne i$; similarly $s\ne j$; so $r$ and $s$ lie strictly between $i$ and
$j$, inside the run.
\end{proof}

\formal{\lean{HPC.BJ59.CircDescr.runChain}, \lean{HPC.BJ59.CircDescr.runChain_crooked}, with \lean{HPC.BJ59.CircDescr.tight_run}, \lean{HPC.BJ59.CircDescr.pat_mem_run}.}

\begin{proposition}\label{prop:cd-fibre-pa}
Every fibre $\mathrm{fib}(y)$ is a pseudo-arc.
\end{proposition}

\begin{proof}
\emph{Continuum:} $\mathrm{fib}(y)=\bigcap_k\cl{W_k}$, with $W_k$ the union of the run at level $k$ (one period); $\cl{W_{k+1}}\subseteq W_k$, and $W_k$ is connected.
\emph{Nondegenerate:} the indices $l$ with $\abs{y-c_k(l)}\le\lambda_k/2-\delta_k$ fill an interval of length $\lambda_kN_k-1\ge4$, so they include some $l$ and $l+2$; by survival the fibre meets
the disjoint links $L_k(l)$ and $L_k(l+2)$. \emph{Snake-like:} the run chains cover the fibre and have mesh $\le\varepsilon_k$. \emph{Hereditarily indecomposable:} Bing's argument
(\vone{Theorem~3.7}) applies verbatim to the run chains, which are crooked in one another by Lemma~\ref{lem:cd-runs}: if a subcontinuum $H=A\cup B$ of the fibre had proper subcontinua
$A$, $B$, take $p\in H\setminus A$, $q\in H\setminus B$ at distance $\ge\delta$ from $A$ and $B$ respectively, and $k$ with $3\varepsilon_k<\delta$; the crookedness of $D^{k+1}(y)$ in $D^k(y)$
gives a contradiction.
\end{proof}

\formal{\lean{HPC.BJ59.CircDescr.isPseudoArc_fib}, from \lean{HPC.BJ59.CircDescr.isContinuum_fib}, \lean{HPC.BJ59.CircDescr.fib_nontrivial},
\lean{HPC.BJ59.CircDescr.isSnakeLike_fib}, \lean{HPC.BJ59.CircDescr.isHereditarilyIndecomposable_fib}.}

\begin{lemma}\label{lem:cd-circlelike}
$M$ is circle-like.
\end{lemma}

\begin{proof}
Given $\varepsilon>0$ take $k$ with $\varepsilon_k<\varepsilon$. The traces $L_k(t)\cap M$, $0\le t<N_k$, are open in $M$ and cover it; two of them meet if and only if their indices are cyclically
adjacent (the meeting rule and Lemma~\ref{lem:cd-meet}), and each is nonempty. So they form a circular chain of $M$ (\vthree{Definition~5.13}) of mesh $<\varepsilon$.
\end{proof}

\formal{\lean{HPC.BJ59.CircDescr.traceChain}, \lean{HPC.BJ59.CircDescr.isCircleLike_M}.}

\begin{proof}[Proof of Theorem~\ref{thm:circdescr}]
$M$ is circle-like (Lemma~\ref{lem:cd-circlelike}), and $\Phi$ is continuous, open and surjective (Proposition~\ref{prop:cd-map}) with pseudo-arc fibres
(Proposition~\ref{prop:cd-fibre-pa}).
\end{proof}

\section*{Exercises}

\begin{exercise}
Show that without (surv) the map $\Phi$ need not be open. (Let all the arcs of level $k+1$ crowd into one half of the arcs of level $k$.)
\end{exercise}

\begin{exercise}
Show that the condition $5\le\lambda_kN_k$ is what makes the fibres nondegenerate, and that $\lambda_kN_k+4\le N_k$ is what makes runs into chains.
\end{exercise}

\begin{exercise}
Show that a point of $M$ lies in at most two links of each level, up to periodicity, and that these two links are consecutive.
\end{exercise}

%======================================================================
\chapter{Sweep patterns}\label{ch:sweep}
%======================================================================

This chapter produces lifted patterns $x\colon\Z\to\Z$ with the properties required of the $\pi_k$ in Definition~\ref{def:circdescr}. The new chain sweeps back
and forth across a window of old links; the window slides once around the circle, one end at a time. Each sweep is a stuttered Bing pattern (\vone{Definitions~4.4
and~4.8}), and a combinatorial lemma, the sweep lemma, makes every run of the new chain crooked.

Since the indices are now integers, we use crookedness of $g\colon\Z\to\Z$ on $[lo,hi]\subseteq\Z$, defined exactly as in Chapter~\ref{ch:circdescr}.

\formal{\lean{HPC.BJ59.PatCrookedOnZ}, with the passage to and from sequences indexed by $\N$, \lean{HPC.BJ59.PatCrookedOnZ.toNat}, \lean{HPC.BJ59.PatCrookedOn.toZ}.}

\section{The sweep lemma}

\begin{definition}\label{def:sweep-walk}
Fix $L>0$. A \emph{sweep walk} is a map $g\colon\Z\to\Z$ together with windows $[b_n,t_n]$ ($t_n>b_n$) and directions (up or down) for the blocks $[nL,(n+1)L]$,
$n\in\Z$ (consecutive blocks share their end points), such that:
\begin{itemize}
\item on block $n$, $g$ takes values in $[b_n,t_n]$ and is crooked; if block $n$ is up it starts at $b_n$ and ends at $t_n$, if down it starts at $t_n$ and ends at $b_n$;
\item directions alternate;
\item after an up block the bottom may grow by one ($b_{n+1}\in\set{b_n,b_n+1}$, $t_{n+1}=t_n$); after a down block the top may grow by one ($t_{n+1}\in\set{t_n,t_n+1}$,
$b_{n+1}=b_n$).
\end{itemize}
A junction $nL$ is an \emph{event} if the windows of blocks $n-1$ and $n$ differ.
\end{definition}

A sweep walk has steps of at most one (inside a block by crookedness of a stuttered pattern, at a junction because both blocks take the same extreme value there).

\begin{theorem}[Sweep lemma]\label{thm:sweep}
A sweep walk is crooked on every interval $[lo,hi]$ containing at most two events.
\end{theorem}

\begin{proof}
Let $lo\le i,j\le hi$ with $g(i)\le h$, $h+3\le k\le g(j)$.

\emph{The ascending case $i<j$.} Let $t_0$ be the last index in $[i,j]$ with $g\le h$, and $s_0$ the first index in $[t_0,j]$ with $g\ge k$. Unit steps give $g(t_0)=h$, $g(s_0)=k$, and
$h<g<k$ strictly between them. It suffices to find $t_0<r<s<s_0$ with $g(r)=k-1$ and $g(s)=h+1$. Let $\sigma$ be the block containing $t_0$ (the later one, if $t_0$ is a junction).

\emph{Case 1: $s_0\in\sigma$.} Use the crookedness of $\sigma$ on $[t_0,s_0]$.

\emph{Case 2: $s_0\notin\sigma$.} Then $t_0$ is not the last index of $\sigma$ (otherwise the next block starts at the value $h$, at an index in $(t_0,s_0]$, against the choice of
$t_0$). So the last index $e$ of $\sigma$ lies in $(t_0,s_0)$ and $h<g(e)<k$. If $\sigma$ were down, $g(e)=b_\sigma\le g(t_0)=h$; so $\sigma$ is up and $g(e)=t_\sigma<k$. The next block
$\sigma'$ is down with values $\le t_\sigma<k$, so $s_0$ lies beyond it; $\sigma'$ is a full block inside $(t_0,s_0)$ ending at its bottom $b_{\sigma'}>h\ge b_\sigma$. So the junction
$\sigma|\sigma'$ is an event, $h=b_\sigma$ and $b_{\sigma'}=h+1$.
\begin{itemize}
\item \emph{Case 2a: $k-1\le t_\sigma$.} Take $r$ the first index after $t_0$ in $\sigma$ with $g(r)=k-1$, and $s$ the last index of $\sigma'$.
\item \emph{Case 2b: $k-1>t_\sigma$.} To reach $k\ge t_\sigma+2$ the top must grow twice after $\sigma'$, two more events; with $\sigma|\sigma'$ that makes three, which is excluded.
\end{itemize}

\emph{The descending case $j<i$.} Let $s_0$ be the last index in $[j,i]$ with $g\ge k$ and $t_0$ the first index after $s_0$ with $g\le h$; then $g(s_0)=k$, $g(t_0)=h$, and it suffices to find
$s_0<s<r<t_0$ with $g(s)=h+1$, $g(r)=k-1$. Suppose $t_0$ is not in the block $\sigma$ of $s_0$. As before the end $e$ of $\sigma$ lies in $(s_0,t_0)$ with $h<g(e)<k$; if $\sigma$ were up,
$g(e)=t_\sigma\ge k$; so $\sigma$ is down and $g(e)=b_\sigma>h$. The next block $\sigma'$ is up with bottom $b_\sigma$ (bottoms change only after up blocks), so its values are $>h$, $t_0\notin\sigma'$, and
$\sigma'$ is a full block ending at its top $\ge t_\sigma\ge k$ inside $(s_0,t_0)$: a contradiction. So $t_0\in\sigma$, and the crookedness of $\sigma$ on $[s_0,t_0]$ gives the witnesses.
\end{proof}

\formal{\lean{HPC.BJ59.sweep_lemma}.}

\begin{remark}
The lemma holds with at most three events: in Case~2b the top then grows exactly once more before reaching $k-1$, the up block after that growth ends at $k-1$, and the down block after
it ends at the bottom $h+1$, which gives $r$ and $s$. Only two events are needed below. Reversing the indices and negating $g$ maps sweep walks to sweep walks, but it maps the
descending case to itself, so the two cases must be argued separately.
\end{remark}

\formal{\lean{HPC.BJ59.sweep_lemma3}, with \lean{HPC.BJ59.sweep_window_case2b}.}

\section{The pattern}

\begin{definition}\label{def:sweepX}
Let $m\ge5$ and $S\ge1$ be integers, $z$ Bing's numbers (\vone{Definition~4.3}), and put
\[
L=S\,z(m-2)\,z(m-1),\qquad P=3L,\qquad Q=2P=6S\,z(m-2)\,z(m-1).
\]
For $j\in\Z$, \emph{period $j$} is the block of indices $[jP,(j+1)P)$, with window
\[
t_j=\fl{(j+m)/2},\qquad b_j=\cel{(j+2-m)/2}
\]
of width $d_j=t_j-b_j$, equal to $m-1$ if $j+m$ is even and $m-2$ if it is odd. Period $j$ consists of three \emph{sweeps} of length $L$, the sweep $n=3j+r$ ($r=0,1,2$) occupying
$[nL,(n+1)L)$. With $K_j=L/z(d_j)$ and $F_d$ Bing's pattern of length $z(d)$ from $0$ to $d$ (\vone{Definition~4.4}), sweep $n$ at position $u\in[0,L)$ has the value
\[
b_j+F_{d_j}(\fl{u/K_j})\ \text{ if $n+m$ is even (up)},\qquad t_j-F_{d_j}(\fl{u/K_j})\ \text{ if $n+m$ is odd (down)}.
\]
This defines the \emph{sweep pattern} $x=x_{m,S}\colon\Z\to\Z$.
\end{definition}

\formal{\lean{HPC.BJ59.sweepX}, with \lean{HPC.BJ59.sweepL}, \lean{HPC.BJ59.sweepQ}, \lean{HPC.BJ59.swT}, \lean{HPC.BJ59.swB}, \lean{HPC.BJ59.swD},
\lean{HPC.BJ59.swK}, \lean{HPC.BJ59.swVal}.}

Each sweep is a stuttered Bing pattern (stutter $K_j\ge S$) across the window of its period. From period $j$ to period $j+1$ the top grows when $j+m$ is odd and the bottom grows when $j+m$ is
even, one end at a time; the period after a top growth starts at the bottom and the period after a bottom growth starts at the top. So $x$ is a sweep walk with blocks of length $L$, whose events lie
at multiples of $P$.

\formal{\lean{HPC.BJ59.sw_win_up}, \lean{HPC.BJ59.sw_win_down}, \lean{HPC.BJ59.swVal_junction}, \lean{HPC.BJ59.sweepX_block_crooked}.}

\begin{theorem}\label{thm:sweep-pattern}
For $m\ge5$, $S\ge1$, the sweep pattern $x$ satisfies, with $y(t)=1/4+t/Q$ and $\rho=m/2-1/4$:
\begin{enumerate}[label=(\alph*)]
\item (period) $x(t+Q)=x(t)+1$;
\item (unit steps) $\abs{x(t+1)-x(t)}\le1$;
\item (slow) if $x(s)+2\le x(t)$ then $\abs{s-t}\ge S$;
\item (nesting) $\abs{x(t)-y(t)}\le\rho$;
\item (survival) if $\abs{y-l}\le m/2-1/2$, there is $t$ with $x(t)=l$ and $\abs{y-y(t)}\le1/4-1/(2Q)$;
\item (crookedness) for every $a\in\Z$, $u\mapsto x(a+u)$ is crooked on $[0,P+3]$.
\end{enumerate}
\end{theorem}

\begin{proof}
(a) $t_{j+2}=t_j+1$, $b_{j+2}=b_j+1$, $d_{j+2}=d_j$, and $Q$ is two periods.

(b) Inside a sweep, by \vone{Lemma~4.9}; at a junction both sweeps take the same extreme value.

(c) The walk is a sequence of maximal constant blocks of length at least $S$ (the stutter is a multiple of $S$), consecutive blocks differing by one.

(d) In period $j$, $y(t)\in[e_j,e_{j+1})$ with $e_j=1/4+j/2$, and $b_j\le x(t)\le t_j$, where $t_j\le(j+m)/2=e_j+\rho$ and $b_j\ge(j+2-m)/2=e_{j+1}-\rho$.

(e) The periods whose window contains $l$ are $j=2l-m,\dots,2l+m-2$, covering the $y$-range $[l-\rho,l+\rho)$, and every sweep of these periods takes the value $l$. In the first,
$j_1=2l-m$, $l=t_{j_1}$ is a new top and the period starts at the bottom; its first sweep reaches $l$ at a $y$ below $l-\rho+1/6$. In the last, $j_2=2l+m-2$, $l=b_{j_2}$ and its third sweep
starts at $l$, at $y=l+\rho-1/6$. Consecutive visits to $l$ are less than $2L$ indices apart, a $y$-gap below $1/3\le1/2-1/Q$. So the intervals $[y(t)-1/4+1/(2Q),\,y(t)+1/4-1/(2Q)]$ over
the visits $t$ cover $[l-\rho+1/4,l+\rho-1/4]=[l-m/2+1/2,\,l+m/2-1/2]$.

(f) An interval of $P+4$ consecutive indices contains at most two multiples of $P$ (as $P\ge4$), hence at most two events. Apply Theorem~\ref{thm:sweep}.
\end{proof}

\formal{\lean{HPC.BJ59.exists_sweepPattern}, from \lean{HPC.BJ59.sweepX_period}, \lean{HPC.BJ59.sweepX_step}, \lean{HPC.BJ59.sweepX_slow}, \lean{HPC.BJ59.sweepX_nest},
\lean{HPC.BJ59.sweepX_surv}, \lean{HPC.BJ59.sweepX_crook}.}

\begin{remark}
Three sweeps per period are needed for survival: with one sweep, consecutive visits to a link can be a whole period apart (a $y$-gap of $1/2$, more than $1/2-1/Q$). In
Chapter~\ref{ch:plane} the new level has $\lambda N=Q/2=P$, so the window $\cel{\lambda_{k+1}N_{k+1}}+3$ of condition (crook) is exactly $P+3$.
\end{remark}

\section*{Exercises}

\begin{exercise}
Compute $Q$ for $m=5$, $S=1$. (Use $z(3)$ and $z(4)$ from \vone{Definition~4.3}.)
\end{exercise}

\begin{exercise}
Give a sweep walk with one sweep per window change which is not crooked on an interval containing four events.
\end{exercise}

\begin{exercise}
Check that the window of period $j$ consists exactly of the integers $l$ with $\abs{l-y}\le\rho$ for every $y\in[e_j,e_{j+1}]$.
\end{exercise}

%======================================================================
\chapter{A circle of pseudo-arcs in the plane}\label{ch:plane}
%======================================================================

We realize the circular descriptions of Chapter~\ref{ch:circdescr} in the plane, using the sweep patterns of Chapter~\ref{ch:sweep}. The construction is
the circular analogue of the plane pseudo-arc of \vone{Chapter~5}: there each new chain was obtained from a standard chain by a homeomorphism of the plane
built from joint positions and a shear (\vone{Proposition~5.14}). Around a circle the linear shear of Volume~1 cannot close up; it is replaced by a shear on the
universal cover of the punctured plane, which commutes with the deck translations and so descends.

The plane carries the sup metric, as in Volume~1.

\section{Standard circular chains}

\begin{definition}\label{def:std-circ}
For $t\in\Z$ let $R(t)=(2t,2t+3)\times(0,1)$ and $Z(t)=[2t+8/5,\,2t+17/5]\times[0,1]$ (the standard rectangles and zones of \vone{Definition~5.1}, with integer index). For
$N\ge1$ let $\cexp_N(u,v)=(e^v\cos(\pi u/N),\,e^v\sin(\pi u/N))$, a covering map of $\R^2$ onto $\R^2\setminus\set0$ with period $2N$ in $u$. The \emph{standard circular chain}
with $N$ links has links $\cexp_N(R(t))$ and zones $\cexp_N(Z(t))$, $t\in\Z$.
\end{definition}

\formal{\lean{HPC.BJ59.stdRectZ}, \lean{HPC.BJ59.zoneZ}, \lean{HPC.BJ59.circExp}, \lean{HPC.BJ59.stdCircLink}, \lean{HPC.BJ59.stdCircZone}.}

The links are open, preconnected (images of rectangles under an open map), lie in the closed annulus $1\le\abs p\le e$, and satisfy $\cexp_N(R(t+N))=\cexp_N(R(t))$.

\formal{\lean{HPC.BJ59.CircStd.isOpen_stdCircLink}, \lean{HPC.BJ59.CircStd.isPreconnected_stdCircLink}, \lean{HPC.BJ59.CircStd.stdCircLink_periodic}.}

\begin{lemma}\label{lem:std-meet}
For $N\ge4$, $\cexp_N(R(s))$ meets $\cexp_N(R(t))$ if and only if $s-t$ is within $1$ of a multiple of $N$.
\end{lemma}

\begin{proof}
Two points have the same image if and only if their $v$-coordinates agree and their $u$-coordinates differ by a multiple of $2N$. The intervals $(2s,2s+3)$ and
$(2t+2Nn,2t+2Nn+3)$ meet if and only if $\abs{s-t-Nn}\le1$.
\end{proof}

\formal{\lean{HPC.BJ59.stdCircLink_meet_iff}.}

\section{The circular planar step}

\begin{theorem}[The circular planar step]\label{thm:circ-step}
Let $\Phi$ be a homeomorphism of the plane, $N\ge4$, $\varepsilon>0$ and $B\in\R$. There is $S\ge12$ with the following property. Let $Q\ge2$, $y_0\in\R$, and $x\colon\Z\to\Z$ with
$x(t+Q)=x(t)+1$, unit steps, $S$-slow ($x(s)+2\le x(t)$ implies $\abs{s-t}\ge S$) and $\abs{x(t)-(y_0+t/Q)}\le B$. Then there is a homeomorphism $\psi$ of the plane such that,
writing $L'_t=\Phi\psi\cexp_{QN}(R(t))$ for the new links and $L_l=\Phi\cexp_N(R(l))$ for the old ones:
\begin{enumerate}[label=(\arabic*)]
\item $\cl{L'_t}\subseteq L_{x(t)}$;
\item if $L'_t$ meets $L_l$, then $x(t')=l+nN$ for some $\abs{t'-t}\le1$ and $n\in\Z$;
\item each $L'_t$ has diameter $\le\varepsilon$ or lies in an old zone $\Phi\cexp_N(Z(k))$;
\item each new zone $\Phi\psi\cexp_{QN}(Z(t))$ has diameter $\le\varepsilon$.
\end{enumerate}
\end{theorem}

\formal{\lean{HPC.BJ59.circ_step}.}

\begin{proof}[Proof in outline]
\emph{Joint positions.} Define $P(t)\in\R$ for $t\in\Z$ by the formulas of \vone{Definition~5.6} applied to the lifted pattern $x$: at a \emph{stay} ($x(t-1)=x(t)$),
$P(t)=2x(t)+3/2+\varphi(t)/4$, where $\varphi(t)=(\ell-r)/(\ell+r)$ is computed from the (doubled, capped at $2S$) distances $\ell$, $r$ to the nearest indices with values $x(t)\mp1$;
at a \emph{transition}, $P(t)=2\min(x(t-1),x(t))+5/2$. The formulas are local, so $P(t+QN)=P(t)+2N$, and the link lemma of \vone{Lemma~5.8}, which concerns a bounded window of
indices, transfers: it describes the horizontal position of the new links relative to the old ones.

\emph{Interpolation.} Let $\Ph\colon\R\to\R$ equal $P(t)$ on $[2t-2/5,\,2t+7/5]$ and be linear from $P(t)$ to $P(t+1)$ on $[2t+7/5,\,2t+8/5]$. It is continuous, $\Ph(u+2QN)=\Ph(u)+2N$,
and $\Ph$ is constant on the $u$-range of each zone.

\emph{The shear on the cover.} For small constants $\varepsilon_0,C>0$ put
\[
\widetilde\Psi(u,v)=\Bigl(\Ph(u)+\varepsilon_0v,\ \tfrac12+C\bigl(\Ph(u)-u/Q-c\bigr)+C\varepsilon_0v\Bigr),\qquad c=2y_0+3/2.
\]
It is a homeomorphism of $\R^2$: it is the vertical shear $(u,v)\mapsto(u,\Ph(u)-u/Q+\varepsilon_0v)$ followed by the linear map $(u,w)\mapsto(u/Q+w,\,1/2+C(w-c))$, of determinant
$C/Q\neq0$. It commutes with the deck translations: $\widetilde\Psi(u+2QN,v)=\widetilde\Psi(u,v)+(2N,0)$.

\emph{Heights.} $\abs{\Ph(u)-u/Q-c}\le2B+4$, so for $C(2B+5)\le1/8$ and $\varepsilon_0\le1/20$, $\widetilde\Psi$ maps $\R\times[0,1]$ into heights in $(1/4,3/4)$. Two passes of the new chain
over the same place are thereby separated vertically by $C$ times the drift between them.

\emph{Descent.} Define $\psi(\cexp_{QN}(p))=\cexp_N(\widetilde\Psi(p))$ and $\psi(0)=0$. By equivariance (of $\widetilde\Psi$ and its inverse) $\psi$ is a well-defined bijection; it is
continuous away from $0$ because $\cexp$ is a local homeomorphism, and continuous at $0$ (with its inverse) because the second coordinate of $\widetilde\Psi(u,v)$ is $C\varepsilon_0v$ plus a
bounded function of $u$, so $\abs{\psi(w)}\le e^{B'}\abs w^{C\varepsilon_0}$.

\emph{The four properties.} On the cover, $\widetilde\Psi(R(t))$ has horizontal projection in $[\min(P(t),P(t+1)),\,\max(P(t),P(t+1))+\varepsilon_0]$ and heights in $(1/4,3/4)$. Containment
and tightness are then the horizontal analysis of the link lemma, read modulo $2N$ through $\cexp_N$. A stay link has image of width $\abs{P(t+1)-P(t)}+\varepsilon_0$ and height
$C(\abs{P(t+1)-P(t)}+3/Q+\varepsilon_0)$, and stay steps are $O(1/S)$, so stay links are small for large $S$; transition links lie in old zones; the image of a zone lies in a box of width
$\varepsilon_0$ and height $C(9/(5Q)+\varepsilon_0)$. Uniform continuity of $\cexp_N$ on $\R\times[0,1]$ (it is periodic) and of $\Phi$ on the compact annulus $1\le\abs p\le e$ converts these
bounds into (3) and (4). The constants are chosen in the order $\varepsilon_0$, $C$, then $S$.
\end{proof}

\formal{The joint positions are \lean{HPC.BJ59.CircJ.jointZ}, with the link lemma \lean{HPC.BJ59.CircJ.link_facts_Z}; the interpolation is \lean{HPC.BJ59.Ramp.phat}; the
shear is \lean{HPC.BJ59.cover}; the descent is \lean{HPC.BJ59.Descent.descHomeo}, with \lean{HPC.BJ59.Descent.E}; the estimates are
\lean{HPC.BJ59.CircStepAux.small_images} and \lean{HPC.BJ59.CircStepAux.phat_height}.}

\section{The assembly}

\begin{definition}\label{def:levels}
A \emph{level} is a quadruple $(N,\Phi,m,o)$ with $N\ge4$, $\Phi$ a homeomorphism of the plane, $m\ge5$, $m+4\le N$, and $o\in\R$. Level $0$ is $(16,\mathrm{id},5,0)$. From level $i$, with
$\eta_i=2^{-i}$ and $B=m/2-1/4$: take $S$ from Theorem~\ref{thm:circ-step} for $(\Phi,N,\eta_i,B)$; take the sweep pattern $x=x_{m,S}$ of Definition~\ref{def:sweepX}, with $Q$ its
period; take $\psi$ from Theorem~\ref{thm:circ-step} with $y_0=1/4$. Level $i+1$ is
\[
\bigl(QN,\ \Phi\circ\psi,\ Q/2,\ o+1/(4N)\bigr).
\]
\end{definition}

Then $Q/2=3Sz(m-2)z(m-1)\ge5$ and $Q/2+4\le QN$, so level $i+1$ is again a level.

\formal{\lean{HPC.BJ59.CircPlane.Lvl}, \lean{HPC.BJ59.CircPlane.L0}, \lean{HPC.BJ59.CircPlane.next}, \lean{HPC.BJ59.CircPlane.lvl}.}

\begin{theorem}\label{thm:circdescr-plane}
Let $(N_k,\Phi_k,m_k,o_k)$ be the levels. The links $L_k(t)=\Phi_k\cexp_{N_k}(R(t))$, the patterns $\pi_k$ (the sweep pattern of step $k$), the offsets $o_k$ and the arc lengths
$\lambda_k=m_k/N_k$ form a circular description in the plane.
\end{theorem}

\begin{proof}
\emph{Links.} $K$ is the closed annulus $1\le\abs p\le e$: level-$0$ links lie in it, and all others by (nest). Openness, preconnectedness, periodicity and the meeting rule come from the standard
circular chains (Lemma~\ref{lem:std-meet}).

\emph{Patterns.} (nest) and (tight) are properties (1), (2) of Theorem~\ref{thm:circ-step}. Degree one: iterating $x(t+Q)=x(t)+1$ $N_k$ times gives $x(t+QN_k)=x(t)+N_k$. Unit steps: Theorem~\ref{thm:sweep-pattern}(b).

\emph{Mesh.} Level-$0$ links and zones lie in the closed sup-ball of radius $e$, so have diameter at most $2e<8$. A level-$(i+1)$ link has diameter $\le\eta_i$ or lies in a level-$i$ zone, whose
diameter is $\le\eta_{i-1}$ (property (4) at the previous step) or $<8$ when $i=0$. So $\diam L_k(t)\le16/2^k$.

\emph{Arcs.} $\lambda_{k+1}=1/(2N_k)\to0$, and $5\le m_k=\lambda_kN_k$, $\lambda_kN_k+4\le N_k$ by the definition of a level. (arcNest): $\abs{c_{k+1}(t)-c_k(x(t))}=\abs{1/4+t/Q-x(t)}/N_k\le(m_k/2-1/4)/N_k=(\lambda_k-\lambda_{k+1})/2$ by
Theorem~\ref{thm:sweep-pattern}(d). (surv): with $y=N_k(z-o_k)$ the hypothesis reads $\abs{y-l}\le m_k/2-1/2$, and Theorem~\ref{thm:sweep-pattern}(e) gives the conclusion multiplied by $N_k$ (note
$\lambda_{k+1}/2=1/(4N_k)$, $\delta_{k+1}=1/(2QN_k)$). (crook): $\lambda_{k+1}N_{k+1}=Q/2=P$, and Theorem~\ref{thm:sweep-pattern}(f).
\end{proof}

\formal{\lean{HPC.BJ59.planeDescr}, \lean{HPC.BJ59.exists_circDescr_plane}, with \lean{HPC.BJ59.CircPlane.link_diam}, \lean{HPC.BJ59.CircPlane.arcNest},
\lean{HPC.BJ59.CircPlane.surv}, \lean{HPC.BJ59.CircPlane.crook}.}

\begin{theorem}[Bing and Jones, \S7]\label{thm:plane-example}
Some plane continuum is a circle of pseudo-arcs.
\end{theorem}

\begin{proof}
Theorems~\ref{thm:circdescr-plane} and~\ref{thm:circdescr}.
\end{proof}

\formal{\lean{HPC.BJ59.bjPlaneExample}, a proof of the statement \lean{HPC.BJ59.BJPlaneExample}.}

\section*{Exercises}

\begin{exercise}
Show that a linear shear $(u,v)\mapsto(u,v+cu)$ with $c\neq0$ does not commute with the translation $u\mapsto u+2QN$, and so does not descend through $\cexp_{QN}$.
\end{exercise}

\begin{exercise}
Verify that $\widetilde\Psi$ is injective directly: if $\widetilde\Psi(u,v)=\widetilde\Psi(u',v')$, subtract $1/C$ times the second coordinate from the first.
\end{exercise}

\begin{exercise}
Show that the continuum of Theorem~\ref{thm:plane-example} separates the plane into exactly two components. (It is homogeneous and decomposable; use \vseven{Chapter~9}.)
\end{exercise}

%======================================================================
\chapter{The classification theorem}\label{ch:classification}
%======================================================================

All the pieces are now in place. A \emph{simple closed curve} is a space homeomorphic to $\Sph$.

\section{The decomposable case}

\begin{theorem}\label{thm:decomp-case}
A decomposable homogeneous plane continuum is a simple closed curve or a circle of pseudo-arcs.
\end{theorem}

\begin{proof}
Let $M$ be such a continuum. By Theorem~\ref{thm:ho-circlelike}, $M$ is circle-like. By \vseven{Theorem~9.8}, $M$ is a simple closed curve, or there is a continuous, open, surjective map
$f\colon M\to\Sph$ whose fibres are pseudo-arcs. In the second case $M$, being circle-like, is a circle of pseudo-arcs (Definition~\ref{def:circle-pa}).
\end{proof}

\formal{\lean{HPC.BJ59.circle_or_circleOfPA_of_decomposable} (with circle-likeness as the hypothesis \lean{HPC.BJ59.HOCircleLike}, discharged by
\lean{HPC.BJ59.hoCircleLike'}), from \lean{HPC.BJ59.circle_or_circPAMap_of_decomposable} and \lean{HPC.BJ59.circle_or_atomic_of_decomposable}.}

\section{The classification}

\begin{theorem}[The classification of homogeneous plane continua]\label{thm:classification}
A nondegenerate plane continuum is homogeneous if and only if it is a simple closed curve, a pseudo-arc, or a circle of pseudo-arcs.
\end{theorem}

\begin{proof}
Let $M$ be a nondegenerate plane continuum.

($\Rightarrow$) Suppose $M$ is homogeneous. If $M$ is decomposable, Theorem~\ref{thm:decomp-case} applies. If $M$ is indecomposable, it is a pseudo-arc by \vsix{Theorem~7.4}.

($\Leftarrow$) A simple closed curve is homogeneous (Proposition~\ref{prop:circle-homog}). A pseudo-arc is homogeneous (\vtwo{Theorem~8.4}). A circle of pseudo-arcs is homogeneous
(Theorem~\ref{thm:homog}).
\end{proof}

\formal{\leansub{HPC.BJ59.ho16_classification}{4}. The indecomposable case in both directions is \lean{HPC.HO16.ho16_indec_homogeneous_iff}, and the converse for pseudo-arcs
\lean{HPC.HO16.ho16_pseudoArc_indec_homogeneous}.}

By Theorem~\ref{thm:dp61}, the third class may equally be described as in Domokos and Prajs: compact spaces with a continuous, open, surjective, atomic map onto $\Sph$ whose fibres are
pseudo-arcs.

\formal{\lean{HPC.BJ59.ho16_classification_atomic'}.}

\begin{theorem}[The plane form]\label{thm:classification-plane}
There are plane continua $C_1$, $C_2$, $C_3$ such that
\begin{enumerate}[label=(\alph*)]
\item $C_1$ is a simple closed curve, $C_2$ is a pseudo-arc and $C_3$ is a circle of pseudo-arcs;
\item each of them is homogeneous;
\item no two of them are homeomorphic;
\item a nondegenerate plane continuum is homogeneous if and only if it is homeomorphic to $C_1$, to $C_2$ or to $C_3$.
\end{enumerate}
\end{theorem}

\begin{proof}
(a) Take the unit circle for $C_1$ (Proposition~\ref{prop:circle-homog}), the plane pseudo-arc of \vone{Theorem~5.18} for $C_2$, and the plane circle of pseudo-arcs of
Theorem~\ref{thm:plane-example} for $C_3$.

(b) Theorem~\ref{thm:classification}, ($\Leftarrow$).

(c) $C_1$ and $C_3$ are decomposable (Proposition~\ref{prop:circle-decomp}, Corollary~\ref{cor:circle-pa-decomp}), while $C_2$ is hereditarily indecomposable; so $C_2$ is homeomorphic to
neither. $C_1$ and $C_3$ are not homeomorphic by Corollary~\ref{cor:circle-not-cpa}.

(d) A continuum homeomorphic to a homogeneous one is homogeneous. Conversely, if $M$ is homogeneous, Theorem~\ref{thm:classification} makes it a simple closed curve (homeomorphic to $C_1$), a
pseudo-arc (homeomorphic to $C_2$ by \vtwo{Theorem~8.3}) or a circle of pseudo-arcs (homeomorphic to $C_3$ by Theorem~\ref{thm:unique}).
\end{proof}

\formal{\leansub{HPC.BJ59.ho16_classification_plane}{4}, with \lean{HPC.BJ59.exists_plane_circle}, \lean{HPC.Plane.isPseudoArc_plane}, \lean{HPC.BJ59.bjPlaneExample},
\lean{HPC.Plane.theorem1_HIC'}.}

\section{How the inputs were discharged}

The formal development first proved the classification with three named statements as hypotheses, and then proved the statements one at a time:
\begin{center}
\begin{tabular}{@{}p{0.45\textwidth}p{0.48\textwidth}@{}}
\toprule
Statement & Proved in this volume as\\
\midrule
Theorem~10 of Bing and Jones (\lean{HPC.BJ59.BJThm10}) & Theorem~\ref{thm:bj10}\\
Decomposable homogeneous plane continua are circle-like (\lean{HPC.BJ59.HOCircleLike}) & Theorem~\ref{thm:ho-circlelike}\\
A plane circle of pseudo-arcs exists (\lean{HPC.BJ59.BJPlaneExample}) & Theorem~\ref{thm:plane-example}\\
\bottomrule
\end{tabular}
\end{center}
The classification with all three as hypotheses is \lean{HPC.BJ59.ho16_classification} (and \lean{HPC.BJ59.ho16_classification_plane}); the circle-likeness statement is then derived from
the other two (\leansub{HPC.BJ59.ho16_classification}{2}), Theorem~10 is discharged (\leansub{HPC.BJ59.ho16_classification}{3}), and finally the plane example
(\leansub{HPC.BJ59.ho16_classification}{4}); likewise for the plane form (\leansub{HPC.BJ59.ho16_classification_plane}{2}, \leansub{HPC.BJ59.ho16_classification_plane}{3},
\leansub{HPC.BJ59.ho16_classification_plane}{4}). The final two theorems carry no hypotheses beyond those in their statements.

\formal{\leansub{HPC.BJ59.ho16_classification}{4}, \leansub{HPC.BJ59.ho16_classification_plane}{4}.}

\section{The whole proof at a glance}

\begin{enumerate}
\item \emph{Indecomposable case} (Volume~6): a homogeneous indecomposable plane continuum is hereditarily indecomposable (Hagopian, Jones), has span zero (Oversteegen and Tymchatyn) and is
snake-like (Hoehn and Oversteegen); hence a pseudo-arc (Volumes~1 and~2).
\item \emph{Decomposable case} (Volume~7): Jones's aposyndetic decomposition gives a continuous decomposition into terminal indecomposable continua with a circle as quotient, or $M$ is a simple
closed curve; the elements are pseudo-arcs.
\item \emph{Circles of pseudo-arcs} (this volume): the quotient map is atomic, so its halves are snake-like (Chapter~\ref{ch:atomic}); Theorem~10 of Bing and Jones (Chapters~\ref{ch:chains}--\ref{ch:thm10})
makes any two such spaces homeomorphic (Chapter~\ref{ch:uniqueness}); a plane example (Chapters~\ref{ch:circdescr}--\ref{ch:plane}) shows that they are circle-like and that the third class is
nonempty.
\item Plane topology (Volumes~4 and~5) and general continuum theory (Volume~3) are used throughout.
\end{enumerate}

\section*{Exercises}

\begin{exercise}
Which of the three continua $C_1$, $C_2$, $C_3$ are snake-like? Which are circle-like? Which are hereditarily indecomposable?
\end{exercise}

\begin{exercise}
Show that the product $P\times\Sph$ of a pseudo-arc and a circle is homogeneous but is not homeomorphic to any plane continuum. (Use Theorem~\ref{thm:classification-plane}.)
\end{exercise}

\begin{exercise}
Show that a homogeneous plane continuum which contains an arc is a simple closed curve.
\end{exercise}


\appendix
%======================================================================
\chapter{Correspondence with the formal proof}\label{app:formal}
%======================================================================

The table lists the definitions and results of this volume that carry a ``Formal counterpart''
line, together with the corresponding declarations of the Lean development. All declaration
names start with \lean{HPC.}. The formal statements are sometimes slightly more general, or are
split differently into lemmas; the table points to the place where the content of each result is
proved. Rows are in the order of the text; rows labelled ``Text of \dots'' belong to formal counterpart lines in the
running text of the named section.

The main theorems of the volume are, in the Lean development (prefix \lean{HPC.} omitted):
\begin{center}
\begin{tabular}{@{}p{0.42\textwidth}p{0.54\textwidth}@{}}
\toprule
Atomic arcs are snake-like (Theorem~\ref{thm:atomic-snake}) & \lean{BJ59.isSnakeLike_of_atomic'}\\
Bing 1951, Theorem~10; Fugate, Lemma~1 (Theorem~\ref{thm:fugate1}) & \lean{Plane.fugate_lemma1}\\
Shrinking to a BJ-chain (Theorem~\ref{thm:shrink}) & \lean{BJ59.exists_bjChain_shrink}\\
Theorem~6, the form used (Theorem~\ref{thm:bj6}) & \lean{BJ59.adjacent_of_properCover}\\
Bing and Jones, Theorems~3 and~4 (Theorems~\ref{thm:bj3}, \ref{thm:bj4}) & \lean{BJ59.theorem3_arc}, \lean{BJ59.bj_theorem4}\\
Theorem~5 is false as printed (Proposition~\ref{prop:bj5-false}) & \lean{BJ59.bj_theorem5_printed_false_general}\\
Theorem~5, repaired (Theorem~\ref{thm:bj5}) & \lean{BJ59.theorem5_repaired}\\
Theorem~7 (Theorems~\ref{thm:bj7}, \ref{thm:bj7-ends}) & \lean{BJ59.bj_theorem7'}, \lean{BJ59.bj_theorem7_ends'}\\
Theorem~8 and its corollary (Theorem~\ref{thm:bj8}, Corollary~\ref{cor:bj8}) & \lean{BJ59.bj_theorem8}, \lean{BJ59.bj_theorem8_cor}\\
The homeomorphism from a zigzag (Theorem~\ref{thm:zigzag}) & \lean{BJ59.zigzag_homeomorph}\\
The step (Theorem~\ref{thm:step}) & \lean{BJ59.bj_step}\\
Bing and Jones, Theorem~10 (Theorem~\ref{thm:bj10}) & \lean{BJ59.bjThm10}\\
Uniqueness of circles of pseudo-arcs (Theorem~\ref{thm:unique}) & \lean{BJ59.homeomorph_of_isCircleOfPA'}\\
Homogeneity of circles of pseudo-arcs (Theorem~\ref{thm:homog}) & \lean{BJ59.isHomogeneous_of_isCircleOfPA'}\\
Domokos and Prajs, Theorem~6.1 (Theorem~\ref{thm:dp61}) & \lean{BJ59.dp22_theorem61'}, \lean{BJ59.dp22_theorem61_iff'}\\
Circle-likeness (Theorem~\ref{thm:ho-circlelike}) & \lean{BJ59.hoCircleLike'}\\
Circular descriptions (Theorem~\ref{thm:circdescr}) & \lean{BJ59.CircDescr.isCircleOfPA}\\
Sweep lemma, sweep patterns (Theorems~\ref{thm:sweep}, \ref{thm:sweep-pattern}) & \lean{BJ59.sweep_lemma}, \lean{BJ59.exists_sweepPattern}\\
The circular planar step (Theorem~\ref{thm:circ-step}) & \lean{BJ59.circ_step}\\
A plane circle of pseudo-arcs (Theorem~\ref{thm:plane-example}) & \lean{BJ59.bjPlaneExample}\\
The decomposable case (Theorem~\ref{thm:decomp-case}) & \lean{BJ59.circle_or_circleOfPA_of_decomposable}, \lean{BJ59.hoCircleLike'}\\
The classification (Theorem~\ref{thm:classification}) & \leansub{BJ59.ho16_classification}{4}\\
The plane form (Theorem~\ref{thm:classification-plane}) & \leansub{BJ59.ho16_classification_plane}{4}\\
\bottomrule
\end{tabular}
\end{center}
All of them are proved in the Lean development without \lean{sorry} and without axioms beyond the
standard ones of Lean and Mathlib. The printed forms that are false (Section~\ref{sec:intro-diff}) are
refuted by the declarations whose names contain \lean{printed_false} or \lean{proof_n3}.

\begin{longtable}{@{}p{0.40\textwidth}p{0.56\textwidth}@{}}
\toprule
In these notes & In the Lean development\\
\midrule
\endhead
\bottomrule
\endfoot
\multicolumn{2}{@{}l}{\emph{Chapter~\ref{ch:circles}}}\\
Definition~\ref{def:arc-continua} & \lean{HPC.BJ59.IsArcOfContinua}, \lean{HPC.BJ59.IsContArcPA}.\\
Definition~\ref{def:circle-pa} & \lean{HPC.BJ59.IsCircleOfPA}, \lean{HPC.BJ59.IsCircPAMap}; the equivalence of ``$X$ is a circle of pseudo-arcs'' with ``$X$ carries a circle-of-pseudo-arcs map'' is \lean{HPC.BJ59.isCircleOfPA_iff}.\\
Example~\ref{ex:not-atomic} & \lean{HPC.BJ59.not_isAtomicMap_snd}, \lean{HPC.BJ59.isAtomicMap_id}.\\
Lemma~\ref{lem:monotone-preimage} & \lean{HPC.BJ59.isPreconnected_preimage_of_monotone'}.\\
Corollary~\ref{cor:G-connected} & \lean{HPC.BJ59.isPreconnected_preimage_Icc}, \lean{HPC.BJ59.isCompact_preimage_Icc}.\\
Definition~\ref{def:halves} & \lean{HPC.Jones.halfC}, \lean{HPC.BJ59.halfParam}, \lean{HPC.BJ59.halfSet}, \lean{HPC.BJ59.halfMap}.\\
Text of ``The two halves'' & \lean{HPC.BJ59.range_halfC_union}, \lean{HPC.BJ59.mem_range_halfC_inter}, \lean{HPC.BJ59.exists_not_mem_range_halfC}, \lean{HPC.BJ59.image_val_halfMap_fibre}, \lean{HPC.BJ59.halfFibreEquiv}.\\
Proposition~\ref{prop:halves} & \lean{HPC.BJ59.isContArcPA_halfMap'} (with \lean{HPC.BJ59.continuous_halfMap}, \lean{HPC.BJ59.isOpenMap_halfMap}, \lean{HPC.BJ59.surjective_halfMap}), \lean{HPC.BJ59.isContinuum_halfSet}, \lean{HPC.BJ59.isAtomicMap_halfMap}.\\
Proposition~\ref{prop:halves-snake} & \lean{HPC.BJ59.isContArcPA_halfMap}, \lean{HPC.BJ59.isSnakeLike_halfSet}, from \lean{HPC.Plane.isSnakeLike_of_subset_circleLike}.\\
Definition~\ref{def:snake-halves} & \lean{HPC.BJ59.IsSnakeHalvesMap}; a circle-of-pseudo-arcs map has snake-like halves by \lean{HPC.BJ59.IsCircPAMap.isSnakeHalvesMap}, and an atomic one on a compact space by \lean{HPC.BJ59.IsAtomicCircPAMap.isSnakeHalvesMap} (Chapter~\ref{ch:atomic}).\\
Text of ``The two halves'' & \lean{HPC.BJ59.IsCircPAMap.mul_left}, \lean{HPC.BJ59.IsAtomicCircPAMap.mul_left}.\\
Proposition~\ref{prop:circle-decomp} & \lean{HPC.BJ59.isDecomposableSpace_circle}.\\
Corollary~\ref{cor:circle-pa-decomp} & \lean{HPC.BJ59.isDecomposableSpace_of_isCircPAMap}.\\
Proposition~\ref{prop:pa-not-in-circle} & \lean{HPC.BJ59.not_isPseudoArc_of_ne_univ_circle}, with \lean{HPC.BJ59.continuousAt_arg_circle}, \lean{HPC.BJ59.arg_injective_circle} and \lean{HPC.BJ59.not_isHereditarilyIndecomposable_of_nontrivial_real}.\\
Corollary~\ref{cor:circle-not-cpa} & Inside the proof of \lean{HPC.BJ59.ho16_classification_plane}, with \lean{HPC.BJ59.isPseudoArc_of_homeomorph}.\\
Proposition~\ref{prop:circle-homog} & \lean{HPC.BJ59.isHomogeneous_circle'}, \lean{HPC.BJ59.isHomogeneous_circle}, \lean{HPC.Span.isHomogeneous_of_homeomorph}, \lean{HPC.BJ59.exists_plane_circle}.\\[4pt]
\multicolumn{2}{@{}l}{\emph{Chapter~\ref{ch:atomic}}}\\
Text of ``chapter introduction'' & \lean{HPC.Plane.EpsCoverable}, \lean{HPC.Plane.isEpsChain_of_refines}.\\
Lemma~\ref{lem:saturation} & \lean{HPC.BJ59.subset_fibre_or_saturated}, \lean{HPC.BJ59.fibre_subset_of_saturated}.\\
Text of ``Atomic maps onto the interval'' & \lean{HPC.BJ59.isPreconnected_preimage_of_fibres}, \lean{HPC.BJ59.isContinuum_univ_of_fibres}.\\
Lemma~\ref{lem:condB} & \lean{HPC.BJ59.condB_preimage_Icc}, with \lean{HPC.BJ59.nested_of_common_bound}.\\
Proposition~\ref{prop:atomic-hu} & \lean{HPC.BJ59.isHereditarilyUnicoherent_of_atomic}, with \lean{HPC.Plane.isHereditarilyUnicoherent_of_isSnakeLike}.\\
Lemma~\ref{lem:brouwer} & \lean{HPC.Plane.exists_minimal_notCoverable}.\\
Text of ``Bing's minimal counterexample'' & \lean{HPC.Plane.IsTerminalSub}, \lean{HPC.Plane.isTerminalSub_singleton_iff}. \lean{HPC.Plane.IsT10Chain}.\\
Lemma~\ref{lem:surgery} & \lean{HPC.Plane.surgery_of_sep} (Bing's original form with condition B is \lean{HPC.Plane.theorem10_1951_of_lastLink}); the trimming is \lean{HPC.Plane.exists_chain_of_weak}.\\
Lemma~\ref{lem:composant-step} & \lean{HPC.Plane.exists_proper_superset_meeting_all_links}, with \lean{HPC.Plane.IsTerminalSub.mono}. \lean{HPC.Plane.fugate_lemma1}, from \lean{HPC.Plane.lemma1_aux}, \lean{HPC.Plane.t10Chain_of_mem_first}, \lean{HPC.Plane.t10Chain_of_mem_last}. Fugate assumes $M$ atriodic and hereditarily unicoherent; the formal proof does not use either hypothesis.\\
Proposition~\ref{prop:glue-halves} & \lean{HPC.Plane.epsCoverable_of_halves}, with Claim~2 as \lean{HPC.Plane.claim2}. Fugate uses $\varepsilon/2$-chains; no mesh is lost.\\
Theorem~\ref{thm:atomic-snake} & \lean{HPC.BJ59.isSnakeLike_of_atomic'}; the case of pseudo-arc fibres is \lean{HPC.BJ59.isSnakeLike_of_atomic}. Openness of $f$ is not needed. Bing's own theorem, for his ``layers'', is \lean{HPC.Plane.theorem11_1951}.\\
Corollary~\ref{cor:atomic-halves} & \lean{HPC.BJ59.isContArcPA_halfMap_of_atomic}, \lean{HPC.BJ59.isSnakeLike_halfSet_of_atomic}, \lean{HPC.BJ59.IsAtomicCircPAMap.isSnakeHalvesMap}.\\[4pt]
\multicolumn{2}{@{}l}{\emph{Chapter~\ref{ch:chains}}}\\
Definition~\ref{def:bjchain} & \lean{HPC.BJ59.IsBJChain}. These are Fugate's \emph{taut} chains: \lean{HPC.BJ59.isBJChain_iff_isTaut}.\\
Lemma~\ref{lem:adj-closure} & \lean{HPC.BJ59.IsBJChain.adjacent_of_closure}.\\
Lemma~\ref{lem:adjacent-meet} & \lean{HPC.BJ59.adjacent_meet_of_cover}.\\
Theorem~\ref{thm:shrink} & \lean{HPC.BJ59.exists_bjChain_shrink}, from \lean{HPC.BJ59.exists_bjChain_shrink'}. The argument is Barrett's, as used for taut chains in \lean{HPC.Plane.exists_taut_epsChain}.\\
Definition~\ref{def:private} & \lean{HPC.BJ59.HasPrivatePoints}.\\
Proposition~\ref{prop:private-interior} & \lean{HPC.BJ59.exists_privatePoint_interior}.\\
Proposition~\ref{prop:private-shrink} & \lean{HPC.BJ59.exists_privatePoint_shrink}, with \lean{HPC.BJ59.exists_injective_choice}.\\
Lemma~\ref{lem:runs} & \lean{HPC.BJ59.exists_minimal_run}, \lean{HPC.BJ59.exists_run_meeting}, \lean{HPC.BJ59.exists_point_only_of_minimal}.\\
Definition~\ref{def:chainmap} & \lean{HPC.BJ59.IsChainMap}.\\
Lemma~\ref{lem:chainmap-ivt} & \lean{HPC.BJ59.chainMap_ivt}, from \lean{HPC.BJ59.exists_eq_of_steps}.\\
Theorem~\ref{thm:bj6} & \lean{HPC.BJ59.adjacent_of_properCover}; Bing and Jones's form is \lean{HPC.BJ59.bj_theorem6} (with \lean{HPC.BJ59.sign_persist}), which is not used.\\
Text of ``Folding a chain along a walk'' & \lean{HPC.BJ59.foldChain}, \lean{HPC.BJ59.isBJChain_foldChain}.\\
Lemma~\ref{lem:unfold} & \lean{HPC.BJ59.unfoldIdx}, \lean{HPC.BJ59.unfoldFin}, \lean{HPC.BJ59.unfoldMain}, \lean{HPC.BJ59.isPattern_unfold}.\\
Lemma~\ref{lem:fibre-spread} & \lean{HPC.BJ59.exists_fibre_spread}.\\
Definition~\ref{def:fibre-fine} & \lean{HPC.BJ59.IsFibreFine}.\\
Lemma~\ref{lem:fibre-fine} & \lean{HPC.BJ59.isFibreFine_of_isEpsChain}.\\[4pt]
\multicolumn{2}{@{}l}{\emph{Chapter~\ref{ch:arcs}}}\\
Theorem~\ref{thm:bj3} & \lean{HPC.BJ59.theorem3_arc} (with explicit hypotheses), \lean{HPC.BJ59.bj_theorem3}; the chain is \lean{HPC.BJ59.chain_assembly}.\\
Text of ``Theorem 3: chains from end to end'' & \lean{HPC.BJ59.theorem2_arc}, \lean{HPC.BJ59.bj_theorem2}.\\
Theorem~\ref{thm:bj4} & \lean{HPC.BJ59.bj_theorem4}, from \lean{HPC.BJ59.theorem4_of_chain} and \lean{HPC.BJ59.theorem4_explicit}.\\
Corollary~\ref{cor:bj-cut} & \lean{HPC.BJ59.bj_cut}.\\
Lemma~\ref{lem:deep} & \lean{HPC.BJ59.exists_deep_radius}, \lean{HPC.BJ59.exists_deep_sets}.\\
Proposition~\ref{prop:bj5-false} & \lean{HPC.BJ59.bj_theorem5_printed_false_general}; a concrete instance on $\I$ is \lean{HPC.BJ59.bj_theorem5_printed_false}.\\
Theorem~\ref{thm:bj5} & \lean{HPC.BJ59.theorem5_repaired}, \lean{HPC.BJ59.bj_theorem5}; the chained parts are \lean{HPC.BJ59.chainedAt_union_point}. The hypotheses ``$\cl U$ in a link'' and ``$E$ covers $X$'' are not used.\\
Lemma~\ref{lem:window-transfer} & \lean{HPC.BJ59.window_transfer}, with \lean{HPC.BJ59.winScale}, \lean{HPC.BJ59.winInv}.\\
Theorem~\ref{thm:bj5-window} & \lean{HPC.BJ59.bj_theorem5_window}, \lean{HPC.BJ59.bj_theorem5_window_aux}.\\
Lemma~\ref{lem:window-set} & \lean{HPC.BJ59.exists_window_set}.\\
Lemma~\ref{lem:meets-all} & \lean{HPC.BJ59.meets_all_of_meets_ends}.\\
Theorem~\ref{thm:step1b} & \lean{HPC.BJ59.bj_step1b}, with Step~1 as \lean{HPC.BJ59.exists_nhd_step1b}.\\[4pt]
\multicolumn{2}{@{}l}{\emph{Chapter~\ref{ch:thm78}}}\\
Theorem~\ref{thm:bj7} & \lean{HPC.BJ59.bj_theorem7'} and \lean{HPC.BJ59.bj_theorem7}; inside $P$, \lean{HPC.BJ59.bj7_compact}, with \lean{HPC.BJ59.traceChain}, \lean{HPC.BJ59.exists_chain_isFrom}, \lean{HPC.BJ59.prependDesc}, \lean{HPC.BJ59.exists_follow_pattern}; the lift uses \lean{HPC.BJ59.exists_bjChain_shrink'}.\\
Theorem~\ref{thm:bj7-ends} & \lean{HPC.BJ59.bj_theorem7_ends'}; inside $P$, \lean{HPC.BJ59.bj7_ends_compact}, with \lean{HPC.BJ59.collar_data}, \lean{HPC.BJ59.exists_collar_chain}, \lean{HPC.BJ59.bufIdx}, \lean{HPC.BJ59.bufPat}.\\
Remark~\ref{rem:bj7-false} & \lean{HPC.BJ59.bj_theorem7_ends_printed_false}, \lean{HPC.BJ59.bj_theorem7_proof_n3}.\\
Theorem~\ref{thm:bj8} & \lean{HPC.BJ59.bj_theorem8}, from \lean{HPC.BJ59.bj_theorem8_le} and \lean{HPC.BJ59.bj_theorem8_ends}.\\
Lemma~\ref{lem:walk-core} & \lean{HPC.BJ59.walk_core}.\\
Corollary~\ref{cor:bj8} & \lean{HPC.BJ59.bj_theorem8_cor}. The irreducibility of $F$ assumed by Bing and Jones is not needed.\\
Text of ``The corollary to Theorem 8'' & \lean{HPC.BJ59.bj_theorem8_cor_sub}, \lean{HPC.BJ59.bj_theorem8_cor_sub_three}, \lean{HPC.BJ59.bj_theorem8_cor_sub_one}, \lean{HPC.BJ59.zigN}.\\[4pt]
\multicolumn{2}{@{}l}{\emph{Chapter~\ref{ch:zigzag}}}\\
Text of ``The data'' & \lean{HPC.BJ59.starImage}.\\
Lemma~\ref{lem:star-props} & \lean{HPC.BJ59.mem_starImage}, \lean{HPC.BJ59.starImage_succ_subset}, \lean{HPC.BJ59.starImage_dist_le}.\\
Theorem~\ref{thm:zigzag} & \lean{HPC.BJ59.zigzag_homeomorph}. Only the chains $E^i$ need to be BJ-chains.\\
Proposition~\ref{prop:zigzag-fibre} & \lean{HPC.BJ59.zigzag_fibre}.\\[4pt]
\multicolumn{2}{@{}l}{\emph{Chapter~\ref{ch:thm10}}}\\
Theorem~\ref{thm:bj10} & \lean{HPC.BJ59.bjThm10}, a proof of the statement \lean{HPC.BJ59.BJThm10}.\\
Definition~\ref{def:conditions} & \lean{HPC.BJ59.IsFibreChainMap} (Condition A), \lean{HPC.BJ59.EndCompat} (Condition C).\\
Text of ``The conditions carried along'' & \lean{HPC.BJ59.IsPatternOn}.\\
Theorem~\ref{thm:step} & The statement is \lean{HPC.BJ59.BJStep}, proved as \lean{HPC.BJ59.bj_step}.\\
Lemma~\ref{lem:end-piece} & \lean{HPC.BJ59.exists_end_piece}, from \lean{HPC.BJ59.exists_pattern}.\\
Lemma~\ref{lem:subarc-piece} & \lean{HPC.BJ59.exists_subarc_piece}.\\
Lemma~\ref{lem:tau} & \lean{HPC.BJ59.exists_tau_run}, \lean{HPC.BJ59.exists_tau_incl}. The gluing (e)--(g) is \lean{HPC.BJ59.bj_glue}; the pieces are \lean{HPC.BJ59.exists_subarc_piece}, \lean{HPC.BJ59.exists_end_piece}; the windows use \lean{HPC.BJ59.adjacent_of_properCover}, \lean{HPC.BJ59.exists_window_set}, \lean{HPC.BJ59.bj_theorem5_window}.\\
Text of ``The induction'' & \lean{HPC.BJ59.bjThm10_of_step}, with \lean{HPC.BJ59.BJStage}, \lean{HPC.BJ59.BJStage.exists_init}, \lean{HPC.BJ59.BJStage.exists_next}. Bing and Jones's Theorem~9 is nevertheless proved: \lean{HPC.BJ59.bj_theorem9}.\\[4pt]
\multicolumn{2}{@{}l}{\emph{Chapter~\ref{ch:uniqueness}}}\\
Lemma~\ref{lem:one-half} & \lean{HPC.BJ59.exists_half_homeomorph'}, with \lean{HPC.BJ59.halfFibreEquiv}.\\
Theorem~\ref{thm:glue} & \lean{HPC.BJ59.exists_homeomorph_glue'}; Theorem~\ref{thm:bj10} enters as the hypothesis \lean{HPC.BJ59.BJThm10}, discharged by \lean{HPC.BJ59.bjThm10}.\\
Corollary~\ref{cor:snakehalves-homeo} & \lean{HPC.BJ59.homeomorph_of_isSnakeHalvesMap}, with \lean{HPC.Plane.theorem1_HIC'}.\\
Theorem~\ref{thm:unique} & \lean{HPC.BJ59.homeomorph_of_isCircleOfPA'}; with Theorem~10 as a hypothesis, \lean{HPC.BJ59.homeomorph_of_isCircleOfPA} and \lean{HPC.BJ59.homeomorph_of_isCircPAMap}.\\
Theorem~\ref{thm:homog} & \lean{HPC.BJ59.isHomogeneous_of_isCircleOfPA'}; with Theorem~10 as a hypothesis, \lean{HPC.BJ59.isHomogeneous_of_isCircleOfPA}, \lean{HPC.BJ59.isHomogeneous_of_isCircPAMap}; the pseudo-arc facts are \lean{HPC.Plane.theorem1_HIC'} and \lean{HPC.Plane.homogeneous_of_isPseudoArc}.\\
Lemma~\ref{lem:circlelike-homeo} & \lean{HPC.BJ59.isCircleLike_circle}, \lean{HPC.BJ59.isCircleLike_univ_of_homeomorph}.\\
Theorem~\ref{thm:dp61} & \lean{HPC.BJ59.dp22_theorem61'} and \lean{HPC.BJ59.dp22_theorem61_iff'}; with the inputs as hypotheses, \lean{HPC.BJ59.dp22_theorem61}, \lean{HPC.BJ59.isCircleOfPA_of_atomic}, \lean{HPC.BJ59.dp22_theorem61_iff}; the converse uses \lean{HPC.BJ59.exists_atomic_of_plane_isCircleOfPA} and \lean{HPC.BJ59.IsAtomicCircPAMap.comp_homeomorph}.\\
Text of ``Recognition'' & \lean{HPC.BJ59.homeomorph_of_atomic}, \lean{HPC.BJ59.isHomogeneous_of_atomic}.\\
Theorem~\ref{thm:ho-circlelike} & \lean{HPC.BJ59.hoCircleLike'}, a proof of the statement \lean{HPC.BJ59.HOCircleLike}; with Theorem~10 and the plane example as hypotheses, \lean{HPC.BJ59.hoCircleLike}.\\[4pt]
\multicolumn{2}{@{}l}{\emph{Chapter~\ref{ch:circdescr}}}\\
Text of ``The definition'' & \lean{HPC.Fold.PatCrookedOn}.\\
Definition~\ref{def:circdescr} & \lean{HPC.BJ59.CircDescr}, \lean{HPC.BJ59.CircDescr.M}; the centres are \lean{HPC.BJ59.CircDescr.c}.\\
Text of ``The definition'' & \lean{HPC.BJ59.CircDescr.fib}, \lean{HPC.BJ59.CircDescr.run_iff}, \lean{HPC.BJ59.CircDescr.run_len}.\\
Theorem~\ref{thm:circdescr} & \lean{HPC.BJ59.CircDescr.isCircleOfPA}.\\
Lemma~\ref{lem:cd-continuum} & \lean{HPC.BJ59.CircDescr.closure_U_succ}, \lean{HPC.BJ59.CircDescr.M_eq}, \lean{HPC.BJ59.CircDescr.isCompact_M}, \lean{HPC.BJ59.CircDescr.isConnected_M}.\\
Lemma~\ref{lem:cd-meet} & \lean{HPC.BJ59.CircDescr.exists_pat_pair}, \lean{HPC.BJ59.CircDescr.inter_M_nonempty}.\\
Lemma~\ref{lem:cd-fibre} & \lean{HPC.BJ59.CircDescr.down}, \lean{HPC.BJ59.CircDescr.fib_subset_M}, \lean{HPC.BJ59.CircDescr.fib_meets}, \lean{HPC.BJ59.CircDescr.fib_add_int}, \lean{HPC.BJ59.CircDescr.exists_fib}, \lean{HPC.BJ59.CircDescr.eq_add_int_of_mem}.\\
Definition~\ref{def:cd-fmap} & \lean{HPC.BJ59.CircDescr.fmap}, \lean{HPC.BJ59.CircDescr.fmap_eq_iff}.\\
Proposition~\ref{prop:cd-map} & \lean{HPC.BJ59.CircDescr.continuous_fmap}, \lean{HPC.BJ59.CircDescr.surjective_fmap}, \lean{HPC.BJ59.CircDescr.isOpenMap_fmap}.\\
Lemma~\ref{lem:cd-runs} & \lean{HPC.BJ59.CircDescr.runChain}, \lean{HPC.BJ59.CircDescr.runChain_crooked}, with \lean{HPC.BJ59.CircDescr.tight_run}, \lean{HPC.BJ59.CircDescr.pat_mem_run}.\\
Proposition~\ref{prop:cd-fibre-pa} & \lean{HPC.BJ59.CircDescr.isPseudoArc_fib}, from \lean{HPC.BJ59.CircDescr.isContinuum_fib}, \lean{HPC.BJ59.CircDescr.fib_nontrivial}, \lean{HPC.BJ59.CircDescr.isSnakeLike_fib}, \lean{HPC.BJ59.CircDescr.isHereditarilyIndecomposable_fib}.\\
Lemma~\ref{lem:cd-circlelike} & \lean{HPC.BJ59.CircDescr.traceChain}, \lean{HPC.BJ59.CircDescr.isCircleLike_M}.\\[4pt]
\multicolumn{2}{@{}l}{\emph{Chapter~\ref{ch:sweep}}}\\
Text of ``chapter introduction'' & \lean{HPC.BJ59.PatCrookedOnZ}, with the passage to and from sequences indexed by $\N$, \lean{HPC.BJ59.PatCrookedOnZ.toNat}, \lean{HPC.BJ59.PatCrookedOn.toZ}.\\
Theorem~\ref{thm:sweep} & \lean{HPC.BJ59.sweep_lemma}.\\
Text of ``The sweep lemma'' & \lean{HPC.BJ59.sweep_lemma3}, with \lean{HPC.BJ59.sweep_window_case2b}.\\
Definition~\ref{def:sweepX} & \lean{HPC.BJ59.sweepX}, with \lean{HPC.BJ59.sweepL}, \lean{HPC.BJ59.sweepQ}, \lean{HPC.BJ59.swT}, \lean{HPC.BJ59.swB}, \lean{HPC.BJ59.swD}, \lean{HPC.BJ59.swK}, \lean{HPC.BJ59.swVal}.\\
Text of ``The pattern'' & \lean{HPC.BJ59.sw_win_up}, \lean{HPC.BJ59.sw_win_down}, \lean{HPC.BJ59.swVal_junction}, \lean{HPC.BJ59.sweepX_block_crooked}.\\
Theorem~\ref{thm:sweep-pattern} & \lean{HPC.BJ59.exists_sweepPattern}, from \lean{HPC.BJ59.sweepX_period}, \lean{HPC.BJ59.sweepX_step}, \lean{HPC.BJ59.sweepX_slow}, \lean{HPC.BJ59.sweepX_nest}, \lean{HPC.BJ59.sweepX_surv}, \lean{HPC.BJ59.sweepX_crook}.\\[4pt]
\multicolumn{2}{@{}l}{\emph{Chapter~\ref{ch:plane}}}\\
Definition~\ref{def:std-circ} & \lean{HPC.BJ59.stdRectZ}, \lean{HPC.BJ59.zoneZ}, \lean{HPC.BJ59.circExp}, \lean{HPC.BJ59.stdCircLink}, \lean{HPC.BJ59.stdCircZone}.\\
Text of ``Standard circular chains'' & \lean{HPC.BJ59.CircStd.isOpen_stdCircLink}, \lean{HPC.BJ59.CircStd.isPreconnected_stdCircLink}, \lean{HPC.BJ59.CircStd.stdCircLink_periodic}.\\
Lemma~\ref{lem:std-meet} & \lean{HPC.BJ59.stdCircLink_meet_iff}.\\
Theorem~\ref{thm:circ-step} & \lean{HPC.BJ59.circ_step}. The joint positions are \lean{HPC.BJ59.CircJ.jointZ}, with the link lemma \lean{HPC.BJ59.CircJ.link_facts_Z}; the interpolation is \lean{HPC.BJ59.Ramp.phat}; the shear is \lean{HPC.BJ59.cover}; the descent is \lean{HPC.BJ59.Descent.descHomeo}, with \lean{HPC.BJ59.Descent.E}; the estimates are \lean{HPC.BJ59.CircStepAux.small_images} and \lean{HPC.BJ59.CircStepAux.phat_height}.\\
Text of ``The assembly'' & \lean{HPC.BJ59.CircPlane.Lvl}, \lean{HPC.BJ59.CircPlane.L0}, \lean{HPC.BJ59.CircPlane.next}, \lean{HPC.BJ59.CircPlane.lvl}.\\
Theorem~\ref{thm:circdescr-plane} & \lean{HPC.BJ59.planeDescr}, \lean{HPC.BJ59.exists_circDescr_plane}, with \lean{HPC.BJ59.CircPlane.link_diam}, \lean{HPC.BJ59.CircPlane.arcNest}, \lean{HPC.BJ59.CircPlane.surv}, \lean{HPC.BJ59.CircPlane.crook}.\\
Theorem~\ref{thm:plane-example} & \lean{HPC.BJ59.bjPlaneExample}, a proof of the statement \lean{HPC.BJ59.BJPlaneExample}.\\[4pt]
\multicolumn{2}{@{}l}{\emph{Chapter~\ref{ch:classification}}}\\
Theorem~\ref{thm:decomp-case} & \lean{HPC.BJ59.circle_or_circleOfPA_of_decomposable} (with circle-likeness as the hypothesis \lean{HPC.BJ59.HOCircleLike}, discharged by \lean{HPC.BJ59.hoCircleLike'}), from \lean{HPC.BJ59.circle_or_circPAMap_of_decomposable} and \lean{HPC.BJ59.circle_or_atomic_of_decomposable}.\\
Theorem~\ref{thm:classification} & \leansub{HPC.BJ59.ho16_classification}{4}. The indecomposable case in both directions is \lean{HPC.HO16.ho16_indec_homogeneous_iff}, and the converse for pseudo-arcs \lean{HPC.HO16.ho16_pseudoArc_indec_homogeneous}.\\
Text of ``The classification'' & \lean{HPC.BJ59.ho16_classification_atomic'}.\\
Theorem~\ref{thm:classification-plane} & \leansub{HPC.BJ59.ho16_classification_plane}{4}, with \lean{HPC.BJ59.exists_plane_circle}, \lean{HPC.Plane.isPseudoArc_plane}, \lean{HPC.BJ59.bjPlaneExample}, \lean{HPC.Plane.theorem1_HIC'}.\\
Text of ``How the inputs were discharged'' & \leansub{HPC.BJ59.ho16_classification}{4}, \leansub{HPC.BJ59.ho16_classification_plane}{4}.\\[4pt]
\end{longtable}


\end{document}
